\nonstopmode{}
\documentclass[a4paper]{book}
\usepackage[times,inconsolata,hyper]{Rd}
\usepackage{makeidx}
\usepackage[utf8]{inputenc} % @SET ENCODING@
% \usepackage{graphicx} % @USE GRAPHICX@
\makeindex{}
\begin{document}
\chapter*{}
\begin{center}
{\textbf{\huge Package `sommer'}}
\par\bigskip{\large \today}
\end{center}
\ifthenelse{\boolean{Rd@use@hyper}}{\hypersetup{pdftitle = {sommer: Solving Mixed Model Equations in R}}}{}
\ifthenelse{\boolean{Rd@use@hyper}}{\hypersetup{pdfauthor = {Giovanny Covarrubias-Pazaran}}}{}
\begin{description}
\raggedright{}
\item[Type]\AsIs{Package}
\item[Title]\AsIs{Solving Mixed Model Equations in R}
\item[Version]\AsIs{4.5.00}
\item[Date]\AsIs{2026-09-30}
\item[Maintainer]\AsIs{Giovanny Covarrubias-Pazaran }\email{cova\_ruber@live.com.mx}\AsIs{}
\item[Description]\AsIs{Structural multivariate-univariate linear mixed model solver for estimation of multiple random effects with unknown variance-covariance structures (e.g., unstructured, diagonal, autoregressive, factor analytic, reduced rank, among others) and known covariance among levels of random effects (e.g., pedigree and genomic relationship matrices) (Covarrubias-Pazaran, 2016 <}\Rhref{https://doi.org/10.1371/journal.pone.0156744}{doi:10.1371/journal.pone.0156744}\AsIs{>; Maier et al., 2015 <}\Rhref{https://doi.org/10.1016/j.ajhg.2014.12.006}{doi:10.1016/j.ajhg.2014.12.006}\AsIs{>; Jensen et al., 1997). REML estimates can be obtained using the Direct-Inversion Newton-Raphson and Direct-Inversion Average Information algorithms for the problems r x r (r being the number of records) or using the Henderson-based average information algorithm for the problem c x c (c being the number of coefficients to estimate).}
\item[Depends]\AsIs{R (>= 3.5.0), Matrix (>= 1.6-2), methods, stats, MASS, crayon,
enhancer}
\item[LazyLoad]\AsIs{yes}
\item[License]\AsIs{GPL-3}
\item[Imports]\AsIs{Rcpp (>= 0.12.19)}
\item[BugReports]\AsIs{}\url{https://github.com/covaruber/sommer/issues}\AsIs{}
\item[URL]\AsIs{}\url{https://github.com/covaruber/sommer}\AsIs{}
\item[LinkingTo]\AsIs{Rcpp, RcppArmadillo, RcppProgress, RcppEigen, Matrix}
\item[SystemRequirements]\AsIs{METIS graph partitioning library (optional; enables
nested-dissection sparse ordering for large relationship
matrices, auto-detected at configure time, falls back to AMD
ordering when absent), NVIDIA CUDA toolkit with cuBLAS and
cuSOLVER (optional; enables GPU kernels on Linux through
SOMMER_CUDA_HOME; CUDA libraries are loaded dynamically)}
\item[Suggests]\AsIs{flexiblas, RhpcBLASctl, parallelly, rmarkdown, knitr, plyr,
parallel, orthopolynom, RSpectra, lattice, emmeans (>= 1.4),
estimability, testthat (>= 3.0.0)}
\item[VignetteBuilder]\AsIs{knitr}
\item[Config/testthat/edition]\AsIs{3}
\item[RoxygenNote]\AsIs{7.3.2}
\item[NeedsCompilation]\AsIs{yes}
\item[Author]\AsIs{Giovanny Covarrubias-Pazaran [aut, cre]
(<}\url{https://orcid.org/0000-0002-7194-3837}\AsIs{>)}
\end{description}
\Rdcontents{\R{} topics documented:}
\inputencoding{utf8}
\HeaderA{sommer-package}{\strong{So}lving \strong{M}ixed \strong{M}odel \strong{E}quations in \strong{R} \\{}}{sommer.Rdash.package}
\aliasA{sommer}{sommer-package}{sommer}
\keyword{R package}{sommer-package}
%
\begin{Description}\relax

Sommer is a structural univariate and multivariate linear mixed-model package
for fitting models with multiple random effects and flexible covariance
structures. Variance parameters are estimated by restricted maximum likelihood
(REML). The package provides two complementary numerical formulations.

\code{\LinkA{mmes}{mmes}} uses Henderson's mixed model equations with an Average
Information REML algorithm. The mixed-model coefficient matrix is represented
as a sparse system and factorized with Eigen/Armadillo routines. Sparse LDLT
factorizations are reused throughout the REML calculations. Selected inverse
elements required for likelihood derivatives and prediction error variances can
be obtained with Takahashi sparse-inverse recursions, so a complete dense
coefficient-matrix inverse is not routinely formed.

\code{\LinkA{mmer}{mmer}} provides the marginal-covariance MNR formulation. Its REML
calculations are organized around the covariance matrix of the observations and
the corresponding REML projection matrix, with Newton-Raphson/Fisher-scoring
and Average Information updates for variance parameters.

Both formulations support multiple random effects and structured covariance
models. Through \code{\LinkA{vsm}{vsm}}, covariance structures can be composed from
an arbitrary number of covariance factors, including identity, diagonal,
unstructured, autoregressive, moving-average, Toeplitz, factor-analytic,
reduced-rank, antedependence, correlation and spatial structures, as well as
user-defined covariance functions.

The marginal formulation in \code{mmer} performs its principal REML
calculations in observation space, whereas \code{mmes} performs them through
the sparse mixed-model coefficient system. Their relative computational
performance therefore depends on the number of observations, the number of
mixed-model coefficients, sparsity, and the covariance structures in the
fitted model.

The numerical algorithms are coded primarily in C++ using Armadillo and Eigen.
Sommer returns REML variance-covariance estimates, BLUEs, BLUPs, residuals,
fitted values, information matrices and, when requested, prediction error
variance information.

\end{Description}
%
\begin{Section}{Functions for genetic analysis}

The package provides kernels to estimate additive (\code{\LinkA{A.mat}{A.mat}}), dominance
(\code{\LinkA{D.mat}{D.mat}}), epistatic (\code{\LinkA{E.mat}{E.mat}}), single-step (\code{\LinkA{H.mat}{H.mat}}) 
relationship matrices for diploid and polyploid organisms. It also provides flexibility 
to fit other genetic models such as full and half diallel models and random regression models.

A good converter from letter code to numeric format is implemented in the function 
\code{\LinkA{atcg1234}{atcg1234}}, which supports higher ploidy levels than diploid. 
Additional functions for genetic analysis have been included such as build a 
genotypic hybrid marker matrix (\code{\LinkA{build.HMM}{build.HMM}}), plot of genetic maps
(\code{\LinkA{map.plot}{map.plot}}), creation of manhattan plots (\code{\LinkA{manhattan}{manhattan}}).
If you need to use pedigree you need to convert your pedigree into a relationship 
matrix (use the `getA` function from the pedigreemm package).

\end{Section}
%
\begin{Section}{Functions for statistical analysis and S3 methods}

The \code{\LinkA{vpredict}{vpredict}} function can be used to estimate standard errors for linear combinations of variance components (e.g. ratios like h2). The \code{\LinkA{r2}{r2}} function calculates reliability. S3 methods are available for some parameter extraction such as:

+ \code{\LinkA{predict.mmes}{predict.mmes}}

+ \code{\LinkA{fitted.mmes}{fitted.mmes}}

+ \code{\LinkA{residuals.mmes}{residuals.mmes}}

+ \code{\LinkA{summary.mmes}{summary.mmes}}

+ \code{\LinkA{coef.mmes}{coef.mmes}}

+ \code{\LinkA{anova.mmes}{anova.mmes}}

+ \code{\LinkA{plot.mmes}{plot.mmes}}

\end{Section}
%
\begin{Section}{Functions for trial analysis}

Recently, spatial modeling has been added added to sommer using the two-dimensional 
spline (\code{\LinkA{spl2Dc}{spl2Dc}}).

\end{Section}
%
\begin{Section}{Keeping sommer updated}
The sommer package is updated on CRAN every 4-months due to CRAN policies but you can find the latest source at https://github.com/covaruber/sommer. This can be easily installed typing the following in the R console:

\code{library(devtools)}

\code{install\_github("covaruber/sommer")}

This is recommended if you reported a bug, was fixed and was immediately pushed to GitHub but not in CRAN until the next update.
\end{Section}
%
\begin{Section}{Tutorials}

\strong{For tutorials} on how to perform different analysis with sommer please look at the vignettes by typing in the terminal:

\code{vignette("sommer.qg")}

\code{vignette("sommer.gxe")}

\code{vignette("sommer.vs.lme4")} 

\code{vignette("sommer.spatial")}

\end{Section}
%
\begin{Section}{Getting started}
The package has been equiped with several datasets to learn how to use the sommer package (and almost to learn all sort of quantitative genetic analysis): 

* \code{\LinkA{DT\_halfdiallel}{DT.Rul.halfdiallel}}, \code{\LinkA{DT\_fulldiallel}{DT.Rul.fulldiallel}} and \code{\LinkA{DT\_mohring}{DT.Rul.mohring}} datasets have examples to fit half and full diallel designs. 

* \code{\LinkA{DT\_h2}{DT.Rul.h2}} to calculate heritability

* \code{\LinkA{DT\_cornhybrids}{DT.Rul.cornhybrids}} and \code{\LinkA{DT\_technow}{DT.Rul.technow}} datasets to perform genomic prediction in hybrid single crosses 

* \code{\LinkA{DT\_wheat}{DT.Rul.wheat}} dataset to do genomic prediction in single crosses in species displaying only additive effects.

* \code{\LinkA{DT\_cpdata}{DT.Rul.cpdata}} dataset to fit genomic prediction models within a biparental population coming from 2 highly heterozygous parents including additive, dominance and epistatic effects. 

* \code{\LinkA{DT\_polyploid}{DT.Rul.polyploid}} to fit genomic prediction and GWAS analysis in polyploids. 

* \code{\LinkA{DT\_gryphon}{DT.Rul.gryphon}} data contains an example of an animal model including pedigree information.

* \code{\LinkA{DT\_btdata}{DT.Rul.btdata}} dataset contains an animal (birds) model.

* \code{\LinkA{DT\_legendre}{DT.Rul.legendre}} simulated dataset for random regression model.

* \code{\LinkA{DT\_sleepstudy}{DT.Rul.sleepstudy}} dataset to know how to translate lme4 models to sommer models.

\end{Section}
%
\begin{Section}{Differences of sommer >= 4.4.1 with previous versions}

Since version 4.4.1, I have unified the use of the two different solving algorithms
into the mmes function by just using the new argument \code{henderson} which by default 
is set to FALSE. Other than that the rest is the same with the addition that now
the identity terms needs to be encapsulated in the \code{\LinkA{ism}{ism}} function. In 
addition, now the multi-trait models need to be fitted in the long format. This are 
few but major changes to the way sommer models are fitted.

\end{Section}
%
\begin{Section}{Differences of sommer >= 4.1.7 with previous versions}

Since version 4.1.7 I have introduced the mmes-based average information function `mmec` which is much faster when dealing with the r > c problem (more records than coefficients to estimate). This introduces its own covariance structure functons such as vsc(), usc(), dsc(), atc(), csc(). Please give it a try, although is in early phase of development.

\end{Section}
%
\begin{Section}{Differences of sommer >= 3.7.0 with previous versions}

Since version 3.7 I have completly redefined the specification of the variance-covariance 
structures to provide more flexibility to the user. This has particularly helped the 
residual covariance structures and the easier combination of custom random effects and 
overlay models. I think that although this will bring some uncomfortable situations at 
the beggining, in the long term this will help users to fit better models. In esence, 
I have abandoned the asreml formulation (not the structures available) given it's 
limitations to combine some of the sommer structures but all covariance structures 
can now be fitted using the `vsm` functions. 

\end{Section}
%
\begin{Section}{Differences of sommer >= 3.0.0 with previous versions}

Since version 3.0 I have decided to focus in developing the multivariate solver and for doing this I have decided to remove the M argument (for GWAS analysis) from the mmes function and move it to it's own function GWAS.

Before the mmes solver had implemented the usm(trait), diag(trait), at(trait) asreml formulation for multivariate models that allow to specify the structure of the trait in multivariate models. Therefore the MVM argument was no longer needed. After version 3.7 now the multi-trait structures can be specified in the \code{Gt} and \code{Gtc} arguments of the \code{\LinkA{vsm}{vsm}} function.

The Average Information algorithm had been removed in the past from the package because of its instability to deal with very complex models without good initial values. Now after 3.7 I have brought it back after I noticed that starting with NR the first three iterations gives enough flexibility to the AI algorithm.

Keep in mind that sommer uses direct inversion (DI) algorithm which can be very slow for datasets with many observations (big 'n'). The package is focused in problems of the type p > n (more random effect(s) levels than observations) and models with dense covariance structures. For example, for experiment with dense covariance structures with low-replication (i.e. 2000 records from 1000 individuals replicated  twice with a covariance structure of 1000x1000) sommer will be faster than MME-based software. Also for genomic problems with large number of random effect levels, i.e. 300 individuals (n) with 100,000 genetic markers (p). On the other hand, for highly replicated trials with small covariance structures or n > p (i.e. 2000 records from 200 individuals replicated 10 times with covariance structure of 200x200) asreml or other MME-based algorithms will be much faster and I recommend you to use that software.

\end{Section}
%
\begin{Section}{Models Enabled}

\strong{General linear mixed model}

Both numerical engines fit the Gaussian linear mixed-model family

\deqn{y = X\beta + Zu + e,}{}

with

\deqn{u \sim N(0,G), \qquad e \sim N(0,R),}{}

and therefore

\deqn{V = \mathrm{Var}(y) = ZGZ^\prime + R.}{}

For several independent random terms this becomes

\deqn{V = \sum_{k=1}^{q} Z_k G_k Z_k^\prime + R.}{}

Here \eqn{y}{} is the response vector, \eqn{X}{} and \eqn{Z}{} are design matrices,
\eqn{\beta}{} contains fixed effects, \eqn{u}{} contains random effects, and
\eqn{e}{} contains residual effects. The same representation covers univariate
and multivariate models after the corresponding responses, design matrices and
covariance structures are assembled.

Covariance models specified through \code{\LinkA{vsm}{vsm}} use a product-level
variance scale and an arbitrary number of covariance-shaping factors:

\deqn{\Sigma = \sigma^2(K_1 \otimes K_2 \otimes \cdots \otimes K_m).}{}

For a random effect with known relationship covariance \eqn{A}{},

\deqn{G = \sigma^2(K_1 \otimes K_2 \otimes \cdots \otimes K_m) \otimes A.}{}

The CovarianceFactor interface separates evaluation of each \eqn{K_j}{}, its
derivatives, parameter reporting and trust limits from the REML solver. New
covariance structures can therefore be added without changing the core
\code{mmes} optimization algorithm.

\strong{Marginal MNR formulation in mmer}

The MNR engine used by \code{\LinkA{mmer}{mmer}} organizes the REML calculations
around the marginal covariance \eqn{V}{}. Define the REML projection matrix

\deqn{P = V^{-1} -
V^{-1}X(X^\prime V^{-1}X)^{-1}X^\prime V^{-1}.}{}

Apart from constants independent of the variance parameters, the Gaussian REML
log-likelihood is

\deqn{\ell_R(\theta) =
-\frac{1}{2}\{\log|V|+\log|X^\prime V^{-1}X|+y^\prime P y\}.}{}

For variance parameter \eqn{\theta_i}{}, let

\deqn{V_i = \frac{\partial V}{\partial\theta_i}.}{}

For covariance models linear in the fitted variance parameters, the REML score
has the standard form

\deqn{s_i =
\frac{1}{2}\{y^\prime P V_i P y-\mathrm{tr}(P V_i)\}.}{}

The MNR C++ implementation forms the projected covariance derivatives
\eqn{PV_i}{}. Its Average Information matrix is

\deqn{AI_{ij} =
\frac{1}{2}y^\prime P V_i P V_j P y,}{}

while its Newton-Raphson/Fisher-scoring path uses trace products of the
projected covariance derivatives. The resulting information system determines
the variance-parameter search direction, and the implementation can blend the
information calculation with an EM-type diagonal stabilization during early
iterations.

This formulation is referred to as \emph{marginal} because its principal
covariance calculations are carried out in observation space through \eqn{V}{}
and \eqn{P}{}. This describes the numerical formulation more accurately than
naming it after a particular matrix operation.

\strong{Sparse Henderson Average Information formulation in mmes}

The \code{\LinkA{mmes}{mmes}} engine uses Henderson's mixed model equations. With
\eqn{W=[X\;Z]}{}, the coefficient system is

\deqn{C =
\left[
\begin{array}{cc}
X^\prime R^{-1}X & X^\prime R^{-1}Z \\
Z^\prime R^{-1}X & Z^\prime R^{-1}Z+G^{-1}
\end{array}
\right],}{}

and the right-hand side is

\deqn{r =
\left[
\begin{array}{c}
X^\prime R^{-1}y \\
Z^\prime R^{-1}y
\end{array}
\right].}{}

The BLUE and BLUP solutions satisfy

\deqn{Cb=r, \qquad
b=(\widehat{\beta}^\prime,\widehat{u}^\prime)^\prime.}{}

The current C++ implementation stores \eqn{C}{} sparsely and uses an Eigen
\code{SimplicialLDLT} factorization. Symbolic factorization information is
reused when the sparsity pattern remains unchanged. Random-effect precision
contributions are inserted blockwise. Structured residual covariance models
use optimized diagonal or repeated-block/Kronecker paths when applicable and a
generic sparse residual-factorization path otherwise.

The REML criterion is evaluated from determinant contributions for the random
covariance structures, residual covariance and mixed-model coefficient matrix,
together with the quadratic term \eqn{y^\prime P y}{}. This is a
mixed-model-equation representation of the same REML objective optimized by
the marginal formulation.

For covariance parameter \eqn{\eta_i}{}, the generic covariance engine supplies

\deqn{\Sigma_i =
\frac{\partial\Sigma}{\partial\eta_i}.}{}

For

\deqn{\Sigma=\sigma^2(K_1\otimes\cdots\otimes K_m),}{}

a parameter belonging to factor \eqn{K_j}{} has derivative

\deqn{\frac{\partial\Sigma}{\partial\eta_{jk}} =
\sigma^2 K_1\otimes\cdots\otimes
\frac{\partial K_j}{\partial\eta_{jk}}
\otimes\cdots\otimes K_m.}{}

For the working scale \eqn{\tau=\log(\sigma^2)}{},

\deqn{\frac{\partial\Sigma}{\partial\tau}=\Sigma.}{}

The score trace terms use selected elements associated with the sparse
coefficient system. These are obtained from the LDLT factors with Takahashi
sparse-inverse recursions; a complete \eqn{C^{-1}}{} is not required during REML
iterations.

Average Information is calculated by analytically differentiating the
factorized mixed-model equations. Differentiating

\deqn{Cb=r}{}

with respect to \eqn{\eta_i}{} gives

\deqn{C b_i = r_i-C_i b,}{}

where

\deqn{C_i=\frac{\partial C}{\partial\eta_i}, \qquad
r_i=\frac{\partial r}{\partial\eta_i}.}{}

The sensitivity systems therefore reuse the same sparse factorization of
\eqn{C}{}. For covariance \eqn{\Sigma}{} with precision
\eqn{\Lambda=\Sigma^{-1}}{} and derivative
\eqn{B_i=\partial\Sigma/\partial\eta_i}{},

\deqn{\Lambda_i=-\Lambda B_i\Lambda.}{}

This factor-differentiation formulation avoids constructing an
observation-space working variate for every random covariance parameter.
Residual/residual and random/residual Average Information terms use the
corresponding mixed-model/Schur-complement identities.

Covariance constructors may provide analytic derivatives in C++ or R. If a
constructor uses numerical differentiation, only its small covariance factor
\eqn{K_j(\eta)}{} is differentiated by central finite differences; the REML
likelihood and downstream score and Average Information calculations are not
numerically differentiated.

\strong{Sparse inverse subsets and prediction error variances}

During optimization, \code{mmes} does not routinely materialize the complete
\eqn{C^{-1}}{}. Takahashi recursions provide the selected inverse subset required
for REML trace calculations. After convergence, \code{control\$output\$inverse=1} obtains the
diagonal information needed for prediction error variances through the sparse
inverse subset, whereas \code{control\$output\$inverse=2} requests the complete
coefficient-matrix inverse.

\strong{Parameter updates and numerical safeguards}

Covariance parameters are optimized in working coordinates chosen to respect
their parameter spaces, such as logarithms for positive scales and
hyperbolic-tangent or bounded-logit mappings for correlations. Each
CovarianceFactor descriptor supplies its natural-scale reporting
transformation and parameter-specific trust limit.

The information system determines a proposed REML update. In the current
\code{mmes} implementation, the proposal is limited in working-coordinate
space and then checked against the exact REML likelihood. If the proposed
global step decreases the accepted likelihood beyond numerical tolerance, the
step is geometrically reduced and reevaluated. A likelihood decrease is
therefore not interpreted as convergence. Boundary handling and robust
information solves are retained for difficult or weakly identified variance
components.

\strong{Choosing between the two formulations}

Both engines estimate the same class of linear mixed models but organize the
numerical work differently. \code{mmer} uses the marginal covariance
representation and is attractive when observation-space calculations are
moderate in size. \code{mmes} uses the sparse Henderson coefficient system and
is particularly attractive when the mixed-model equations are sparse and
their factorization can be reused efficiently. The best choice depends on the
dimensions, sparsity and covariance structures of the fitted model rather than
on a universal rule based only on the number of records or coefficients.

See \code{\LinkA{mmes}{mmes}}, \code{\LinkA{mmer}{mmer}}, and \code{\LinkA{vsm}{vsm}} for the
interfaces and available covariance structures.

\end{Section}
%
\begin{Section}{Bug report and contact}

If you have any questions or suggestions please post it in https://stackoverflow.com or https://stats.stackexchange.com

I'll be glad to help or answer any question. I have spent a valuable amount of time developing this package. Please cite this package in your publication. Type 'citation("sommer")' to know how to cite it.


\end{Section}
%
\begin{Author}\relax
Giovanny Covarrubias-Pazaran
\end{Author}
%
\begin{References}\relax

Covarrubias-Pazaran G. 2016. Genome assisted prediction of quantitative traits using the R package sommer. PLoS ONE 11(6): doi:10.1371/journal.pone.0156744 

Covarrubias-Pazaran G. 2018. Software update: Moving the R package sommer to multivariate mixed models for genome-assisted prediction. doi: https://doi.org/10.1101/354639

Sanderson, C., \& Curtin, R. (2025). Armadillo: An Efficient Framework for Numerical Linear Algebra. arXiv preprint arXiv:2502.03000.

Bernardo Rex. 2010. Breeding for quantitative traits in plants. Second edition. Stemma Press. 390 pp.

Gilmour et al. 1995. Average Information REML: An efficient algorithm for variance parameter estimation in linear mixed models. Biometrics 51(4):1440-1450.

Henderson C.R. 1975. Best Linear Unbiased Estimation and Prediction under a Selection Model. Biometrics vol. 31(2):423-447.

Kang et al. 2008. Efficient control of population structure in model organism association mapping. Genetics 178:1709-1723.

Lee et al. 2015. MTG2: An efficient algorithm for multivariate linear mixed model analysis based on genomic information. Cold Spring Harbor. doi: http://dx.doi.org/10.1101/027201.

Maier et al. 2015. Joint analysis of psychiatric disorders increases accuracy of risk prediction for schizophrenia, bipolar disorder, and major depressive disorder. Am J Hum Genet; 96(2):283-294.

Searle. 1993. Applying the EM algorithm to calculating ML and REML estimates of variance components. Paper invited for the 1993 American Statistical Association Meeting, San Francisco.

Yu et al. 2006. A unified mixed-model method for association mapping that accounts for multiple levels of relatedness. Genetics 38:203-208.

Tunnicliffe W. 1989. On the use of marginal likelihood in time series model estimation. JRSS 51(1):15-27.

\end{References}
%
\begin{Examples}
\begin{ExampleCode}

####=========================================####
#### For CRAN time limitations most lines in the 
#### examples are silenced with one '#' mark, 
#### remove them and run the examples
####=========================================####

####=========================================####
#### EXAMPLES
#### Different models with sommer
####=========================================####

data(DT_example, package="enhancer")





DT <- DT_example
DT=DT[with(DT, order(Env)), ]
head(DT)

####=========================================####
#### Univariate homogeneous variance models  ####
####=========================================####

## Compound simmetry (CS) model
ans1 <- mmes(Yield~Env,
             random= ~ Name + Env:Name,
             rcov= ~ units,
             data=DT)
summary(ans1)

####===========================================####
#### Univariate heterogeneous variance models  ####
####===========================================####
## Compound simmetry (CS) + Diagonal (DIAG) model
ans3 <- mmes(Yield~Env,
             random= ~Name + vsm(dsm(Env),ism(Name)),
             rcov= ~ vsm(dsm(Env),ism(units)),
             data=DT)
summary(ans3)





\end{ExampleCode}
\end{Examples}
\inputencoding{utf8}
\HeaderA{A.mat}{Additive relationship matrix}{A.mat}
%
\begin{Description}\relax
Calculates the realized additive relationship matrix. Currently is the C++ implementation of van Raden (2008). 
\end{Description}
%
\begin{Usage}
\begin{verbatim}
A.mat(X,min.MAF=0,return.imputed=FALSE)
\end{verbatim}
\end{Usage}
%
\begin{Arguments}
\begin{ldescription}
\item[\code{X}] 
Matrix (\eqn{n \times m}{}) of unphased genotypes for \eqn{n}{} lines and \eqn{m}{} biallelic markers, 
coded as \{-1,0,1\}. Fractional (imputed) and missing values (NA) are allowed.

\item[\code{min.MAF}] 
Minimum minor allele frequency. The A matrix is not sensitive to rare alleles, so by default only monomorphic markers are removed.

\item[\code{return.imputed}] 
When TRUE, the imputed marker matrix is returned.  

\end{ldescription}
\end{Arguments}
%
\begin{Details}\relax

For vanraden method: the marker matrix is centered by subtracting column means \eqn{M= X - ms}{} where ms is the coumn means. Then \eqn{A=M M'/c}{}, where \eqn{c = \sum_k{d_k}/k}{}, the mean value of the diagonal values of the \eqn{M M'}{} portion.

\end{Details}
%
\begin{Value}
If return.imputed = FALSE, the \eqn{n \times n}{} additive relationship matrix is returned.

If return.imputed = TRUE, the function returns a list containing
\begin{description}

\item[\$A] the A matrix
\item[\$X] the imputed marker matrix

\end{description}

\end{Value}
%
\begin{References}\relax

Endelman, J.B., and J.-L. Jannink. 2012. Shrinkage estimation of the realized relationship matrix. G3:Genes, Genomes, Genetics. 2:1405-1413. doi: 10.1534/g3.112.004259

Covarrubias-Pazaran G (2016) Genome assisted prediction of quantitative traits using the R package sommer. PLoS ONE 11(6): doi:10.1371/journal.pone.0156744 

\end{References}
%
\begin{Examples}
\begin{ExampleCode}
####=========================================####
#### random population of 200 lines with 1000 markers
####=========================================####
X <- matrix(rep(0,200*1000),200,1000)
for (i in 1:200) {
  X[i,] <- ifelse(runif(1000)<0.5,-1,1)
}

A <- A.mat(X)

####=========================================####
#### take a look at the Genomic relationship matrix 
#### (just a small part)
####=========================================####
# colfunc <- colorRampPalette(c("steelblue4","springgreen","yellow"))
# hv <- heatmap(A[1:15,1:15], col = colfunc(100),Colv = "Rowv")
# str(hv)
\end{ExampleCode}
\end{Examples}
\inputencoding{utf8}
\HeaderA{analyze\_mme\_crossovers}{Analyze Crossover Frontiers from MME Benchmarks}{analyze.Rul.mme.Rul.crossovers}
%
\begin{Description}\relax
Identifies problem configurations (factors, levels, dimensions) where experimental solver
paths provide speedups over baseline methods, fitting empirical crossover conditions.
\end{Description}
%
\begin{Usage}
\begin{verbatim}
analyze_mme_crossovers(
  benchmark,
  baseline_engine = NULL,
  candidate_engine = NULL
)
\end{verbatim}
\end{Usage}
%
\begin{Arguments}
\begin{ldescription}
\item[\code{benchmark}] An object of class \code{sommer\_mme\_benchmark}.

\item[\code{baseline\_engine}] Name of the reference baseline engine. Default is the first engine tested.

\item[\code{candidate\_engine}] Name of the experimental engine to evaluate.
\end{ldescription}
\end{Arguments}
%
\begin{Value}
An object of class \code{"sommer\_mme\_crossover"}.
\end{Value}
\inputencoding{utf8}
\HeaderA{anova.mmes}{anova form a GLMM fitted with mmes}{anova.mmes}
\aliasA{anova.mmer}{anova.mmes}{anova.mmer}
\keyword{models}{anova.mmes}
%
\begin{Description}\relax
\code{anova} method for class \code{"mmes"}.
\end{Description}
%
\begin{Usage}
\begin{verbatim}
## S3 method for class 'mmes'
anova(object, object2=NULL, ...)
\end{verbatim}
\end{Usage}
%
\begin{Arguments}
\begin{ldescription}
\item[\code{object}] an object of class \code{"mmes"}
\item[\code{object2}] an object of class \code{"mmes"}. If supplied, a likelihood ratio test between the two models is returned (not available for PQL fits). If \code{NULL}, Wald tests for the fixed effects of \code{object} are returned by \code{\LinkA{wald\_mmes}{wald.Rul.mmes}}.
\item[\code{...}] Further arguments passed to \code{\LinkA{wald\_mmes}{wald.Rul.mmes}} (e.g. \code{ssType}, \code{denDF}) when \code{object2} is \code{NULL}.
\end{ldescription}
\end{Arguments}
%
\begin{Value}
vector of anova

\end{Value}
%
\begin{Author}\relax
Giovanny Covarrubias
\end{Author}
%
\begin{SeeAlso}\relax
\code{\LinkA{anova}{anova}}, \code{\LinkA{mmes}{mmes}}, \code{\LinkA{wald\_mmes}{wald.Rul.mmes}}
\end{SeeAlso}
\inputencoding{utf8}
\HeaderA{antem}{Antedependence covariance structure}{antem}
%
\begin{Description}\relax
Antedependence covariance structure. The returned dimensionless covariance shape is intended for use inside \code{vsm}, which supplies the single overall variance scale.
\end{Description}
%
\begin{Usage}
\begin{verbatim}
antem(x, order = 1L, beta = NULL, innovations = NULL, fixed = NULL)
\end{verbatim}
\end{Usage}
%
\begin{Arguments}
\begin{ldescription}
\item[\code{x}] Ordered variable defining q covariance levels.
\item[\code{order}] Antedependence order, 1 <= order < q.
\item[\code{beta}] Optional starting regression coefficients for allowed subdiagonals.
\item[\code{innovations}] Optional q positive innovation variances.
\item[\code{fixed}] Logical vector controlling all beta coefficients and q-1 innovation-variance ratios.
\end{ldescription}
\end{Arguments}
%
\begin{Details}\relax
The modified-Cholesky representation is \deqn{Ty=e,\qquad \mathrm{Cov}(e)=D,\qquad K=T^{-1}DT^{-\mathsf T}.}{} The matrix \eqn{T}{} is unit lower triangular, with nonzero regression coefficients only within the requested number of preceding levels.  The first innovation variance is the scale reference; the remaining innovation variances are positive ratios.
\end{Details}
%
\begin{Value}
A list containing the incidence/design matrix in \code{Z} and a compiled CovarianceFactor v2 descriptor in \code{covFactor}, for use inside \code{vsm}.
\end{Value}
%
\begin{SeeAlso}\relax
\code{\LinkA{ar1m}{ar1m}}, \code{\LinkA{toeplitzm}{toeplitzm}}, \code{\LinkA{vsm}{vsm}}.
\end{SeeAlso}
%
\begin{Examples}
\begin{ExampleCode}
## Not run: 
vsm(antem(time, order=2), ism(id))

## End(Not run)
\end{ExampleCode}
\end{Examples}
\inputencoding{utf8}
\HeaderA{APY}{Construct an APY Relationship Precision Matrix}{APY}
%
\begin{Description}\relax
Constructs the Algorithm for Proven and Young (APY) approximation for a
positive-definite relationship matrix and a user-specified core set. In the
genomic evaluation literature, the input is the genomic relationship matrix
\eqn{G}{} for genotyped animals, usually calculated from SNP marker data. The
core and non-core partition is within that genotyped set; it is not a
pedigree-versus-genomic partition. The non-core conditional covariance is
approximated by its diagonal. The returned matrix is a sparse precision matrix
suitable for the \code{Gu} argument of \code{mmes()} in a genomic BLUP model.
\end{Description}
%
\begin{Usage}
\begin{verbatim}
APY(G, core, tol = 1e-10, return.details = FALSE)
\end{verbatim}
\end{Usage}
%
\begin{Arguments}
\begin{ldescription}
\item[\code{G}] A finite, symmetric, positive-definite numeric relationship matrix.
For the standard APY genetic evaluation use, this is the genomic relationship
matrix among genotyped animals, commonly obtained from marker genotypes with
\code{A.mat()}. Row and column names, when present, must agree exactly. This
implementation receives \code{G} as a dense matrix.
\item[\code{core}] Unique one-based row indices or row names selecting the core
individuals.
\item[\code{tol}] Positive tolerance below one, used for symmetry validation and
the minimum acceptable conditional variance relative to the relationship
scale. Small or non-positive conditional variances cause an error; they are
not silently regularized.
\item[\code{return.details}] If \code{TRUE}, return a list containing the sparse
precision matrix as \code{Gu}, core and non-core indices, regression matrix
\code{B}, and non-core conditional variances.
\end{ldescription}
\end{Arguments}
%
\begin{Details}\relax
Partition \eqn{G}{} into core and non-core blocks. With
\eqn{B = G_{nc}G_{cc}^{-1}}{} and
\eqn{D = \operatorname{diag}(G_{nn} - BG_{cn})}{}, APY uses the precision
\eqn{\left[\begin{smallmatrix}
G_{cc}^{-1} + B' D^{-1} B & -B' D^{-1} \\
-D^{-1} B & D^{-1}
\end{smallmatrix}\right]}{}. It preserves the core block and core-to-non-core
relationships while replacing non-core conditional correlations with zero.
Thus APY is exact when the non-core conditional covariance is diagonal;
otherwise it is an approximation to the supplied relationship model.

APY was developed to provide a computationally economical approximation to
the inverse genomic relationship matrix for large genomic evaluations
(Misztal, 2016; Fragomeni et al., 2015). The algorithm partitions genotyped
animals into core and non-core individuals. The core contains the relationships
that are modeled jointly; each non-core animal is regressed on the core, with
its remaining conditional variance treated as independent of other non-core
animals. Thus non-core animals are not independent of the core: their
core-to-non-core relationships are retained.

Pedigree relationships are a distinct input. For pedigree-only models, the
pedigree relationship matrix \eqn{A}{} and its sparse inverse \eqn{A^{-1}}{} are
normally used directly; APY is not required. In single-step genomic BLUP,
pedigree and genomic information are combined, conventionally through
\eqn{H^{-1} = A^{-1} + \left[\begin{smallmatrix}0&0\\0&G^{-1}-A_{22}^{-1}\end{smallmatrix}\right]}{},
where genotyped animals are ordered last and \eqn{A_{22}}{} is their pedigree
relationship block. APY approximates the genomic term \eqn{G^{-1}}{} in that
correction; it does not replace the pedigree term or construct \eqn{H^{-1}}{}.
This function returns only the APY genomic precision. Do not use it as the
complete single-step precision for a population that also contains
non-genotyped animals.

This function constructs APY from an already available relationship matrix. It
does not avoid the cost of constructing a full dense relationship matrix from
markers, and it does not change the solver algorithm. Marker-based block
construction and single-step matrix assembly are not implemented here.
\end{Details}
%
\begin{Value}
A sparse symmetric APY precision matrix with individual dimnames and
\code{attr(x, "inverse") = TRUE}, ready for \code{Gu}. Metadata describing
the core and conditional variances is stored in \code{attr(x, "APY")}. If
\code{return.details = TRUE}, returns the diagnostic list described above.
\end{Value}
%
\begin{References}\relax
Misztal I (2016). Inexpensive Computation of the Inverse of the Genomic
Relationship Matrix in Populations with Small Effective Population Size.
\emph{Genetics}, 202(2), 401--409. doi:10.1534/genetics.115.182089.

Fragomeni BO, Lourenco DA, Tsuruta S, Masuda Y, Aguilar I, Legarra A, Lawlor
TJ, Misztal I (2015). Use of genomic recursions and algorithm for proven and
young animals for single-step genomic BLUP analyses--a simulation study.
\emph{Journal of Animal Breeding and Genetics}, 132(5), 340--345.
doi:10.1111/jbg.12161.
\end{References}
%
\begin{SeeAlso}\relax
\code{\LinkA{mmes}{mmes}}, \code{\LinkA{A.mat}{A.mat}}, \code{\LinkA{H.mat}{H.mat}}
\end{SeeAlso}
%
\begin{Examples}
\begin{ExampleCode}
## Marker-derived GBLUP for genotyped animals.
set.seed(2)
ids <- paste0("animal", seq_len(12))
markers <- matrix(sample(c(-1, 0, 1), 12 * 80, replace = TRUE), nrow = 12,
                  dimnames = list(ids, NULL))
G <- A.mat(markers)

## Blend slightly with the identity so this small illustrative G is positive
## definite, then construct APY precision for a six-animal core.
G <- 0.95 * G + 0.05 * diag(nrow(G))
dimnames(G) <- list(ids, ids)
Ginv.apy <- APY(G, core = ids[seq_len(6)])

phenotypes <- data.frame(
  id = factor(rep(ids, each = 3), levels = ids),
  y = rnorm(length(ids) * 3)
)
fit <- mmes(y ~ 1,
  random = ~vsm(ism(id), Gu = Ginv.apy),
  rcov = ~units,
  data = phenotypes,
  control = mmes_control(diagnostics = list(verbose = FALSE)))

## In a single-step evaluation, this APY approximation supplies G^{-1} for
## the genotyped block correction; it is not itself the complete H^{-1}.
\end{ExampleCode}
\end{Examples}
\inputencoding{utf8}
\HeaderA{ar1m}{First-order autoregressive covariance structure}{ar1m}
%
\begin{Description}\relax
First-order autoregressive covariance structure. The returned dimensionless covariance shape is intended for use inside \code{vsm}, which supplies the single overall variance scale.
\end{Description}
%
\begin{Usage}
\begin{verbatim}
ar1m(x, rho = 0.30, fixed = FALSE,
     variance = c("homogeneous", "heterogeneous"), values = NULL)
\end{verbatim}
\end{Usage}
%
\begin{Arguments}
\begin{ldescription}
\item[\code{x}] Ordered factor or design defining q ordered covariance levels.
\item[\code{rho}] Starting AR(1) correlation in (-1,1).
\item[\code{fixed}] For homogeneous variance, a logical indicating whether rho is fixed. For heterogeneous variance, a logical vector of length q: rho followed by q-1 variance-ratio parameters.
\item[\code{variance}] Whether the AR correlation has homogeneous or heterogeneous marginal variances.
\item[\code{values}] Optional q positive starting variances when \code{variance = "heterogeneous"}.
\end{ldescription}
\end{Arguments}
%
\begin{Details}\relax
The homogeneous AR(1) shape is \deqn{K_{ij}=\rho^{|i-j|}.}{} Heterogeneous mode uses \deqn{K=D R(\rho)D,}{} with the first variance fixed as the relative-scale reference. The working correlation coordinate is \eqn{\eta=\operatorname{atanh}(\rho)}{}, ensuring \eqn{|\rho|<1}{}.

Factor levels define the lag order, so numeric-looking levels must be in increasing numeric order; \code{ar1m()} warns otherwise. Inside \code{vsm()} one correlation is shared by every level of the other Kronecker factors. For a separate correlation per trial, wrap the \code{vsm()} term in \code{\LinkA{dsumm}{dsumm}}.
\end{Details}
%
\begin{Value}
A list containing the incidence/design matrix in \code{Z} and a compiled CovarianceFactor v2 descriptor in \code{covFactor}, for use inside \code{vsm}.
\end{Value}
%
\begin{SeeAlso}\relax
\code{\LinkA{ar2m}{ar2m}}, \code{\LinkA{ar3m}{ar3m}}, \code{\LinkA{toeplitzm}{toeplitzm}}, \code{\LinkA{vsm}{vsm}},
\code{\LinkA{dsumm}{dsumm}}.
\end{SeeAlso}
%
\begin{Examples}
\begin{ExampleCode}
## Not run: 
vsm(ar1m(time), ism(id))

# Separable row-by-range residual shared across trials:
# rcov = ~ vsm(dsm(trial), ar1m(range), ar1m(row), ism(units))
# Trial-specific correlations and variances:
# rcov = ~ dsumm(vsm(ar1m(range), ar1m(row), ism(units)), by=trial)

## End(Not run)
\end{ExampleCode}
\end{Examples}
\inputencoding{utf8}
\HeaderA{ar2m}{Second-order autoregressive covariance structure}{ar2m}
%
\begin{Description}\relax
Second-order autoregressive covariance structure. The returned dimensionless covariance shape is intended for use inside \code{vsm}, which supplies the single overall variance scale.
\end{Description}
%
\begin{Usage}
\begin{verbatim}
ar2m(x, pacf = c(0.20, 0.10), fixed = NULL,
     variance = c("homogeneous", "heterogeneous"), values = NULL)
\end{verbatim}
\end{Usage}
%
\begin{Arguments}
\begin{ldescription}
\item[\code{x}] Ordered factor defining covariance levels.
\item[\code{pacf}] Length-two vector of starting partial autocorrelations, each strictly between -1 and 1.
\item[\code{fixed}] For homogeneous variance, a logical vector of length two. For heterogeneous variance, a logical vector of length q+1: two PACFs followed by q-1 variance-ratio parameters.
\item[\code{variance}] Whether the AR correlation has homogeneous or heterogeneous marginal variances.
\item[\code{values}] Optional q positive starting variances when \code{variance = "heterogeneous"}.
\end{ldescription}
\end{Arguments}
%
\begin{Details}\relax
A stationary AR(2) covariance is parameterized through reflection coefficients / partial autocorrelations. Each PACF is mapped from an unconstrained working coordinate by \code{tanh}, then converted to stable AR coefficients by the Levinson/Durbin recursion. Heterogeneous mode uses \deqn{K=D R D}{} and estimates q-1 positive variance ratios in addition to the PACFs.
\end{Details}
%
\begin{Value}
A list containing the incidence/design matrix in \code{Z} and a compiled CovarianceFactor v2 descriptor in \code{covFactor}, for use inside \code{vsm}.
\end{Value}
%
\begin{SeeAlso}\relax
\code{\LinkA{ar1m}{ar1m}}, \code{\LinkA{ar3m}{ar3m}}, \code{\LinkA{toeplitzm}{toeplitzm}}, \code{\LinkA{vsm}{vsm}}.
\end{SeeAlso}
%
\begin{Examples}
\begin{ExampleCode}
## Not run: 
vsm(ar2m(time), ism(id))

## End(Not run)
\end{ExampleCode}
\end{Examples}
\inputencoding{utf8}
\HeaderA{ar3m}{Third-order autoregressive covariance structure}{ar3m}
%
\begin{Description}\relax
Third-order autoregressive covariance structure. The returned dimensionless covariance shape is intended for use inside \code{vsm}, which supplies the single overall variance scale.
\end{Description}
%
\begin{Usage}
\begin{verbatim}
ar3m(x, pacf = c(0.20, 0.10, 0.05), fixed = NULL,
     variance = c("homogeneous", "heterogeneous"), values = NULL)
\end{verbatim}
\end{Usage}
%
\begin{Arguments}
\begin{ldescription}
\item[\code{x}] Ordered factor defining covariance levels.
\item[\code{pacf}] Length-three vector of starting partial autocorrelations, each strictly between -1 and 1.
\item[\code{fixed}] For homogeneous variance, a logical vector of length three. For heterogeneous variance, a logical vector of length q+2: three PACFs followed by q-1 variance-ratio parameters.
\item[\code{variance}] Whether the AR correlation has homogeneous or heterogeneous marginal variances.
\item[\code{values}] Optional q positive starting variances when \code{variance = "heterogeneous"}.
\end{ldescription}
\end{Arguments}
%
\begin{Details}\relax
A stationary AR(3) covariance is parameterized through three partial autocorrelations. The PACF parameterization guarantees stationarity while allowing unconstrained working coordinates. Heterogeneous mode uses \deqn{K=D R D}{} and estimates q-1 positive variance ratios in addition to the PACFs.
\end{Details}
%
\begin{Value}
A list containing the incidence/design matrix in \code{Z} and a compiled CovarianceFactor v2 descriptor in \code{covFactor}, for use inside \code{vsm}.
\end{Value}
%
\begin{SeeAlso}\relax
\code{\LinkA{ar1m}{ar1m}}, \code{\LinkA{ar2m}{ar2m}}, \code{\LinkA{toeplitzm}{toeplitzm}}, \code{\LinkA{vsm}{vsm}}.
\end{SeeAlso}
%
\begin{Examples}
\begin{ExampleCode}
## Not run: 
vsm(ar3m(time), ism(id))

## End(Not run)
\end{ExampleCode}
\end{Examples}
\inputencoding{utf8}
\HeaderA{atm}{Selected-level diagonal covariance structure}{atm}
%
\begin{Description}\relax
Selected-level diagonal covariance structure. The returned dimensionless covariance shape is intended for use inside \code{vsm}, which supplies the single overall variance scale.
\end{Description}
%
\begin{Usage}
\begin{verbatim}
atm(x, levs, values = NULL, fixed = NULL)
\end{verbatim}
\end{Usage}
%
\begin{Arguments}
\begin{ldescription}
\item[\code{x}] Variable defining all available levels.
\item[\code{levs}] Character names or numeric positions of the levels retained in this covariance factor.
\item[\code{values}] Optional positive starting variances for the selected levels.
\item[\code{fixed}] Logical vector of length length(levs)-1 controlling the selected variance ratios.
\end{ldescription}
\end{Arguments}
%
\begin{Details}\relax
\code{atm} is a selected-level version of \code{dsm}.  The design is restricted to \code{levs}; observations outside those levels receive zeros for this factor.  The first selected level is the scale reference and the remaining selected levels are positive variance ratios.  This representation is most naturally used for random effects; residual factors must still assign exactly one covariance-product coordinate to every observation.
\end{Details}
%
\begin{Value}
A list containing the incidence/design matrix in \code{Z} and a compiled CovarianceFactor v2 descriptor in \code{covFactor}, for use inside \code{vsm}.
\end{Value}
%
\begin{SeeAlso}\relax
\code{\LinkA{dsm}{dsm}}, \code{\LinkA{vsm}{vsm}}, \code{\LinkA{mmes}{mmes}}.
\end{SeeAlso}
%
\begin{Examples}
\begin{ExampleCode}
## Not run: 
vsm(atm(environment, c("E1","E3")), ism(genotype))

## End(Not run)
\end{ExampleCode}
\end{Examples}
\inputencoding{utf8}
\HeaderA{atr}{Selected-level diagonal covariance structure for \code{mmer} and \code{vsr}}{atr}
%
\begin{Description}\relax
Creates a diagonal covariance-component selector for specified columns or levels of a covariance dimension. It is intended for the \code{mmer} covariance-model interface, typically inside \code{\LinkA{vsr}{vsr}}.
\end{Description}
%
\begin{Usage}
\begin{verbatim}
atr(x, levs)
\end{verbatim}
\end{Usage}
%
\begin{Arguments}
\begin{ldescription}
\item[\code{x}] 
A factor, character vector, numeric vector, or design/incidence matrix defining the covariance dimension. Factor and character inputs are expanded to incidence columns. A matrix is used directly.

\item[\code{levs}] 
Character vector identifying the columns or levels whose diagonal entries are selected. If omitted, all available column names are selected.

\end{ldescription}
\end{Arguments}
%
\begin{Details}\relax
\code{atr()} constructs a design matrix \eqn{Z}{} together with a diagonal \code{thetaC} selector. Let \eqn{q}{} be the number of columns of \eqn{Z}{}, and define
\deqn{
a_j =
\begin{cases}
1, & \text{if level }j\text{ is included in \code{levs}},\\
0, & \text{otherwise}.
\end{cases}
}{}
The returned constraint matrix is
\deqn{\mathrm{thetaC}=\mathrm{diag}(a_1,\ldots,a_q).}{}

Thus \code{atr()} is useful when a covariance component is to be associated only with a specified subset of levels or design columns. Levels not selected in \code{levs} receive zero on the corresponding diagonal of \code{thetaC}.

This function belongs to the \code{mmer}/\code{vsr} covariance-structure interface. The related \code{\LinkA{atm}{atm}} function belongs to the newer \code{mmes}/\code{vsm} CovarianceFactor interface and uses a different parameterization.

For a matrix input, its column names define the available levels. For factor or character input, the incidence-matrix column names define the levels. If \code{levs} is missing, all available columns are selected.
\end{Details}
%
\begin{Value}
A list with components:
\begin{itemize}

\item{} \code{Z}: the design/incidence representation of \code{x}.
\item{} \code{thetaC}: a diagonal zero/one matrix selecting the levels specified by \code{levs}.

\end{itemize}

\end{Value}
%
\begin{SeeAlso}\relax
\code{\LinkA{vsr}{vsr}}, \code{\LinkA{dsr}{dsr}}, \code{\LinkA{usr}{usr}}, \code{\LinkA{csr}{csr}},
\code{\LinkA{atm}{atm}}, \code{\LinkA{mmer}{mmer}}
\end{SeeAlso}
%
\begin{Examples}
\begin{ExampleCode}
# Select levels B and C
A <- atr(factor(c("A","B","A","C")), levs = c("B","C"))
A$Z
A$thetaC
\end{ExampleCode}
\end{Examples}
\inputencoding{utf8}
\HeaderA{binm}{Binomial proportion response with its number of trials}{binm}
\keyword{models}{binm}
%
\begin{Description}\relax
Builds a single proportion response for binomial \code{\LinkA{mmes}{mmes}} fits from
counts of successes and failures (or trials). The number of trials travels with
the proportion as an attribute and is used as the binomial prior weight, so the
single-response interface of \code{mmes} is preserved.
\end{Description}
%
\begin{Usage}
\begin{verbatim}
binm(successes, failures = NULL, trials = NULL)
\end{verbatim}
\end{Usage}
%
\begin{Arguments}
\begin{ldescription}
\item[\code{successes}] Numeric vector of successes.
\item[\code{failures}] Numeric vector of failures. Provide either \code{failures} or \code{trials}.
\item[\code{trials}] Numeric vector with the number of trials.
\end{ldescription}
\end{Arguments}
%
\begin{Details}\relax
The result is \code{successes/trials} with class \code{"binm"} and attribute
\code{"trials"}. Subsetting keeps the trials aligned, so the vector can be stored
as a data-frame column. Missing counts give a missing response, handled by the
\code{naMethodY} rules of \code{mmes}. Fitting a proportion without \code{binm}
treats every value as one trial, which is wrong unless all trials equal one.
\end{Details}
%
\begin{Value}
A numeric vector of proportions of class \code{"binm"}.
\end{Value}
%
\begin{SeeAlso}\relax
\code{\LinkA{mmes}{mmes}}
\end{SeeAlso}
%
\begin{Examples}
\begin{ExampleCode}
set.seed(1)
d <- data.frame(x = rnorm(60), g = factor(rep(1:12, each = 5)))
d$n <- sample(5:10, 60, TRUE)
d$s <- rbinom(60, d$n, plogis(0.3 * d$x))
fit <- mmes(binm(s, trials = n) ~ x,
  random = ~ g,
  data = d,
  family = binomial(),
  control = mmes_control(update = list(blendWeight = 0),
  diagnostics = list(verbose = FALSE)))
fit$b
\end{ExampleCode}
\end{Examples}
\inputencoding{utf8}
\HeaderA{car}{Proper conditional autoregressive spatial covariance structure}{car}
%
\begin{Description}\relax
Proper conditional autoregressive spatial covariance structure. The returned dimensionless covariance shape is intended for use inside \code{vsm}, which supplies the single overall variance scale.
\end{Description}
%
\begin{Usage}
\begin{verbatim}
car(x, W, rho = 0.10, fixed = FALSE)
\end{verbatim}
\end{Usage}
%
\begin{Arguments}
\begin{ldescription}
\item[\code{x}] Factor or design defining q spatial levels.
\item[\code{W}] Symmetric nonnegative adjacency/weights matrix with zero diagonal, aligned to the levels of x.
\item[\code{rho}] Starting proper-CAR dependence parameter.
\item[\code{fixed}] Logical indicating whether rho is fixed.
\end{ldescription}
\end{Arguments}
%
\begin{Details}\relax
Let \eqn{D=\mathrm{diag}(W\mathbf{1})}{}.  The proper CAR precision and covariance are \deqn{Q=D-\rho W,\qquad M=Q^{-1},\qquad K=M/M_{11}.}{} The constructor requires positive row sums, so isolated levels are not allowed.  The admissible open interval for \code{rho} is obtained from the eigenvalues of \eqn{D^{-1/2}WD^{-1/2}}{} and is enforced through a bounded-logit working coordinate.  This is a proper, nonsingular CAR model; an intrinsic singular CAR is not used by the current precision-based Henderson implementation.  The covariance derivative is analytic.
\end{Details}
%
\begin{Value}
A list containing the incidence/design matrix in \code{Z} and a compiled CovarianceFactor v2 descriptor in \code{covFactor}, for use inside \code{vsm}.
\end{Value}
%
\begin{SeeAlso}\relax
\code{\LinkA{sar}{sar}}, \code{\LinkA{maternm}{maternm}}, \code{\LinkA{vsm}{vsm}}.
\end{SeeAlso}
%
\begin{Examples}
\begin{ExampleCode}
## Not run: 
vsm(car(location, W), ism(genotype))

## End(Not run)
\end{ExampleCode}
\end{Examples}
\inputencoding{utf8}
\HeaderA{coef.mmes}{coef form a GLMM fitted with mmes}{coef.mmes}
\aliasA{coef.mmer}{coef.mmes}{coef.mmer}
\keyword{models}{coef.mmes}
%
\begin{Description}\relax
\code{coef} method for class \code{"mmes"}.
\end{Description}
%
\begin{Usage}
\begin{verbatim}
## S3 method for class 'mmes'
coef(object, ...)
\end{verbatim}
\end{Usage}
%
\begin{Arguments}
\begin{ldescription}
\item[\code{object}] an object of class \code{"mmes"}
\item[\code{...}] Further arguments to be passed
\end{ldescription}
\end{Arguments}
%
\begin{Value}
vector of coef

\end{Value}
%
\begin{Author}\relax
Giovanny Covarrubias
\end{Author}
%
\begin{SeeAlso}\relax
\code{\LinkA{coef}{coef}}, \code{\LinkA{mmes}{mmes}} 
\end{SeeAlso}
\inputencoding{utf8}
\HeaderA{corgm}{General positive-definite correlation structure}{corgm}
%
\begin{Description}\relax
General positive-definite correlation structure. The returned dimensionless covariance shape is intended for use inside \code{vsm}, which supplies the single overall variance scale.
\end{Description}
%
\begin{Usage}
\begin{verbatim}
corgm(x, theta = NULL, fixed = NULL,
      variance = c("homogeneous", "heterogeneous"), values = NULL)
\end{verbatim}
\end{Usage}
%
\begin{Arguments}
\begin{ldescription}
\item[\code{x}] Variable defining q covariance levels.
\item[\code{theta}] Optional positive-definite q by q starting covariance or correlation matrix.
With \code{variance="heterogeneous"} its diagonal is used as starting
variances when \code{values} is not given.
\item[\code{fixed}] Logical vector flagging parameters held at their starting value:
length q(q-1)/2 (correlations) for the homogeneous mode and
q(q-1)/2 + q-1 (correlations, then variance ratios of levels 2..q) for the
heterogeneous mode.
\item[\code{variance}] \code{"homogeneous"} (default) gives a pure correlation
matrix scaled by the single variance of \code{vsm}.
\code{"heterogeneous"} adds one variance per level.
\item[\code{values}] Optional positive starting variances (length q) for the
heterogeneous mode. Only their ratios to the first level matter; the first
level's variance is the \code{vsm} scale. When omitted, the starting ratios
are taken from the data.
\end{ldescription}
\end{Arguments}
%
\begin{Details}\relax
An unrestricted SPD correlation matrix is represented using a unit-diagonal lower factor \eqn{A}{}.  The intermediate matrix \eqn{S=AA^{\mathsf T}}{} is standardized to correlation scale, \deqn{K_{ij}=S_{ij}/\sqrt{S_{ii}S_{jj}}.}{} The unrestricted working parameters are the strict lower-triangular entries of \eqn{A}{}.

With \code{variance="heterogeneous"} the shape is \eqn{K = DRD}{} with
\eqn{D=\mathrm{diag}(1, \sqrt{r_2}, \ldots, \sqrt{r_q})}{}, where
\eqn{r_l=\sigma^2_l/\sigma^2_1 = \exp(\phi_l)}{} are variance ratios relative to
the first level; the \code{vsm} scale is \eqn{\sigma^2_1}{}. The model is the
same as \code{\LinkA{usm}{usm}} but parameterized as correlations and variances,
which allows individual variances to be fixed. This is what mixed-family
multi-trait models (\code{\LinkA{familym}{familym}}) need: the residual variance of a
binomial or Poisson trait is fixed to one while the correlations and the
variances of Gaussian traits are estimated. Reported parameters
(\code{\LinkA{covparams\_mmes}{covparams.Rul.mmes}}) are the correlations and the variances of each
level.
\end{Details}
%
\begin{Value}
A list containing the incidence/design matrix in \code{Z} and a compiled CovarianceFactor v2 descriptor in \code{covFactor}, for use inside \code{vsm}.
\end{Value}
%
\begin{SeeAlso}\relax
\code{\LinkA{csm}{csm}}, \code{\LinkA{usm}{usm}}, \code{\LinkA{vsm}{vsm}},
\code{\LinkA{covmatrix\_mmes}{covmatrix.Rul.mmes}}, \code{\LinkA{familym}{familym}}.
\end{SeeAlso}
%
\begin{Examples}
\begin{ExampleCode}
data(DT_example, package = "enhancer")
DT <- DT_example
# genetic correlations between environments with a homogeneous variance
m1 <- mmes(Yield ~ Env,
  random = ~ vsm(corgm(Env), ism(Name)),
  rcov = ~ units,
  data = DT,
  control = mmes_control(diagnostics = list(verbose = FALSE)))
covparams_mmes(m1, 1)

# heterogeneous variances (equivalent to usm, parameterized as corr/var)
m2 <- mmes(Yield ~ Env,
  random = ~ vsm(corgm(Env, variance = "heterogeneous"), ism(Name)),
  rcov = ~ units,
  data = DT,
  control = mmes_control(diagnostics = list(verbose = FALSE)))
covparams_mmes(m2, 1)
covmatrix_mmes(m2, 1)$correlation
\end{ExampleCode}
\end{Examples}
\inputencoding{utf8}
\HeaderA{corImputation}{Imputing a matrix using correlations}{corImputation}
%
\begin{Description}\relax
corImputation imputes missing data based on the correlation that exists between row levels.
\end{Description}
%
\begin{Usage}
\begin{verbatim}
  corImputation(wide, Gu=NULL, nearest=10, roundR=FALSE)
\end{verbatim}
\end{Usage}
%
\begin{Arguments}
\begin{ldescription}
\item[\code{wide}] numeric matrix with individuals in rows and time variable in columns (e.g., environments, genetic markers, etc.).
\item[\code{Gu}] optional correlation matrix between the individuals or row levels. If NULL it will be computed as the correlation of t(wide).
\item[\code{nearest}] integer value describing how many nearest neighbours (the ones showing the highest correlation) should be used to average and return the imputed value.
\item[\code{roundR}] a TRUE/FALSE statement describing if the average result should be rounded or not. This may be specifically useful for categorical data in the form of numbers (e.g., -1,0,1).
\end{ldescription}
\end{Arguments}
%
\begin{Value}
\begin{description}


\item[\$res]  a list with the imputed matrix and the original matrix.


\end{description}

\end{Value}
%
\begin{Author}\relax
Giovanny Covarrubias-Pazaran
\end{Author}
%
\begin{References}\relax
Covarrubias-Pazaran G (2016) Genome assisted prediction of quantitative traits using the R package sommer. PLoS ONE 11(6): doi:10.1371/journal.pone.0156744 

\end{References}
%
\begin{Examples}
\begin{ExampleCode}

####################################
### imputing genotype data example
####################################
# data(DT_cpdata, package="enhancer")
# X <- GT_cpdata
# # add missing data
# v <- sample(1:length(X), 500)
# Xna <- X
# Xna[v]<- NA
# ## impute (can take some time)
# Y <- corImputation(wide=Xna, Gu=NULL, nearest=20, roundR=TRUE) 
# cm <- table(Y$imputed[v],X[v])
# ## calculate accuracy
# sum(diag(cm))/length(v)
####################################
### imputing phenotypic data example
####################################
# data(DT_h2, package="enhancer")
# X <- reshape(DT_h2[,c("Name","Env","y")], direction = "wide", idvar = "Name",
#                 timevar = "Env", v.names = "y", sep= "_")
# rownames(X) <- X$Name
# X <- as.matrix(X[,-1])
# head(X)
# # add missing data
# v <- sample(1:length(X), 50)
# Xna <- X
# Xna[v]<- NA
# ## impute
# Y <- corImputation(wide=Xna, Gu=NULL, nearest=20, roundR=TRUE)
# plot(y=Y$imputed[v],x=X[v], xlab="true",ylab="predicted")
# cor(Y$imputed[v],X[v], use = "complete.obs")


\end{ExampleCode}
\end{Examples}
\inputencoding{utf8}
\HeaderA{covm}{Covariance Between Two Random Effects}{covm}
%
\begin{Description}\relax
\code{covm} combines two random-effect structures created with
\code{\LinkA{vsm}{vsm}} into a single correlated random-effect structure.
It fits an unstructured \eqn{2 \times 2}{} covariance matrix between the
two random effects while allowing them to have different incidence
matrices. The two effects must act on the same coefficient space and
use the same relationship or precision matrix. The resulting structure
can be fitted with the \code{\LinkA{mmes}{mmes}} solver.
\end{Description}
%
\begin{Usage}
\begin{verbatim}
covm(ran1, ran2, thetaC = NULL, theta = NULL,
fixed = NULL, fixedSigma2 = FALSE,
labels = c("ran1", "ran2"), tol = 1e-10)
\end{verbatim}
\end{Usage}
%
\begin{Arguments}
\begin{ldescription}

\item[\code{ran1}] 
A random-effect structure returned by \code{\LinkA{vsm}{vsm}} for the
first random effect. Currently, \code{covm} requires a simple
\code{vsm} structure with one covariance-product coordinate, for
example \code{vsm(ism(focal))} or
\code{vsm(ism(focal), Gu = Ai)}.


\item[\code{ran2}] 
A random-effect structure returned by \code{\LinkA{vsm}{vsm}} for the
second random effect. It must have the same coefficient dimension
and coefficient levels as \code{ran1}, although its incidence
matrix may differ.


\item[\code{thetaC}] 
Deprecated and not supported by the current CovarianceFactor
parameterization. The former cell-wise constraint matrix cannot in
general be translated exactly into the normalized-Cholesky
parameterization used by \code{covm}. Use \code{fixedSigma2} and
\code{fixed} instead.


\item[\code{theta}] 
An optional symmetric \eqn{2 \times 2}{} positive-definite matrix
containing initial values for the variance-covariance matrix of the
two random effects. Values are supplied directly on the natural
variance-covariance scale.

```
The diagonal elements define the initial variances of the two random
effects and the off-diagonal element defines their initial covariance.

If \code{theta = NULL}, the default starting matrix is

\begin{alltt}
```

diag(2) * 0.15 + matrix(0.015, 2, 2)
\end{alltt}


```
giving initial variances of 0.165 and an initial covariance of 0.015.
```



\item[\code{fixed}] 
An optional logical vector of length two indicating whether the two
normalized-Cholesky coordinates describing the covariance structure
should be fixed at their starting values. The coordinates correspond
to the off-diagonal Cholesky element and the second diagonal Cholesky
element, respectively. The default is \code{c(FALSE, FALSE)}, so both
are estimated.


\item[\code{fixedSigma2}] 
Logical value indicating whether the variance scale associated with
the first random effect should be fixed at its starting value.
The default is \code{FALSE}.


\item[\code{labels}] 
Character vector of length two giving labels for the two random
effects in the covariance descriptor. The default is
\code{c("ran1", "ran2")}. The labels must be different and non-empty.


\item[\code{tol}] 
Numerical tolerance used when checking positive definiteness and
whether the relationship/precision matrices supplied by the two
random effects are equal. The default is \code{1e-10}.

\end{ldescription}
\end{Arguments}
%
\begin{Details}\relax

\code{covm} is intended for models in which two distinct random effects
are expected to be correlated. A common example is the joint modeling
of direct and indirect genetic effects.

If the coefficient vector for the first random effect is \eqn{u_1}{} and
that for the second is \eqn{u_2}{}, \code{covm} represents their covariance
structure as

\deqn{
\mathrm{Var}([u_1^\prime, u_2^\prime]^\prime)
= \sigma^2 K_{\mathrm{effect}} \otimes A,
}{}

where \eqn{A}{} represents the covariance structure associated with the
coefficient levels, and \eqn{K_{\mathrm{effect}}}{} is a
\eqn{2 \times 2}{} unstructured covariance shape describing the
relationship between the two random effects.

Internally, \code{covm} uses the same CovarianceFactor interface as
\code{\LinkA{vsm}{vsm}}. The \eqn{2 \times 2}{} effect covariance is represented
using a normalized-Cholesky parameterization. This guarantees a
positive-definite covariance matrix during optimization.

The variance of the first random effect is represented by the overall
scale parameter \code{sigma2}. The remaining variance ratio and
covariance are represented through the normalized-Cholesky factor.

The two random effects may have different incidence matrices. However,
they must refer to identical coefficient levels in the same order and
must use the same relationship/precision matrix.

When a relationship matrix is supplied through \code{Gu}, the
Henderson solver requires it to be an inverse/precision matrix with
\code{attr(Gu, "inverse") = TRUE}. For example:

\begin{alltt}
attr(Ai, "inverse") <- TRUE

covm(
vsm(ism(focal), Gu = Ai),
vsm(ism(neighbour), Gu = Ai)
)
\end{alltt}


The current implementation combines simple \code{vsm} random effects
only. Each input must contain one covariance-product coordinate. More
complex covariance-factor products should therefore not currently be
supplied independently inside \code{ran1} and \code{ran2}.
\end{Details}
%
\begin{Value}

A list describing a correlated random-effect structure compatible with
the CovarianceFactor-v2 interface used by \code{\LinkA{vsm}{vsm}} and
\code{\LinkA{mmes}{mmes}}. The main elements are:

\begin{description}


```
\item[\code{Z}] 
A list containing the two incidence matrices, one for each random
effect.


\item[\code{Gu}] 
The common sparse relationship precision matrix associated with
the coefficient levels. It is stored with
\code{attr(Gu, "inverse") = TRUE} for use by the Henderson solver.


\item[\code{covStruct}] 
A CovarianceFactor-v2 descriptor containing the overall variance
scale and the normalized-Cholesky representation of the
\eqn{2 \times 2}{} unstructured covariance between the two random
effects.


\item[\code{residualLocalIndex}] 
\code{NULL}, since \code{covm} describes a random-effect structure
rather than a residual covariance structure.


\item[\code{productDesign}] 
\code{NULL} for the current implementation.


\item[\code{partitionsR}] 
\code{NULL} for the current implementation.


\item[\code{covm}] 
Logical value equal to \code{TRUE}, identifying the object as one
constructed by \code{covm}.


\item[\code{covm\_labels}] 
The two labels used for the correlated random effects.

```


\end{description}

\end{Value}
%
\begin{Author}\relax
Giovanny Covarrubias-Pazaran
\end{Author}
%
\begin{References}\relax

Covarrubias-Pazaran G (2016).
Genome assisted prediction of quantitative traits using the R package
sommer. PLoS ONE 11(6).
\Rhref{https://doi.org/10.1371/journal.pone.0156744}{doi:10.1371\slash{}journal.pone.0156744}

Bijma, P. (2014).
The quantitative genetics of indirect genetic effects: a selective
review of modelling issues.
Heredity, 112(1), 61--69.
\end{References}
%
\begin{SeeAlso}\relax
\code{\LinkA{vsm}{vsm}} for constructing random-effect covariance structures
and \code{\LinkA{mmes}{mmes}} for fitting mixed models using the Henderson
solver.
\end{SeeAlso}
%
\begin{Examples}
\begin{ExampleCode}

data(DT_ige, package = "enhancer")
DT <- DT_ige

## Correlate two random effects with identity relationship matrices

covRes <- with(
DT,
covm(
vsm(ism(focal)),
vsm(ism(neighbour))
)
)

str(covRes)

## Custom initial variance-covariance matrix

covRes2 <- with(
DT,
covm(
vsm(ism(focal)),
vsm(ism(neighbour)),
theta = matrix(
c(0.5, 0.1,
0.1, 0.3),
2, 2
)
)
)

## Relationship precision matrix

## Ai must have matching row and column names and be marked as inverse.

##

## attr(Ai, "inverse") <- TRUE

##

## covRes3 <- with(

## DT,

## covm(

## vsm(ism(focal), Gu = Ai),

## vsm(ism(neighbour), Gu = Ai)

## )

## )

## See the DT_ige help page for a complete model-fitting example.

\end{ExampleCode}
\end{Examples}
\inputencoding{utf8}
\HeaderA{covmatrix\_mmes}{Covariance and correlation matrix of a fitted term with standard errors}{covmatrix.Rul.mmes}
\keyword{models}{covmatrix\_mmes}
%
\begin{Description}\relax
Returns the full covariance matrix of a random or residual term of a fitted
\code{\LinkA{mmes}{mmes}} model over its covariance coordinates (for example the
genetic trait-by-trait covariance matrix of a multi-trait model), together
with the implied correlation matrix and delta-method standard errors of every
element.
\end{Description}
%
\begin{Usage}
\begin{verbatim}
covmatrix_mmes(object, term, se = TRUE, max.dim = 500L, rel_step = 1e-5)
\end{verbatim}
\end{Usage}
%
\begin{Arguments}
\begin{ldescription}
\item[\code{object}] A fitted \code{\LinkA{mmes}{mmes}} model.
\item[\code{term}] Name (as in \code{names(object\$covStruct)}) or position of the
term. The residual term is the last one.
\item[\code{se}] If \code{TRUE}, add standard errors.
\item[\code{max.dim}] Refuse terms whose covariance coordinate space is larger than
this (the matrices are dense).
\item[\code{rel\_step}] Relative step of the central differences used for the
Jacobian.
\end{ldescription}
\end{Arguments}
%
\begin{Details}\relax
The matrix is \eqn{\sigma^2 K_1 \otimes \cdots \otimes K_m}{} built from all the
covariance factors of the \code{vsm} term (the relationship matrix
\code{Gu} is not included), so for \code{vsm(usm(trait), ism(id))} it is the
trait covariance matrix, and for \code{vsm(dsm(env), usm(trait), ism(id))} the
environment-by-trait matrix. Row and column names are the coordinate levels.

Standard errors use the delta method,
\eqn{\mathrm{Var}(g(\hat\theta)) \approx J V J^\top}{}, where \eqn{V}{} is
\code{object\$theta\_se} (the covariance of the reported parameters, as used by
\code{\LinkA{covparams\_mmes\_se}{covparams.Rul.mmes.Rul.se}}) and \eqn{J}{} the Jacobian of each element with
respect to the reported parameters. Fixed parameters contribute no variance.
Diagonal standard errors of the correlation matrix are zero. Terms split by
\code{\LinkA{dsumm}{dsumm}} sections are not supported.
\end{Details}
%
\begin{Value}
A list with \code{covariance} and \code{correlation} and, if
\code{se=TRUE}, \code{covariance.se} and \code{correlation.se}.
\end{Value}
%
\begin{SeeAlso}\relax
\code{\LinkA{covparams\_mmes}{covparams.Rul.mmes}}, \code{\LinkA{covparams\_mmes\_se}{covparams.Rul.mmes.Rul.se}},
\code{\LinkA{vpredict}{vpredict}}, \code{\LinkA{stackTraits}{stackTraits}}
\end{SeeAlso}
%
\begin{Examples}
\begin{ExampleCode}
set.seed(17)
n <- 60
DT <- data.frame(id = factor(rep(seq_len(n), each = 3)))
u <- matrix(rnorm(n * 2), n) 
DT$trait_a <- u[DT$id, 1] + rnorm(nrow(DT), sd = 0.7)
DT$trait_b <- u[DT$id, 2] + rnorm(nrow(DT), sd = 1)
DTL <- stackTraits(DT, traits = c("trait_a", "trait_b"), keep = "id")
fit <- mmes(value ~ trait,
  random = ~ vsm(usm(trait), ism(id)),
  rcov = ~ vsm(usm(trait), ism(record)),
  data = DTL,
  control = mmes_control(diagnostics = list(verbose = FALSE)))
G <- covmatrix_mmes(fit, 1)       # genetic
G$covariance
G$correlation
G$correlation.se
R <- covmatrix_mmes(fit, 2)       # residual
R$correlation
\end{ExampleCode}
\end{Examples}
\inputencoding{utf8}
\HeaderA{covparams\_mmes}{Extract native-scale covariance parameters from an mmes fit}{covparams.Rul.mmes}
%
\begin{Description}\relax
Returns covariance parameters in the native, model-facing scale defined by
each CovarianceFactor descriptor, rather than the unconstrained working
coordinates or normalized variance ratios used internally by the optimizer.
\end{Description}
%
\begin{Usage}
\begin{verbatim}
covparams_mmes(object, term = NULL)
\end{verbatim}
\end{Usage}
%
\begin{Arguments}
\begin{ldescription}
\item[\code{object}] A fitted object of class \code{"mmes"}.
\item[\code{term}] Optional covariance-term name or numeric index. By default,
all random and residual covariance terms are returned.
\end{ldescription}
\end{Arguments}
%
\begin{Details}\relax
Every CovarianceFactor descriptor carries a native reporting callback. Built-in
structures use it to report meaningful quantities such as level-specific
variances for \code{dsm()}, variances and covariances for \code{usm()}, and
variance-correlation parameters for correlation structures. Structural zeros
and internal coordinates such as log standard deviations, normalized Cholesky
coefficients, and variance ratios are not returned where the covariance model
provides a more direct native representation.

All built-in CovarianceFactor families provide native reporters, including
identity, diagonal, unstructured, AR(1--3), compound-symmetry, moving-average,
general-correlation, factor-analytic, antedependence, reduced-rank, Matern,
Toeplitz, SAR, and CAR structures. Heterogeneous AR and compound-symmetry
models report their dependence parameter together with level-specific
variances rather than an overall variance plus variance ratios.

Both fixed and estimated covariance parameters are included. For an arbitrary
Kronecker product, the product-level variance scale is absorbed by the first
covariance-shaping factor. Subsequent factors are necessarily reported on their
normalized relative scale because the product has only one identifiable overall
scale.

User-defined \code{ownm()} structures can supply a \code{native\_report}
callback. If none is supplied, their overall variance and naturally transformed
parameter values are returned.

Use \code{covparams\_mmes\_se()} when native-scale delta-method standard errors
are also required.
\end{Details}
%
\begin{Value}
A data frame with columns \code{term}, \code{factor}, \code{section},
\code{parameter}, and \code{estimate}. \code{section} identifies the
\code{\LinkA{dsumm}{dsumm}} section a parameter belongs to and is \code{NA} otherwise.
\end{Value}
%
\begin{SeeAlso}\relax
\code{\LinkA{covparams\_mmes\_se}{covparams.Rul.mmes.Rul.se}}, \code{\LinkA{mmes}{mmes}}, \code{\LinkA{vsm}{vsm}},
\code{\LinkA{ownm}{ownm}}
\end{SeeAlso}
%
\begin{Examples}
\begin{ExampleCode}
## Not run: 
fit <- mmes(Yield ~ Env,
            random = ~ vsm(dsm(Env), ism(Name)),
            rcov = ~ units, data = DT_example)
covparams_mmes(fit)
fit$covParNative

## End(Not run)
\end{ExampleCode}
\end{Examples}
\inputencoding{utf8}
\HeaderA{covparams\_mmes\_se}{Native-scale covariance parameters and standard errors}{covparams.Rul.mmes.Rul.se}
%
\begin{Description}\relax
Returns descriptor-defined native covariance parameters with delta-method
standard errors computed from the fitted covariance-parameter uncertainty
matrix.
\end{Description}
%
\begin{Usage}
\begin{verbatim}
covparams_mmes_se(object, term = NULL, rel_step = 1e-6)
\end{verbatim}
\end{Usage}
%
\begin{Arguments}
\begin{ldescription}
\item[\code{object}] A fitted object of class \code{"mmes"}.
\item[\code{term}] Optional covariance-term name or numeric index. By default,
all random and residual covariance terms are returned.
\item[\code{rel\_step}] Positive relative step used to numerically differentiate
each CovarianceFactor native reporting callback.
\end{ldescription}
\end{Arguments}
%
\begin{Details}\relax
The fitted \code{theta\_se} matrix describes uncertainty in the reported
\code{covPar} coordinates. For each covariance structure,
\code{covparams\_mmes\_se} numerically evaluates the Jacobian of its descriptor's
native reporting callback and applies the multivariate delta method,
\deqn{V_{native} = J V_{covPar} J^{\mathsf T}.}{}
Consequently, coupled transformations use the complete covariance information.
For example, the standard error of a \code{dsm()} level variance accounts for
uncertainty and covariance in both the overall variance and its variance ratio.
Likewise, \code{usm()} standard errors account for all relevant normalized
Cholesky parameters.

The transformation is available for every built-in CovarianceFactor family
supported by \code{covparams\_mmes()}; user-defined \code{ownm()} structures
use their descriptor-provided \code{native\_report} callback.

Parameters marked fixed in the covariance descriptor contribute no independent
sampling uncertainty: their rows and columns in \code{theta\_se} are excluded
before transformation. A native quantity involving both fixed and free
parameters can still have a nonzero standard error through its free parameters.
Both fixed and estimated native quantities are returned.

The same table is stored automatically in \code{object\$covParNativeSE} by
\code{mmes()}.
\end{Details}
%
\begin{Value}
A data frame containing \code{term}, \code{factor}, \code{parameter},
\code{estimate}, \code{StdError}, and \code{Zratio}. A zero standard error is
reported for a fully fixed native quantity, with \code{Zratio=NA}.
\end{Value}
%
\begin{SeeAlso}\relax
\code{\LinkA{covparams\_mmes}{covparams.Rul.mmes}}, \code{\LinkA{mmes}{mmes}}, \code{\LinkA{vsm}{vsm}}
\end{SeeAlso}
%
\begin{Examples}
\begin{ExampleCode}
## Not run: 
fit <- mmes(Yield ~ Env,
            random = ~ vsm(dsm(Env), ism(Name)),
            rcov = ~ units, data = DT_example)
covparams_mmes_se(fit)
fit$covParNativeSE

## End(Not run)
\end{ExampleCode}
\end{Examples}
\inputencoding{utf8}
\HeaderA{csm}{Compound-symmetry covariance structure}{csm}
%
\begin{Description}\relax
Compound-symmetry covariance structure. The returned dimensionless covariance shape is intended for use inside \code{vsm}, which supplies the single overall variance scale.
\end{Description}
%
\begin{Usage}
\begin{verbatim}
csm(x, rho = 0.10, fixed = FALSE,
    variance = c("homogeneous", "heterogeneous"), values = NULL)
\end{verbatim}
\end{Usage}
%
\begin{Arguments}
\begin{ldescription}
\item[\code{x}] Variable defining q covariance levels.
\item[\code{rho}] Starting common off-diagonal correlation.
\item[\code{fixed}] For homogeneous variance, a logical indicating whether rho is fixed. For heterogeneous variance, a logical vector of length q: rho followed by q-1 variance-ratio parameters.
\item[\code{variance}] Whether the correlation has homogeneous or heterogeneous marginal variances.
\item[\code{values}] Optional q positive starting variances when \code{variance = "heterogeneous"}.
\end{ldescription}
\end{Arguments}
%
\begin{Details}\relax
With \code{variance = "homogeneous"}, \code{csm} uses unit diagonal and a common off-diagonal correlation, \deqn{K_{ii}=1, \quad K_{ij}=\rho\ (i\ne j).}{} With \code{variance = "heterogeneous"}, it uses \deqn{K=D C(\rho)D,}{} where the first variance is the reference and the remaining q-1 variances are estimated as positive ratios. Positive definiteness requires \eqn{-1/(q-1)<\rho<1}{}. The product-level variance remains owned by \code{vsm}.
\end{Details}
%
\begin{Value}
A list containing the incidence/design matrix in \code{Z} and a compiled CovarianceFactor v2 descriptor in \code{covFactor}, for use inside \code{vsm}.
\end{Value}
%
\begin{SeeAlso}\relax
\code{\LinkA{corgm}{corgm}}, \code{\LinkA{dsm}{dsm}}, \code{\LinkA{vsm}{vsm}}.
\end{SeeAlso}
%
\begin{Examples}
\begin{ExampleCode}
## Not run: 
vsm(csm(environment), ism(genotype))

## End(Not run)
\end{ExampleCode}
\end{Examples}
\inputencoding{utf8}
\HeaderA{csr}{User-specified covariance-component structure for \code{mmer} and \code{vsr}}{csr}
%
\begin{Description}\relax
Combines a design/incidence representation with a user-supplied covariance-component constraint matrix for use by the \code{mmer} covariance-model interface, typically inside \code{\LinkA{vsr}{vsr}}.
\end{Description}
%
\begin{Usage}
\begin{verbatim}
csr(x, mm)
\end{verbatim}
\end{Usage}
%
\begin{Arguments}
\begin{ldescription}
\item[\code{x}] 
A factor, character vector, numeric vector, or design/incidence matrix defining the covariance dimension. Factor and character inputs are expanded to incidence columns. A matrix is used directly.

\item[\code{mm}] 
A square covariance-component constraint matrix. Its dimensions must correspond to the columns of the design matrix generated from \code{x}. The matrix is returned as \code{thetaC}, with its row and column names set to the column names of the resulting design matrix.

\end{ldescription}
\end{Arguments}
%
\begin{Details}\relax
\code{csr()} provides the general constraint-matrix interface underlying custom covariance patterns for \code{mmer}/\code{vsr}. Unlike \code{\LinkA{dsr}{dsr}} and \code{\LinkA{usr}{usr}}, it does not construct a predefined diagonal or unstructured pattern. Instead, the supplied matrix \code{mm} determines the covariance-component structure through \code{thetaC}.

Conceptually, \code{thetaC} identifies which elements of the covariance matrix are associated with estimated variance-covariance components and which constraints are shared among elements. The interpretation of the entries therefore follows the \code{thetaC} convention used by \code{\LinkA{vsr}{vsr}}.

When \code{x} is a matrix it is used directly as \eqn{Z}{}. Otherwise, factors and character vectors are converted to incidence matrices, while numeric non-factor inputs are treated as a single design column.

The function assigns the column names of \eqn{Z}{} to both dimensions of \code{mm}; therefore \code{mm} should be conformable with the number of columns of the resulting design matrix.
\end{Details}
%
\begin{Value}
A list with components:
\begin{itemize}

\item{} \code{Z}: the design/incidence matrix associated with the covariance dimension.
\item{} \code{thetaC}: the user-supplied covariance-component constraint matrix, with dimensions named according to \code{Z}.

\end{itemize}

\end{Value}
%
\begin{SeeAlso}\relax
\code{\LinkA{vsr}{vsr}}, \code{\LinkA{dsr}{dsr}}, \code{\LinkA{usr}{usr}}, \code{\LinkA{atr}{atr}},
\code{\LinkA{mmer}{mmer}}
\end{SeeAlso}
%
\begin{Examples}
\begin{ExampleCode}
# A custom two-level covariance-component pattern
x <- factor(c("A","B","A","B"))
mm <- matrix(c(1, 2,
               2, 1), 2, 2, byrow = TRUE)
C <- csr(x, mm)
C$Z
C$thetaC
\end{ExampleCode}
\end{Examples}
\inputencoding{utf8}
\HeaderA{D.mat}{Dominance relationship matrix}{D.mat}
%
\begin{Description}\relax
C++ implementation of the dominance matrix. Calculates the realized dominance relationship matrix. Can help to increase the prediction accuracy when 2 conditions are met; 1) The trait has intermediate to high heritability, 2) The population contains a big number of individuals that are half or full sibs (HS \& FS).
\end{Description}
%
\begin{Usage}
\begin{verbatim}
D.mat(X,nishio=TRUE,min.MAF=0,return.imputed=FALSE)
\end{verbatim}
\end{Usage}
%
\begin{Arguments}
\begin{ldescription}
\item[\code{X}] 
Matrix (\eqn{n \times m}{}) of unphased genotypes for \eqn{n}{} lines and \eqn{m}{} biallelic markers, 
coded as \{-1,0,1\}. Fractional (imputed) and missing values (NA) are allowed.

\item[\code{nishio}] 
If TRUE, Nishio and Satoh (2014), which uses the orthogonal parameterization of Vitezica et al. (2013). Otherwise Su et al. (2012). See Details and references.

\item[\code{min.MAF}] 
Minimum minor allele frequency. The D matrix is not sensitive to rare alleles, so by default only monomorphic markers are removed.

\item[\code{return.imputed}] 
When TRUE, the imputed marker matrix is returned.  



\end{ldescription}
\end{Arguments}
%
\begin{Details}\relax

Missing values are replaced by the column (marker) mean before the matrix is computed. Markers with minor allele frequency less than or equal to \code{min.MAF} (and therefore all monomorphic markers) are removed. For each marker \eqn{k}{}, \eqn{p_k}{} is the frequency of the allele counted by the code 1, estimated as the column mean of \eqn{(X+1)/2}{}, and \eqn{q_k = 1-p_k}{}.

For the nishio method (Nishio and Satoh, 2014; Vitezica et al., 2013): genotypes are coded as \eqn{w_{ik} = -2q_k^2}{} for \eqn{x_{ik}=1}{}, \eqn{2p_kq_k}{} for \eqn{x_{ik}=0}{} and \eqn{-2p_k^2}{} for \eqn{x_{ik}=-1}{}, computed as \eqn{w_{ik} = -x_{ik}^2 + (p_k-q_k)x_{ik} + 2p_kq_k}{}. Fractional values are passed through the same formula; note that a mean-imputed value \eqn{x_{ik} = p_k-q_k}{} receives the heterozygote code \eqn{2p_kq_k}{}. Then \eqn{D = W W'/c}{}, where \eqn{c = \sum_k (2p_kq_k)^2}{}. Under Hardy-Weinberg equilibrium this coding is orthogonal to the additive coding.

For the Su method (Su et al., 2012): heterozygosity is coded as \eqn{h_{ik} = 1-|x_{ik}|}{} and centered as \eqn{m_{ik} = h_{ik} - 2p_kq_k}{}. Then \eqn{D = M M'/c}{}, where \eqn{c = \sum_k 2p_kq_k(1-2p_kq_k)}{}.

Versions of sommer earlier than 4.4.9 used a different formula for \code{nishio=TRUE}: \eqn{h_{ik}}{} centered by its column mean and divided by \eqn{\sum_k (2p_kq_k)^2}{}. That was not the Nishio and Satoh (2014) method, so results from those versions will differ.

\end{Details}
%
\begin{Value}
If return.imputed = FALSE, the \eqn{n \times n}{} dominance relationship matrix is returned.

If return.imputed = TRUE, the function returns a list containing
\begin{description}

\item[\$D] the D matrix
\item[\$X] the imputed marker matrix

\end{description}

\end{Value}
%
\begin{References}\relax

Covarrubias-Pazaran G (2016) Genome assisted prediction of quantitative traits using the R package sommer. PLoS ONE 11(6): doi:10.1371/journal.pone.0156744 

Nishio M and Satoh M. 2014. Including Dominance Effects in the Genomic BLUP Method for Genomic Evaluation. Plos One 9(1), doi:10.1371/journal.pone.0085792

Vitezica ZG, Varona L, Legarra A. 2013. On the Additive and Dominant Variance and Covariance of Individuals Within the Genomic Selection Scope. Genetics 195(4): 1223-1230. doi:10.1534/genetics.113.155176

Su G, Christensen OF, Ostersen T, Henryon M, Lund MS. 2012. Estimating Additive and Non-Additive Genetic Variances and Predicting Genetic Merits Using Genome-Wide Dense Single Nucleotide Polymorphism Markers. PLoS ONE 7(9): e45293. doi:10.1371/journal.pone.0045293

\end{References}
%
\begin{Examples}
\begin{ExampleCode}
####=========================================####
#### EXAMPLE 1
####=========================================####
####random population of 200 lines with 1000 markers
X <- matrix(rep(0,200*1000),200,1000)
for (i in 1:200) {
  X[i,] <- sample(c(-1,0,0,1), size=1000, replace=TRUE)
}

D <- D.mat(X)

\end{ExampleCode}
\end{Examples}
\inputencoding{utf8}
\HeaderA{dfToMatrix}{data frame to matrix}{dfToMatrix}
%
\begin{Description}\relax
This function takes a matrix that is in data frame format and transforms it into a matrix. Other packages that allows you to obtain an additive relationship matrix from a pedigree is the `pedigreemm` package.
\end{Description}
%
\begin{Usage}
\begin{verbatim}
  dfToMatrix(x, row="Row",column="Column",
             value="Ainverse", returnInverse=FALSE, 
             bend=1e-6)
\end{verbatim}
\end{Usage}
%
\begin{Arguments}
\begin{ldescription}
\item[\code{x}] ginv element, output from the Ainverse function.
\item[\code{row}] name of the column in x that indicates the row in the original relationship matrix.
\item[\code{column}] name of the column in x that indicates the column in the original relationship matrix.
\item[\code{value}] name of the column in x that indicates the value for a given row and column in the original relationship matrix.
\item[\code{returnInverse}] a TRUE/FALSE value indicating if the inverse of the x matrix should be computed once the data frame x is converted into a matrix.
\item[\code{bend}] a numeric value to add to the diagonal matrix in case matrix is singular for inversion.
\end{ldescription}
\end{Arguments}
%
\begin{Value}

\begin{ldescription}
\item[\code{K}]  pedigree transformed in a relationship matrix.
\item[\code{Kinv}]  inverse of the pedigree transformed in a relationship matrix.

\end{ldescription}
\end{Value}
%
\begin{Author}\relax
Giovanny Covarrubias-Pazaran
\end{Author}
%
\begin{References}\relax
Covarrubias-Pazaran G (2016) Genome assisted prediction of quantitative traits using the R package sommer. PLoS ONE 11(6): doi:10.1371/journal.pone.0156744 

\end{References}
%
\begin{Examples}
\begin{ExampleCode}

library(Matrix)
m <- matrix(1:9,3,3)
m <- tcrossprod(m)

mdf <- as.data.frame(as.table(m))
mdf

dfToMatrix(mdf, row = "Var1", column = "Var2", 
            value = "Freq",returnInverse=FALSE )


\end{ExampleCode}
\end{Examples}
\inputencoding{utf8}
\HeaderA{dsm}{Diagonal heterogeneous covariance structure}{dsm}
%
\begin{Description}\relax
Diagonal heterogeneous covariance structure. The returned dimensionless covariance shape is intended for use inside \code{vsm}, which supplies the single overall variance scale.
\end{Description}
%
\begin{Usage}
\begin{verbatim}
dsm(x, values = NULL, fixed = NULL, theta = NULL)
\end{verbatim}
\end{Usage}
%
\begin{Arguments}
\begin{ldescription}
\item[\code{x}] Variable or design defining the covariance levels.
\item[\code{values}] Optional positive starting variances, one per level.
\item[\code{fixed}] Logical vector of length q-1 indicating which variance-ratio parameters are fixed.
\item[\code{theta}] Optional q by q diagonal starting covariance matrix; if supplied, its diagonal replaces values.
\end{ldescription}
\end{Arguments}
%
\begin{Details}\relax
For \eqn{q}{} levels, \code{dsm} uses the first variance as the internal scale reference and estimates \eqn{q-1}{} positive variance ratios, \deqn{K=\mathrm{diag}(1,r_2,\ldots,r_q).}{} The ratios are optimized on log scales and reported on their natural positive scale.
\end{Details}
%
\begin{Value}
A list containing the incidence/design matrix in \code{Z} and a compiled CovarianceFactor v2 descriptor in \code{covFactor}, for use inside \code{vsm}.
\end{Value}
%
\begin{SeeAlso}\relax
\code{\LinkA{vsm}{vsm}}, \code{\LinkA{mmes}{mmes}}.
\end{SeeAlso}
%
\begin{Examples}
\begin{ExampleCode}
## Not run: 
vsm(dsm(environment), ism(genotype))

## End(Not run)
\end{ExampleCode}
\end{Examples}
\inputencoding{utf8}
\HeaderA{dsr}{Diagonal covariance structure for \code{mmer} and \code{vsr}}{dsr}
%
\begin{Description}\relax
Creates a diagonal variance-component structure for use by the \code{mmer} covariance-model interface, typically inside \code{\LinkA{vsr}{vsr}}. Each column or level represented by \code{x} receives its own diagonal variance component, while all off-diagonal covariance components are constrained to zero.
\end{Description}
%
\begin{Usage}
\begin{verbatim}
dsr(x)
\end{verbatim}
\end{Usage}
%
\begin{Arguments}
\begin{ldescription}
\item[\code{x}] 
A factor, character vector, numeric vector, or design/incidence matrix defining the covariance dimension. Factor and character inputs are expanded to incidence columns. A matrix is used directly.

\end{ldescription}
\end{Arguments}
%
\begin{Details}\relax
\code{dsr()} is part of the covariance-structure interface used by \code{mmer}/\code{vsr}; it is distinct from \code{\LinkA{dsm}{dsm}}, which belongs to the newer \code{mmes}/\code{vsm} CovarianceFactor interface.

The function returns a design matrix \eqn{Z}{} and a covariance-parameter constraint matrix \code{thetaC}. If \eqn{q}{} columns or levels are represented by \code{x}, then
\deqn{\mathrm{thetaC}=I_q.}{}
Consequently, the associated covariance matrix has the form
\deqn{\Sigma=\mathrm{diag}(\sigma_1^2,\ldots,\sigma_q^2),}{}
so that the \eqn{q}{} variances are estimated separately and all covariances are fixed to zero.

For a factor or character vector, missing observations are retained through the construction of the incidence matrix. If only one non-missing level is present, a one-column incidence matrix is constructed explicitly.

For a numeric non-factor input, \code{x} is treated as a single design column rather than expanded into factor levels.
\end{Details}
%
\begin{Value}
A list with components:
\begin{itemize}

\item{} \code{Z}: the design/incidence matrix associated with the covariance dimension.
\item{} \code{thetaC}: a diagonal matrix defining the diagonal covariance-component pattern used by \code{vsr}.

\end{itemize}

\end{Value}
%
\begin{SeeAlso}\relax
\code{\LinkA{vsr}{vsr}}, \code{\LinkA{usr}{usr}}, \code{\LinkA{csr}{csr}}, \code{\LinkA{atr}{atr}},
\code{\LinkA{dsm}{dsm}}, \code{\LinkA{mmer}{mmer}}
\end{SeeAlso}
%
\begin{Examples}
\begin{ExampleCode}
# Diagonal covariance structure across levels
D <- dsr(factor(c("A","B","A","C")))
D$Z
D$thetaC
\end{ExampleCode}
\end{Examples}
\inputencoding{utf8}
\HeaderA{dsumm}{Direct sum of residual covariance structures across sections}{dsumm}
%
\begin{Description}\relax
\code{dsumm} replicates a residual \code{\LinkA{vsm}{vsm}} structure independently for
every level (section) of a grouping variable, such as a trial, location, or
experiment. Each section receives its own variance and its own copy of every
covariance parameter of the inner structure.
\end{Description}
%
\begin{Usage}
\begin{verbatim}
dsumm(x, by, levels = NULL)
\end{verbatim}
\end{Usage}
%
\begin{Arguments}
\begin{ldescription}
\item[\code{x}] A residual \code{vsm()} term, for example
\code{vsm(ar1m(range), ar1m(row), ism(units))}. Its covariance-shaping
factors define the structure fitted within each section.
\item[\code{by}] Observation-level grouping variable defining the sections. Missing
values of \code{by} exclude the observation.
\item[\code{levels}] Optional character vector of \code{by} levels that receive the
structured covariance of \code{x}. All other levels receive an independent
and identically distributed residual with their own variance. The default
\code{NULL} applies the structure to every level.
\end{ldescription}
\end{Arguments}
%
\begin{Details}\relax
\strong{Direct sum versus Kronecker product.} A \code{vsm()} term is a Kronecker
product with a single overall variance. Including a diagonal section factor,
\deqn{\mathrm{vsm}(\mathrm{dsm}(t), \mathrm{ar1m}(r), \mathrm{ar1m}(c), \mathrm{ism}(units)):\quad
R = \sigma^2\, D_t \otimes K_r(\rho_r) \otimes K_c(\rho_c),}{}
gives each section its own variance, but a single correlation shared by all
sections, because each factor of a Kronecker product appears only once. In
contrast, \code{dsumm} builds the block-diagonal direct sum
\deqn{R = \bigoplus_{t} \sigma^2_t\, K_r(\rho_{r,t}) \otimes K_c(\rho_{c,t}),}{}
so every section has its own variance and its own correlations. Sections are
mutually independent, so their parameters are estimated only from their own
observations; with section-specific fixed effects the joint fit reproduces
separate per-section fits.

\strong{Choosing between the two.} Use the \code{vsm()} form when trials are
small or similar and a common spatial correlation is adequate (fewer
parameters). Use \code{dsumm} when correlations are expected to differ between
trials. With \eqn{T}{} sections and \eqn{p}{} inner covariance parameters,
\code{dsumm} estimates \eqn{T(p+1)}{} residual parameters.

\strong{Unlisted levels.} Levels omitted from \code{levels} are modelled as
\eqn{\sigma^2_t I}{}. Their observations are retained even when the variables of
the inner structure (for example row or range) are missing, so trials without
spatial coordinates can be analysed together with spatially designed trials.

\strong{Computation.} Each section is embedded in the rectangle spanned by its
observed levels of the inner factors. Covariance applications use the small
factor inverses, and missing cells (e.g. plots with a missing response) are
absorbed exactly by a Schur complement, so the dense observation-level
covariance matrix is never formed. This requires sections that are close to
complete rectangles and no \code{W} weights.

\strong{Reporting.} \code{\LinkA{covparams\_mmes}{covparams.Rul.mmes}} returns one row per parameter
and section, with the section in the \code{section} column. The variance rows
belong to the \code{by} factor; correlation rows are labelled by their inner
factor, e.g. \code{ar1m(range)}.
\end{Details}
%
\begin{Value}
A residual covariance structure, in the same form as a residual \code{vsm()}
term, for use in the \code{rcov} argument of \code{\LinkA{mmes}{mmes}}.
\end{Value}
%
\begin{Section}{Limitations}
\code{dsumm} is currently available only for residual (\code{rcov}) structures,
without \code{W} weights, and cannot be nested.
\end{Section}
%
\begin{SeeAlso}\relax
\code{\LinkA{vsm}{vsm}}, \code{\LinkA{ar1m}{ar1m}}, \code{\LinkA{dsm}{dsm}}, \code{\LinkA{mmes}{mmes}},
\code{\LinkA{covparams\_mmes}{covparams.Rul.mmes}}.
\end{SeeAlso}
%
\begin{Examples}
\begin{ExampleCode}
## Not run: 

mix <- mmes(yield ~ trial,
            rcov = ~ dsumm(vsm(ar1m(range), ar1m(row), ism(units)), by = trial),
            data = dat)
covparams_mmes(mix)

# Spatial structure only for two trials; remaining trials are iid.
mix2 <- mmes(yield ~ trial,
             rcov = ~ dsumm(vsm(ar1m(range), ar1m(row), ism(units)),
                            by = trial, levels = c("T1", "T2")),
             data = dat)

# Site-specific unstructured covariance among environments.
mix3 <- mmes(y ~ site,
             rcov = ~ dsumm(vsm(usm(env), ism(units)), by = site),
             data = dat2)

## End(Not run)
\end{ExampleCode}
\end{Examples}
\inputencoding{utf8}
\HeaderA{E.mat}{Epistatic relationship matrix}{E.mat}
%
\begin{Description}\relax
Calculates the realized epistatic relationship matrix of second order (additive x additive, additive x dominance, or dominance x dominance) using hadamard products with the C++ Armadillo library.
\end{Description}
%
\begin{Usage}
\begin{verbatim}
E.mat(X,nishio=TRUE,type="A#A",min.MAF=0.02)
\end{verbatim}
\end{Usage}
%
\begin{Arguments}
\begin{ldescription}
\item[\code{X}] 
Matrix (\eqn{n \times m}{}) of unphased genotypes for \eqn{n}{} lines and \eqn{m}{} biallelic markers, 
coded as \{-1,0,1\}. Fractional (imputed) and missing values (NA) are allowed.

\item[\code{nishio}] 
If TRUE, Nishio and Satoh (2014) (orthogonal parameterization of Vitezica et al. 2013), otherwise Su et al. (2012) (see Details in the D.mat help page).

\item[\code{type}] 
An argument specifying the type of epistatic relationship matrix desired. The default is the second order epistasis (additive x additive) type="A\#A". Other options are additive x dominant (type="A\#D"), or dominant by dominant (type="D\#D").  

\item[\code{min.MAF}] 
Minimum minor allele frequency. The A matrix is not sensitive to rare alleles, so by default only monomorphic markers are removed.

\end{ldescription}
\end{Arguments}
%
\begin{Details}\relax
it is computed as the Hadamard product of the epistatic relationship matrix; E=A\#A, E=A\#D, E=D\#D. 
\end{Details}
%
\begin{Value}
The epistatic relationship matrix is returned.

\end{Value}
%
\begin{References}\relax

Covarrubias-Pazaran G (2016) Genome assisted prediction of quantitative traits using the R package sommer. PLoS ONE 11(6): doi:10.1371/journal.pone.0156744 

Endelman, J.B., and J.-L. Jannink. 2012. Shrinkage estimation of the realized relationship matrix. G3:Genes, Genomes, Genetics. 2:1405-1413. doi: 10.1534/g3.112.004259

Nishio M and Satoh M. 2014. Including Dominance Effects in the Genomic BLUP Method for Genomic Evaluation. Plos One 9(1), doi:10.1371/journal.pone.0085792

Su G, Christensen OF, Ostersen T, Henryon M, Lund MS. 2012. Estimating Additive and Non-Additive Genetic Variances and Predicting Genetic Merits Using Genome-Wide Dense Single Nucleotide Polymorphism Markers. PLoS ONE 7(9): e45293. doi:10.1371/journal.pone.0045293

\end{References}
%
\begin{Examples}
\begin{ExampleCode}
####=========================================####
####random population of 200 lines with 1000 markers
####=========================================####
X <- matrix(rep(0,200*1000),200,1000)
for (i in 1:200) {
  X[i,] <- sample(c(-1,0,0,1), size=1000, replace=TRUE)
}

E <- E.mat(X, type="A#A") 
# if heterozygote markers are present can be used "A#D" or "D#D"
\end{ExampleCode}
\end{Examples}
\inputencoding{utf8}
\HeaderA{fam}{Factor-analytic covariance structure}{fam}
%
\begin{Description}\relax
Factor-analytic covariance structure. The returned dimensionless covariance shape is intended for use inside \code{vsm}, which supplies the single overall variance scale.
\end{Description}
%
\begin{Usage}
\begin{verbatim}
fam(x, k = 1L, loadings = NULL, specific = NULL, fixed = NULL,
    variance = c("heterogeneous", "homogeneous"))
\end{verbatim}
\end{Usage}
%
\begin{Arguments}
\begin{ldescription}
\item[\code{x}] Variable defining q covariance levels.
\item[\code{k}] Factor-analytic rank, satisfying 1 <= k < q.
\item[\code{loadings}] Optional finite q by k starting loading matrix.
\item[\code{specific}] Optional q positive starting specific variances. With
\code{variance="homogeneous"}, supply either one positive value or q equal
values. The default is 0.5 for every level.
\item[\code{fixed}] Logical vector controlling all loading and specific-variance-ratio parameters.
\item[\code{variance}] Whether to estimate separate specific variances for each level
(\code{"heterogeneous"}, the default) or a common specific variance
(\code{"homogeneous"}).
\end{ldescription}
\end{Arguments}
%
\begin{Details}\relax
The factor-analytic shape is based on \deqn{M=\Lambda\Lambda^{\mathsf T}+\Psi,}{} where \eqn{\Psi}{} is diagonal positive, followed by \eqn{K=M/M_{11}}{}.  For rotational identification, the leading \eqn{k\times k}{} loading block is lower triangular and its diagonal loadings are positive.  The first specific variance is an internal reference; remaining specific variances are positive ratios.  This removes the otherwise redundant internal scale because \code{vsm} supplies the overall \eqn{\sigma^2}{}.

With \code{variance="heterogeneous"}, \code{fixed} has length
\eqn{qk-k(k-1)/2+q-1}{}: triangular loadings followed by q-1 specific-variance
ratios. With \code{variance="homogeneous"}, all specific ratios are one,
\eqn{\Psi=I_q}{} internally, and \code{fixed} has length \eqn{qk-k(k-1)/2}{}.
The common specific variance on the fitted scale is estimated, not fixed
at one, and is reported as \code{common\_specific\_variance}. This homogeneous
model reproduces the covariance family previously fitted by \code{rrm()}.
For a genuine reduced-rank model with all specific variances fixed at zero,
use \code{\LinkA{rrm}{rrm}}.
\end{Details}
%
\begin{Value}
A list containing the incidence/design matrix in \code{Z} and a compiled CovarianceFactor v2 descriptor in \code{covFactor}, for use inside \code{vsm}.
\end{Value}
%
\begin{SeeAlso}\relax
\code{\LinkA{rrm}{rrm}}, \code{\LinkA{usm}{usm}}, \code{\LinkA{vsm}{vsm}}.
\end{SeeAlso}
%
\begin{Examples}
\begin{ExampleCode}
## Not run: 
vsm(fam(environment, k=2), ism(genotype))
vsm(fam(environment, k=2, variance="homogeneous"), ism(genotype))

## End(Not run)
\end{ExampleCode}
\end{Examples}
\inputencoding{utf8}
\HeaderA{familym}{Different distribution families for the traits of a multi-trait model}{familym}
\keyword{models}{familym}
%
\begin{Description}\relax
Assigns a \code{stats::family()} object to each trait of a long-format
multi-trait \code{\LinkA{mmes}{mmes}} model, e.g. a binary disease score together
with a Gaussian yield, so that their genetic and residual correlations can be
estimated jointly by penalized quasi-likelihood (PQL).
\end{Description}
%
\begin{Usage}
\begin{verbatim}
familym(by, ...)
\end{verbatim}
\end{Usage}
%
\begin{Arguments}
\begin{ldescription}
\item[\code{by}] Name (character string) of the column of \code{data} that
identifies the traits, typically \code{"trait"} as created by
\code{\LinkA{stackTraits}{stackTraits}}.
\item[\code{...}] Family objects named by the levels of the \code{by} column,
e.g. \code{Disease = binomial(), Yield = gaussian()}. Every level present in
the data needs a family.
\end{ldescription}
\end{Arguments}
%
\begin{Details}\relax
The object is passed as the \code{family} argument of \code{mmes}. Each PQL
iteration builds the working response and working weights of every row with
the family of its trait, and fits a weighted Gaussian multi-trait model.

The residual variance of binomial, Poisson and negative binomial traits is not
a free parameter (their dispersion is one), while that of Gaussian and other
families is estimated. The residual structure must therefore allow individual
trait variances to be fixed: use
\code{rcov = \textasciitilde{} vsm(corgm(trait, variance = "heterogeneous"), ism(record))}
(trait correlations and heterogeneous variances, see \code{\LinkA{corgm}{corgm}}) or
\code{rcov = \textasciitilde{} vsm(dsm(trait), ism(units))} (uncorrelated residuals).
The first level of the trait factor must be a fixed-dispersion trait,
because its variance is the \code{vsm} scale; reorder the levels
otherwise. \code{mmes} then fixes the residual variance of every
fixed-dispersion trait to one and estimates the rest.

As for single-family PQL fits, the likelihood, \code{AIC} and \code{BIC} are
\code{NA}. The fitted object has a named \code{dispersion} vector (one per
trait: one for fixed-dispersion traits, the REML residual variance for the
others), and \code{fitted()} and \code{residuals()} apply the inverse link and
variance function of each trait to its rows. Random-effect variances are on
the link scale of each trait.
\end{Details}
%
\begin{Value}
An object of class \code{"sommer\_familym"}.
\end{Value}
%
\begin{SeeAlso}\relax
\code{\LinkA{mmes}{mmes}}, \code{\LinkA{corgm}{corgm}}, \code{\LinkA{stackTraits}{stackTraits}},
\code{\LinkA{covmatrix\_mmes}{covmatrix.Rul.mmes}}
\end{SeeAlso}
%
\begin{Examples}
\begin{ExampleCode}
set.seed(12)
ng <- 100; nr <- 4
G <- matrix(c(1, 0.4, 0.4, 0.8), 2)
u <- matrix(rnorm(ng * 2), ng) %*% chol(G)
d <- data.frame(id = factor(rep(paste0("g", 1:ng), each = nr)))
d$disease <- rbinom(nrow(d), 1, plogis(-0.5 + u[as.integer(d$id), 1]))
d$yield <- 10 + u[as.integer(d$id), 2] + rnorm(nrow(d))
L <- stackTraits(d, traits = c("disease", "yield"))

fit <- mmes(value ~ trait,
  random = ~ vsm(usm(trait), ism(id)),
  rcov = ~ vsm(corgm(trait, variance = "heterogeneous"), ism(record)),
  family = familym("trait", disease = binomial(), yield = gaussian()),
  data = L,
  control = mmes_control(diagnostics = list(verbose = FALSE)))
fit$dispersion
covmatrix_mmes(fit, 1)$correlation   # genetic correlation (link scale)
head(fitted(fit))                    # probabilities for disease rows
\end{ExampleCode}
\end{Examples}
\inputencoding{utf8}
\HeaderA{fitted.mmes}{fitted form a LMM fitted with mmes}{fitted.mmes}
\aliasA{fitted.mmer}{fitted.mmes}{fitted.mmer}
\keyword{models}{fitted.mmes}
%
\begin{Description}\relax
\code{fitted} method for class \code{"mmes"}.
\end{Description}
%
\begin{Usage}
\begin{verbatim}
## S3 method for class 'mmes'
fitted(object, type=c("response", "link"), ...)
\end{verbatim}
\end{Usage}
%
\begin{Arguments}
\begin{ldescription}
\item[\code{object}] an object of class \code{"mmes"}
\item[\code{type}] For a non-Gaussian PQL fit, return fitted means on the response
scale (default) or linear predictors on the link scale.
\item[\code{...}] Further arguments to be passed to the mmes function
\end{ldescription}
\end{Arguments}
%
\begin{Value}
For Gaussian fits, fitted values of the form y.hat = Xb + Zu. For non-Gaussian
PQL fits, response-scale conditional means by default, or Xb + Zu with
\code{type="link"}.

\end{Value}
%
\begin{Author}\relax
Giovanny Covarrubias
\end{Author}
%
\begin{SeeAlso}\relax
\code{\LinkA{fitted}{fitted}}, \code{\LinkA{mmes}{mmes}}
\end{SeeAlso}
%
\begin{Examples}
\begin{ExampleCode}
# data(DT_cpdata, package="enhancer")
# DT <- DT_cpdata
# GT <- GT_cpdata
# MP <- MP_cpdata
# #### create the variance-covariance matrix
# A <- A.mat(GT) # additive relationship matrix
# #### look at the data and fit the model
# head(DT)
# mix1 <- mmes(Yield~1,
#               random=~vsm(ism(id),Gu=A)
#                       + Rowf + Colf + spl2Dc(Row,Col),
#               rcov=~units,
#               data=DT)
# 
# ff=fitted(mix1)
# 
# colfunc <- colorRampPalette(c("steelblue4","springgreen","yellow"))
# lattice::wireframe(`u:Row.fitted`~Row*Col, data=ff$dataWithFitted,  
#           aspect=c(61/87,0.4), drape=TRUE,# col.regions = colfunc,
#           light.source=c(10,0,10))
# lattice::levelplot(`u:Row.fitted`~Row*Col, data=ff$dataWithFitted, col.regions = colfunc)

\end{ExampleCode}
\end{Examples}
\inputencoding{utf8}
\HeaderA{fixm}{fixed indication matrix}{fixm}
%
\begin{Description}\relax
\code{fixm} creates a square matrix with 3's in the diagnals and off-diagonals to quickly specify a fixed constraint in the Gtc argument of the \code{\LinkA{vsm}{vsm}} function.
\end{Description}
%
\begin{Usage}
\begin{verbatim}
  fixm(x, reps=NULL)
\end{verbatim}
\end{Usage}
%
\begin{Arguments}
\begin{ldescription}
\item[\code{x}] integer specifying the number of traits to be fitted for a given random effect.
\item[\code{reps}] integer specifying the number of times the matrix should be repeated in a list format to provide easily the constraints in complex models that use the ds(), us() or cs() structures.
\end{ldescription}
\end{Arguments}
%
\begin{Value}
\begin{description}


\item[\$res]  a matrix or a list of matrices with the constraints to be provided in the Gtc argument of the \code{\LinkA{vsm}{vsm}} function.

\end{description}

\end{Value}
%
\begin{Author}\relax
Giovanny Covarrubias-Pazaran
\end{Author}
%
\begin{References}\relax
Covarrubias-Pazaran G (2016) Genome assisted prediction of quantitative traits using the R package sommer. PLoS ONE 11(6): doi:10.1371/journal.pone.0156744 
\end{References}
%
\begin{Examples}
\begin{ExampleCode}
fixm(4)
fixm(4,2)
\end{ExampleCode}
\end{Examples}
\inputencoding{utf8}
\HeaderA{generate\_mme\_workload}{Generate Parameterized MME Workloads for Solver Evaluation}{generate.Rul.mme.Rul.workload}
%
\begin{Description}\relax
Instantiates reproducible synthetic mixed model datasets and formulas spanning
factor dimensions, covariance structures, relationship densities, and tasks.
\end{Description}
%
\begin{Usage}
\begin{verbatim}
generate_mme_workload(name = NULL,
                      n_ind = 60L,
                      n_levels = 5L,
                      structure = c("ar1", "ar2", "us", "diag", "fa",
                                    "rr", "spatial", "sar", "car"),
                      relationship = c("none", "dense", "sparse"),
                      missing_rate = 0.0,
                      reml = TRUE,
                      computeCi = 0L,
                      seed = 42L)
\end{verbatim}
\end{Usage}
%
\begin{Arguments}
\begin{ldescription}
\item[\code{name}] Descriptive label for the workload.
\item[\code{n\_ind}] Number of unique individuals / genotypes. Default is 60.
\item[\code{n\_levels}] Number of factor levels (e.g. environments, timepoints, or sites). Default is 5.
\item[\code{structure}] Covariance structure type. Supported: \code{"ar1"}, \code{"ar2"},
\code{"us"}, \code{"diag"}, \code{"fa"}, \code{"rr"}, \code{"spatial"}, \code{"sar"}, \code{"car"}.
\item[\code{relationship}] Individual relationship matrix: \code{"none"} (identity),
\code{"dense"} (genomic relationship inverse), or \code{"sparse"} (pedigree inverse).
\item[\code{missing\_rate}] Fraction of missing response values (0.0 to 0.5). Default is 0.0.
\item[\code{reml}] Logical; whether to fit by REML (TRUE) or ML (FALSE). Default is TRUE.
\item[\code{computeCi}] Final inverse demand: 0 (none), 1 (Takahashi selected inverse), 2 (full inverse).
\item[\code{seed}] Random seed for reproducible dataset simulation.
\end{ldescription}
\end{Arguments}
%
\begin{Value}
A list of class \code{"mme\_workload"} containing data, formulas, REML and
computeCi settings, and workload metadata.
\end{Value}
%
\begin{SeeAlso}\relax
\code{\LinkA{mmes}{mmes}}, \code{\LinkA{mmes\_control}{mmes.Rul.control}}
\end{SeeAlso}
%
\begin{Examples}
\begin{ExampleCode}
# Generate a small 1D autoregressive workload
wl <- generate_mme_workload(
  name = "demo_ar1",
  n_ind = 20,
  n_levels = 4,
  structure = "ar1"
)
\end{ExampleCode}
\end{Examples}
\inputencoding{utf8}
\HeaderA{GWAS}{Genome wide association study analysis}{GWAS}
%
\begin{Description}\relax

THIS FUNCTION IS DEPRECATED. Fits a multivariate/univariate linear mixed model GWAS by likelihood methods (REML), see the Details section below. It uses the \code{\LinkA{mmer}{mmer}} function and its core coded in C++ using the Armadillo library to optimize dense matrix operations common in the derect-inversion algorithms. After the model fit extracts the inverse of the phenotypic variance matrix to perform the association test for the "p" markers. Please check the Details section (Model enabled) if you have any issue with making the function run.

The sommer package is updated on CRAN every 3-months due to CRAN policies but you can find the latest source at https://github.com/covaruber/sommer . This can be easily installed typing the following in the R console:

library(devtools)

install\_github("covaruber/sommer")

This is recommended since bugs fixes will be immediately available in the GitHub source. \strong{For tutorials} on how to perform different analysis with sommer please look at the vignettes by typing in the terminal:

\strong{vignette("v1.sommer.quick.start")} 

\strong{vignette("v2.sommer.changes.and.faqs")}

\strong{vignette("v3.sommer.qg")}

\strong{vignette("v4.sommer.gxe")}

or visit \strong{https://covaruber.github.io}

\end{Description}
%
\begin{Usage}
\begin{verbatim}

GWAS(fixed, random, rcov, data, weights, W,
    nIters=20, tolParConvLL = 1e-03, tolParInv = 1e-06, 
    init=NULL, constraints=NULL,method="NR", 
    getPEV=TRUE,naMethodX="exclude",
    naMethodY="exclude",returnParam=FALSE, 
    dateWarning=TRUE,date.warning=TRUE,verbose=FALSE,
    stepWeight=NULL, emWeight=NULL,
    M=NULL, gTerm=NULL, n.PC = 0, min.MAF = 0.05, 
    P3D = TRUE)
\end{verbatim}
\end{Usage}
%
\begin{Arguments}
\begin{ldescription}
\item[\code{fixed}] A formula specifying the \strong{response variable(s)} \strong{and fixed effects}, i.e:

\emph{response \textasciitilde{} covariate} for univariate models

\emph{cbind(response.i,response.j) \textasciitilde{} covariate} for multivariate models

The fcm function can be used to constrain fixed effects in multi-response models.



\item[\code{random}] a formula specifying the name of the \strong{random effects}, i.e. \emph{random= \textasciitilde{} genotype + year}.

Useful functions can be used to fit heterogeneous variances and other special models (\emph{see 'Special Functions' in the Details section for more information}):

\code{\LinkA{vsr}{vsr}(...,Gu,Gt,Gtc)} is the main function to specify variance models and special structures for random effects. On the ... argument you provide the unknown variance-covariance structures (i.e. usr,dsr,at,csr) and the random effect where such covariance structure will be used (the random effect of interest). Gu is used to provide known covariance matrices among the levels of the random effect, Gt initial values and Gtc for constraints. Auxiliar functions for building the variance models are:


\item[\code{rcov}] a formula specifying the name of the \strong{error term}, i.e. \emph{rcov= \textasciitilde{} units}.

The functions that can be used to fit heterogeneous residual variances are the same used on the random term but the random effect is always "units", i.e. \emph{rcov=\textasciitilde{}vsr(dsr(Location),units)}


\item[\code{data}] a data frame containing the variables specified in the formulas for response, fixed, and random effects.
\item[\code{weights}] name of the covariate for weights. To be used for the product R = Wsi*R*Wsi, where * is the matrix product, Wsi is the square root of the inverse of W and R is the residual matrix.
\item[\code{W}] Alternatively, instead of providing a vector of weights the user can specify an entire W matrix (e.g., when covariances exist). To be used first to produce Wis = solve(chol(W)), and then calculate R = Wsi*R*Wsi.t(), where * is the matrix product, and R is the residual matrix. Only one of the arguments weights or W should be used. If both are indicated W will be given the preference.
\item[\code{nIters}] Maximum number of iterations allowed.  Default value is 15.
\item[\code{tolParConvLL}] Convergence criteria.
\item[\code{tolParInv}] tolerance parameter for matrix inverse used when singularities are encountered.
\item[\code{init}] initial values for the variance components. By default this is NULL and variance components are estimated by the method selected, but in case the user want to provide initial values for ALL var-cov components this argument is functional. It has to be provided as a list or an array, where each list element is one variance component and if multitrait model is pursued each element of the list is a matrix of variance covariance components among traits. Initial values can also be provided in the Gt argument of the \LinkA{vsr}{vsr} function.Is highly encouraged to use the Gt and Gtc arguments of the \LinkA{vsr}{vsr} function instead of this argument
\item[\code{constraints}] when initial values are provided these have to be accompanied by their constraints. See the \LinkA{vsr}{vsr} function for more details on the constraints. Is highly encouraged to use the Gt and Gtc arguments of the \LinkA{vsr}{vsr} function instead of this argument.
\item[\code{method}] this refers to the method or algorithm to be used for estimating variance components. Direct-inversion Newton-Raphson \strong{NR} and Average Information \strong{AI} (Tunnicliffe 1989; Gilmour et al. 1995; Lee et al. 2015).
\item[\code{getPEV}] a TRUE/FALSE value indicating if the program should return the predicted error variance and variance for random effects. This option is provided since this can take a long time for certain models where p > n by a big extent.
\item[\code{naMethodX}] one of the two possible values; "include" or "exclude". If "include" is selected then the function will impute the X matrices for fixed effects with the median value. If "exclude" is selected it will get rid of all rows with missing values for the X (fixed) covariates. The default is "exclude". The "include" option should be used carefully.
\item[\code{naMethodY}] one of the three possible values; "include", "include2" or "exclude". If "include" is selected then the function will impute the response variables with the median value. The difference between "include" and "include2" is only available in the multitrait models when the imputation can happen for the entire matrix of responses or only for complete cases ("include2"). If "exclude" is selected it will get rid of rows in responses where missing values are present for the estimation of variance components. The default is "exclude".
\item[\code{returnParam}] a TRUE/FALSE value to indicate if the program should return the parameters used for modeling without fitting the model.
\item[\code{dateWarning}] a TRUE/FALSE value to indicate if the program should warn you when is time to update the sommer package.
\item[\code{date.warning}] a TRUE/FALSE value to indicate if the program should warn you when is time to update the sommer package.
\item[\code{verbose}] a TRUE/FALSE value to indicate if the program should return the progress of the iterative algorithm.
\item[\code{stepWeight}] A vector of values (of length equal to the number of iterations) indicating the weight used to multiply the update (delta) for variance components at each iteration. If NULL the 1st iteration will be multiplied by 0.5, the 2nd by 0.7, and the rest by 0.9. This argument can help to avoid that variance components go outside the parameter space in the initial iterations which doesn't happen very often with the NR method but it can be detected by looking at the behavior of the likelihood. In that case you may want to give a smaller weight to the initial 8-10 iterations.
\item[\code{emWeight}] A vector of values (of length equal to the number of iterations) indicating with values between 0 and 1 the weight assigned to the EM information matrix. And the values 1 - emWeight will be applied to the NR/AI information matrix to produce a joint information matrix. If NULL weights for EM information matrix are zero and 1 for the NR/AI information matrix.
\item[\code{M}] The marker matrix containing the marker scores for each level of the random effect selected in the gTerm argument, coded as numeric based on the number of reference alleles in the genotype call, e.g. (-1,0,1) = (aa,Aa,AA), levels in diploid individuals. Individuals in rows and markers in columns. No additional columns should be provided, is a purely numerical matrix. Similar logic applies to polyploid individuals, e.g. (-3,-2,-1,0,1,2,3) = (aaaa,aaaA,aaAA,Aaaa,AAaa,AAAa,AAAA). 
\item[\code{gTerm}] a character vector indicating the random effect linked to the marker matrix M (i.e. the genetic term) in the model. The random effect selected should have the same number of levels than the number of rows of M. When fitting only a random effect without a special covariance structure (e.g., dsr, usr, etc.) you will need to add the call 'u:' to the name of the random effect given the behavior of the naming rules of the solver when having a simple random effect without covariance structure. 
\item[\code{n.PC}] Number of principal components to include as fixed effects. Default is 0 (equals K model).
\item[\code{min.MAF}] Specifies the minimum minor allele frequency (MAF). If a marker has a MAF less than min.MAF, it is assigned a zero score.
\item[\code{P3D}] When P3D=TRUE, variance components are estimated by REML only once, without any markers in the model and then a for loop for hypothesis testing is performed. When P3D=FALSE, variance components are estimated by REML for each marker separately. The latter can be quite time consuming. As many models will be run as number of marker.

\end{ldescription}
\end{Arguments}
%
\begin{Details}\relax

\strong{Citation}

Type \emph{citation("sommer")} to know how to cite the sommer package in your publications. 

\strong{Models Enabled}

For details about the models enabled and more information about the covariance structures please check the help page of the package (\code{\LinkA{sommer}{sommer}}). In general the GWAS model implemented in sommer to obtain marker effect is a generalized linear model of the form:

b = (X'V-X)X'V-y

with X = ZMi

where:
b is the marker effect (dimensions 1 x mt)
y is the response variable (univariate or multivariate) (dimensions 1 x nt)
V- is the inverse of the phenotypic variance matrix (dimensions nt x nt)
Z is the incidence matrix for the random effect selected (gTerm argument) to perform the GWAS (dimensions nt x ut)
Mi is the ith column of the marker matrix (M argument) (dimensions u x m)

for t traits, n observations, m markers and u levels of the random effect. Depending if P3D is TRUE or FALSE the V- matrix will be calculated once and used for all marker tests (P3D=TRUE) or estimated through REML for each marker (P3D=FALSE).


\emph{vignette('sommer.start')}

\strong{Bug report and contact}

If you have any technical questions or suggestions please post it in https://stackoverflow.com or https://stats.stackexchange.com. 

If you have any bug report please go to https://github.com/covaruber/sommer or send me an email to address it asap.

\end{Details}
%
\begin{Value}
If all parameters are correctly indicated the program will return a list with the following information:

\begin{ldescription}
\item[\code{Vi}] the inverse of the phenotypic variance matrix V\textasciicircum{}- = (ZGZ+R)\textasciicircum{}-1
\item[\code{sigma}] a list with the values of the variance-covariance components with one list element for each random effect.
\item[\code{sigma\_scaled}] a list with the values of the scaled variance-covariance components with one list element for each random effect.
\item[\code{sigmaSE}] Hessian matrix containing the variance-covariance for the variance components. SE's can be obtained taking the square root of the diagonal values of the Hessian.
\item[\code{Beta}] a data frame for trait BLUEs (fixed effects).
\item[\code{VarBeta}] a variance-covariance matrix for trait BLUEs
\item[\code{U}] a list (one element for each random effect) with a data frame for trait BLUPs.
\item[\code{VarU}] a list (one element for each random effect) with the variance-covariance matrix for trait BLUPs.
\item[\code{PevU}] a list (one element for each random effect) with the predicted error variance matrix for trait BLUPs.
\item[\code{fitted}]  Fitted values y.hat=XB
\item[\code{residuals}] Residual values e = Y - XB
\item[\code{AIC}] Akaike information criterion
\item[\code{BIC}] Bayesian information criterion
\item[\code{convergence}] a TRUE/FALSE statement indicating if the model converged.
\item[\code{monitor}] The values of log-likelihood and variance-covariance components across iterations during the REML estimation.
\item[\code{scores}]  marker scores (-log\_(10)p) for the traits
\item[\code{method}] The method for extimation of variance components specified by the user.
\item[\code{constraints}] contraints used in the mixed models for the random effects.

\end{ldescription}
\end{Value}
%
\begin{Author}\relax
Giovanny Covarrubias-Pazaran
\end{Author}
%
\begin{References}\relax

Covarrubias-Pazaran G. Genome assisted prediction of quantitative traits using the R package sommer. PLoS ONE 2016, 11(6): doi:10.1371/journal.pone.0156744   

Covarrubias-Pazaran G. 2018. Software update: Moving the R package sommer to multivariate mixed models for genome-assisted prediction. doi: https://doi.org/10.1101/354639

Bernardo Rex. 2010. Breeding for quantitative traits in plants. Second edition. Stemma Press. 390 pp.

Gilmour et al. 1995. Average Information REML: An efficient algorithm for variance parameter estimation in linear mixed models. Biometrics 51(4):1440-1450.

Kang et al. 2008. Efficient control of population structure in model organism association mapping. Genetics 178:1709-1723.

Lee, D.-J., Durban, M., and Eilers, P.H.C. (2013). Efficient two-dimensional smoothing with P-spline ANOVA mixed models and nested bases. Computational Statistics and Data Analysis, 61, 22 - 37.

Lee et al. 2015. MTG2: An efficient algorithm for multivariate linear mixed model analysis based on genomic information. Cold Spring Harbor. doi: http://dx.doi.org/10.1101/027201.

Maier et al. 2015. Joint analysis of psychiatric disorders increases accuracy of risk prediction for schizophrenia, bipolar disorder, and major depressive disorder. Am J Hum Genet; 96(2):283-294.

Rodriguez-Alvarez, Maria Xose, et al. Correcting for spatial heterogeneity in plant breeding experiments with P-splines. Spatial Statistics 23 (2018): 52-71.

Searle. 1993. Applying the EM algorithm to calculating ML and REML estimates of variance components. Paper invited for the 1993 American Statistical Association Meeting, San Francisco.

Yu et al. 2006. A unified mixed-model method for association mapping that accounts for multiple levels of relatedness. Genetics 38:203-208.

Tunnicliffe W. 1989. On the use of marginal likelihood in time series model estimation. JRSS 51(1):15-27.

Zhang et al. 2010. Mixed linear model approach adapted for genome-wide association studies. Nat. Genet. 42:355-360.

\end{References}
%
\begin{Examples}
\begin{ExampleCode}

####=========================================####
#### For CRAN time limitations most lines in the 
#### examples are silenced with one '#' mark, 
#### remove them and run the examples using
#### command + shift + C |OR| control + shift + C
####=========================================####
#####========================================####
##### potato example
#####========================================####
# 
# data(DT_polyploid, package="enhancer")
# DT <- DT_polyploid
# GT <- GT_polyploid
# MP <- MP_polyploid
# ####=========================================####
# ####### convert markers to numeric format
# ####=========================================####
# numo <- atcg1234(data=GT, ploidy=4);
# numo$M[1:5,1:5];
# numo$ref.allele[,1:5]
# 
# ###=========================================####
# ###### plants with both genotypes and phenotypes
# ###=========================================####
# common <- intersect(DT$Name,rownames(numo$M))
# 
# ###=========================================####
# ### get the markers and phenotypes for such inds
# ###=========================================####
# marks <- numo$M[common,]; marks[1:5,1:5]
# DT2 <- DT[match(common,DT$Name),];
# DT2 <- as.data.frame(DT2)
# DT2[1:5,]
# 
# ###=========================================####
# ###### Additive relationship matrix, specify ploidy
# ###=========================================####
# A <- A.mat(marks)
# ###=========================================####
# ### run it as GWAS model
# ###=========================================####
# ans2 <- GWAS(tuber_shape~1,
#              random=~vsr(Name,Gu=A),
#              rcov=~units,
#              gTerm = "u:Name",
#              M=marks, data=DT2)
# plot(ans2$scores[,1])
\end{ExampleCode}
\end{Examples}
\inputencoding{utf8}
\HeaderA{H.mat}{Combined relationship matrix H}{H.mat}
%
\begin{Description}\relax
Given a matrix A and a matrix G returns a H matrix with the C++ Armadillo library.
\end{Description}
%
\begin{Usage}
\begin{verbatim}
H.mat(A, G, tau = 1, omega = 1, tolparinv=1e-6)
\end{verbatim}
\end{Usage}
%
\begin{Arguments}
\begin{ldescription}
\item[\code{A}] 
Additive relationship matrix based on pedigree.

\item[\code{G}] 
Additive relationship matrix based on marker data.  

\item[\code{tau}] 
As described by Martini et al. (2018).

\item[\code{omega}] 
As described by Martini et al. (2018). 

\item[\code{tolparinv}] 
Tolerance parameter for matrix inverse used when singularities are encountered in the estimation procedure.

\end{ldescription}
\end{Arguments}
%
\begin{Details}\relax
See references 
\end{Details}
%
\begin{Value}

H Matrix with the relationship between the individuals based on pedigree and corrected by molecular information

\end{Value}
%
\begin{References}\relax

Covarrubias-Pazaran G (2016) Genome assisted prediction of quantitative traits using the R package sommer. PLoS ONE 11(6): doi:10.1371/journal.pone.0156744 

Martini, J. W., Schrauf, M. F., Garcia-Baccino, C. A., Pimentel, E. C., Munilla, S., Rogberg-Munoz, A., ... \& Simianer, H. (2018). The effect of the H-1 scaling factors tau and omega on the structure of H in the single-step procedure. Genetics Selection Evolution, 50(1), 16.

\end{References}
%
\begin{Examples}
\begin{ExampleCode}
####=========================================####
####random population of 200 lines with 1000 markers
####=========================================####
M <- matrix(rep(0,200*1000),200,1000)
for (i in 1:200) {
  M[i,] <- sample(c(-1,0,0,1), size=1000, replace=TRUE)
}
rownames(M) <- 1:nrow(M)
v <- sample(1:nrow(M),100)
M2 <- M[v,]

A <- A.mat(M) # assume this is a pedigree-based matrix for the sake of example
G <- A.mat(M2)

H <- H.mat(A,G)
# colfunc <- colorRampPalette(c("steelblue4","springgreen","yellow"))
# hv <- heatmap(H[1:15,1:15], col = colfunc(100),Colv = "Rowv")
\end{ExampleCode}
\end{Examples}
\inputencoding{utf8}
\HeaderA{ism}{Identity covariance structure}{ism}
%
\begin{Description}\relax
Identity covariance structure. The returned dimensionless covariance shape is intended for use inside \code{vsm}, which supplies the single overall variance scale.
\end{Description}
%
\begin{Usage}
\begin{verbatim}
ism(x)
\end{verbatim}
\end{Usage}
%
\begin{Arguments}
\begin{ldescription}
\item[\code{x}] Factor, character vector, numeric design, or matrix defining covariance levels.
\end{ldescription}
\end{Arguments}
%
\begin{Details}\relax
\code{ism} returns \eqn{K=I}{}.  It has no covariance-shape parameters; the only unknown covariance scale is the \eqn{\sigma^2}{} supplied by \code{vsm}.  It is commonly used as the final term of \code{vsm} to indicate an independent main effect.
\end{Details}
%
\begin{Value}
A list containing the incidence/design matrix in \code{Z} and a compiled CovarianceFactor v2 descriptor in \code{covFactor}, for use inside \code{vsm}.
\end{Value}
%
\begin{SeeAlso}\relax
\code{\LinkA{vsm}{vsm}}, \code{\LinkA{mmes}{mmes}}.
\end{SeeAlso}
%
\begin{Examples}
\begin{ExampleCode}
## Not run: 
vsm(ism(id))
vsm(dsm(environment), ism(genotype))

## End(Not run)
\end{ExampleCode}
\end{Examples}
\inputencoding{utf8}
\HeaderA{loadings\_mmes}{Extract factor-analytic loadings from a fitted mmes model}{loadings.Rul.mmes}
%
\begin{Description}\relax
Reconstructs the normalized loadings matrix and specific variances of a
\code{fam()} or \code{rrm()} covariance-shaping factor fitted inside
\code{vsm()}, on the same scale as the fitted covariance matrix returned in
\code{object\$theta}.
\end{Description}
%
\begin{Usage}
\begin{verbatim}
loadings_mmes(object, term = NULL, varianceScale=TRUE, rotation=TRUE)
\end{verbatim}
\end{Usage}
%
\begin{Arguments}
\begin{ldescription}

\item[\code{object}] a fitted model of class \code{"mmes"}.
\item[\code{term}] character name of the random term (matching a name in
\code{object\$covStruct}) built with a single \code{fam()} or \code{rrm()}
covariance-shaping factor. If \code{NULL}, the unique such term in the model
is used automatically; an error is raised if none or more than one exist.
\item[\code{varianceScale}] Whether to return loadings multiplied by
\code{sqrt(sigma2)} and specific variances multiplied by \code{sigma2}.
When \code{FALSE}, both describe the dimensionless covariance shape.
\item[\code{rotation}] a logical value to indicate if loadings should be rotated by its
singular vectors.

\end{ldescription}
\end{Arguments}
%
\begin{Details}\relax
\code{fam()}/\code{rrm()} remove their otherwise redundant internal scale so
that \code{vsm()} owns the single overall variance; the loadings and specific
variances reported by \code{loadings\_mmes} are re-normalized so that
\deqn{\Sigma = \sigma^2(\Lambda\Lambda^{\mathsf T}+\Psi)}{}
reproduces \code{object\$theta[[term]]} exactly when
\code{varianceScale=FALSE}, where \eqn{\sigma^2}{} is
\code{object\$covPar[[term]][1]}. With \code{varianceScale=TRUE}, the
returned loadings crossproduct plus the diagonal of the returned specific
variances directly reproduces that covariance.
\end{Details}
%
\begin{Value}
A list containing:
\begin{description}

\item[loadings] a levels-by-factor matrix \eqn{\Lambda}{}.
\item[specific] a named vector of specific variances \eqn{\Psi}{} (identically
zero for every level in a \code{rrm()} term).
\item[sigma2] the fitted overall variance scale.
\item[model] \code{"fa"} or \code{"rr"}.
\item[term] the resolved term name.

\end{description}

\end{Value}
%
\begin{SeeAlso}\relax
\code{\LinkA{scores\_mmes}{scores.Rul.mmes}}, \code{\LinkA{fam}{fam}}, \code{\LinkA{rrm}{rrm}}, \code{\LinkA{vsm}{vsm}}, \code{\LinkA{mmes}{mmes}}.
\end{SeeAlso}
%
\begin{Examples}
\begin{ExampleCode}
## Not run: 
mix <- mmes(BLUEs ~ trial,
            random = ~ vsm(fam(trial, 2), ism(genotype)),
            rcov = ~ units, data = dt)
loadings_mmes(mix)

## End(Not run)
\end{ExampleCode}
\end{Examples}
\inputencoding{utf8}
\HeaderA{mam}{Moving-average covariance structures of order one or two}{mam}
\aliasA{ma1m}{mam}{ma1m}
\aliasA{ma2m}{mam}{ma2m}
%
\begin{Description}\relax
Moving-average covariance structures of order one or two. The returned dimensionless covariance shape is intended for use inside \code{vsm}, which supplies the single overall variance scale.
\end{Description}
%
\begin{Usage}
\begin{verbatim}
mam(x, order = 1L, theta = NULL, fixed = NULL)
ma1m(x, theta = 0.15, fixed = FALSE)
ma2m(x, theta = c(0.15, 0.05), fixed = NULL)
\end{verbatim}
\end{Usage}
%
\begin{Arguments}
\begin{ldescription}
\item[\code{x}] Ordered factor defining covariance levels.
\item[\code{order}] Moving-average order; currently 1 or 2.
\item[\code{theta}] Starting MA polynomial coefficients.  Defaults to 0.15 for each order.
\item[\code{fixed}] Logical vector with one value per MA coefficient.
\end{ldescription}
\end{Arguments}
%
\begin{Details}\relax
Using the conventional polynomial \eqn{e_t+\theta_1e_{t-1}+\theta_2e_{t-2}}{}, the autocovariance at lag \eqn{h}{} is proportional to \eqn{\sum_{j=0}^{p-h}\theta_j\theta_{j+h}}{} with \eqn{\theta_0=1}{}.  The covariance is normalized to correlation scale and is zero beyond lag \eqn{p}{}.  Invertibility of the MA polynomial is not required merely to define this covariance matrix.  \code{ma1m} and \code{ma2m} are convenience wrappers.
\end{Details}
%
\begin{Value}
A list containing the incidence/design matrix in \code{Z} and a compiled CovarianceFactor v2 descriptor in \code{covFactor}, for use inside \code{vsm}.
\end{Value}
%
\begin{SeeAlso}\relax
\code{\LinkA{ar1m}{ar1m}}, \code{\LinkA{toeplitzm}{toeplitzm}}, \code{\LinkA{vsm}{vsm}}.
\end{SeeAlso}
%
\begin{Examples}
\begin{ExampleCode}
## Not run: 
vsm(ma1m(time), ism(id))
vsm(ma2m(time), ism(id))

## End(Not run)
\end{ExampleCode}
\end{Examples}
\inputencoding{utf8}
\HeaderA{maternm}{Matern spatial covariance structure}{maternm}
%
\begin{Description}\relax
Matern spatial covariance structure. The returned dimensionless covariance shape is intended for use inside \code{vsm}, which supplies the single overall variance scale.
\end{Description}
%
\begin{Usage}
\begin{verbatim}
maternm(x, range = NULL, nu = 0.5, fixed = c(FALSE, FALSE),
        distance = NULL, anisotropy = c("none", "geometric"),
        angle = pi/2, ratio = 1.5)
\end{verbatim}
\end{Usage}
%
\begin{Arguments}
\begin{ldescription}
\item[\code{x}] Numeric coordinate vector or matrix/data frame whose rows give spatial coordinates.
\item[\code{range}] Positive starting range parameter.  If NULL, the median positive pairwise distance is used.
\item[\code{nu}] Positive starting Matern smoothness parameter.
\item[\code{fixed}] Logical of length one or two (range, nu), or four with geometric anisotropy (range, nu, angle, ratio); length one is recycled.
\item[\code{distance}] Optional finite symmetric q by q distance matrix with zero diagonal, aligned to the unique spatial locations.
\item[\code{anisotropy}] \code{"geometric"} rotates two-dimensional coordinates by \code{angle} and divides the second rotated axis by \code{ratio} before computing distances (ASReml anisotropic Matern).
\item[\code{angle}] Starting anisotropy angle in (0, pi).
\item[\code{ratio}] Starting positive anisotropy ratio.
\end{ldescription}
\end{Arguments}
%
\begin{Details}\relax
For distance \eqn{d}{}, the correlation is \deqn{K(d)=\frac{2^{1-\nu}}{\Gamma(\nu)}z^{\nu}K_{\nu}(z),\qquad z=\frac{\sqrt{2\nu}d}{\phi},}{} where \eqn{\phi}{} is \code{range}.  Both range and smoothness are positive and are optimized on log scales.  The implementation uses the generic R evaluator and central factor-level numerical derivatives.  Repeated coordinate rows are mapped to the same covariance level.
\end{Details}
%
\begin{Value}
A list containing the incidence/design matrix in \code{Z} and a compiled CovarianceFactor v2 descriptor in \code{covFactor}, for use inside \code{vsm}.
\end{Value}
%
\begin{SeeAlso}\relax
\code{\LinkA{metricm}{metricm}}, \code{\LinkA{sar}{sar}}, \code{\LinkA{car}{car}}, \code{\LinkA{vsm}{vsm}}.
\end{SeeAlso}
%
\begin{Examples}
\begin{ExampleCode}
## Not run: 
vsm(maternm(cbind(xcoord,ycoord)), ism(genotype))
vsm(maternm(time), ism(id))

## End(Not run)
\end{ExampleCode}
\end{Examples}
\inputencoding{utf8}
\HeaderA{meffects\_mmes}{Marker effects recovered from a fitted relationship-based term}{meffects.Rul.mmes}
\keyword{models}{meffects\_mmes}
%
\begin{Description}\relax
Back-transforms the BLUPs of a random term fitted with a genomic relationship
matrix (GBLUP) into BLUPs of marker effects, optionally with standard errors
and back-solved GWAS-type tests. This is the post-fit equivalent of ASReml's
\code{mef} argument and does not require refitting the model.
\end{Description}
%
\begin{Usage}
\begin{verbatim}
meffects_mmes(object, term, M, method = c("A.mat", "custom"),
              min.MAF = 0, Z = NULL, scale = NULL, blend = 0, Gu = NULL,
              se = FALSE, chunk = 1000)
\end{verbatim}
\end{Usage}
%
\begin{Arguments}
\begin{ldescription}
\item[\code{object}] A fitted \code{\LinkA{mmes}{mmes}} model.
\item[\code{term}] Name (as in \code{names(object\$uList)}) or position of the random term fitted with \code{Gu}.
\item[\code{M}] Marker matrix (individuals by markers, coded -1/0/1 as for
\code{\LinkA{A.mat}{A.mat}}) with rownames matching the levels of the term. Use the
same matrix (same rows) that was used to build the relationship matrix,
because allele frequencies are computed from it. Rows that are not levels of
the term are ignored.
\item[\code{method}] \code{"A.mat"} reproduces exactly how \code{\LinkA{A.mat}{A.mat}}
builds G from \code{M} (imputation, \code{min.MAF} and monomorphic filters,
centring by \eqn{2p}{} and scaling by \eqn{2\sum p_jq_j}{}). \code{"custom"} uses
\code{Z} and \code{scale}.
\item[\code{min.MAF}] The \code{min.MAF} used in \code{A.mat}.
\item[\code{Z}] For \code{method="custom"}: transformed marker matrix with rownames
such that \eqn{G = scale \cdot ZZ^\top}{}.
\item[\code{scale}] For \code{method="custom"}: the scalar \eqn{s}{} in \eqn{G = sZZ^\top}{}.
\item[\code{blend}] \eqn{\lambda}{} if the fitted matrix was \eqn{(1-\lambda)G+\lambda I}{}.
Small bending constants added only for invertibility can be left at 0.
\item[\code{Gu}] Optional relationship matrix (or inverse, with \code{attr(Gu,"inverse")=TRUE}).
By default the inverse used in the fit is rebuilt from the model inputs.
\item[\code{se}] If \code{TRUE}, add standard errors, z statistics and p-values.
\item[\code{chunk}] Number of markers processed per block when \code{se=TRUE}.
\end{ldescription}
\end{Arguments}
%
\begin{Details}\relax
If the fitted relationship matrix is \eqn{G=(1-\lambda)sZZ^\top+\lambda I}{}, the
BLUP of the marker effects is
\deqn{\hat a=(1-\lambda)s\,Z^\top G^{-1}\hat u,}{}
using the same \eqn{G^{-1}}{} as the fit, so \eqn{Z\hat a=\hat u-\lambda G^{-1}\hat u}{}
(exactly \eqn{\hat u}{} when \eqn{\lambda=0}{}). For structured terms such as
\code{vsm(dsm(env), ism(id), Gu=Gi)} effects are returned per coordinate
(environment, trait, ...).

With \code{se=TRUE}, following Gualdron-Duarte et al. (2014),
\deqn{\mathrm{Var}(\hat a)=H^\top(\sigma^2_c G - C^{uu})H,\qquad H=(1-\lambda)sG^{-1}Z,}{}
where \eqn{\sigma^2_c}{} is the variance of the coordinate and \eqn{C^{uu}}{} the
prediction error covariance of the term; \eqn{z=\hat a/\mathrm{SE}}{} gives
back-solved GWAS tests. Only diagonal elements are formed: the \eqn{C^{uu}}{}
part uses solves with the stored coefficient matrix in blocks of \code{chunk}
markers. The cost grows as \eqn{O(mn^2)}{} for m markers and n individuals.

The rebuild of \eqn{G^{-1}}{} requires the data used in the original call; supply
\code{Gu} otherwise. Markers removed by the filters are returned as \code{NA}.
\end{Details}
%
\begin{Value}
A data frame with columns \code{marker}, \code{coordinate} (structured
terms only), \code{effect} and, if \code{se=TRUE}, \code{se}, \code{z} and
\code{p.value}. Attributes \code{scale} and \code{markersKept}.
\end{Value}
%
\begin{References}\relax
Gualdron-Duarte JL, Cantet RJC, Bates RO, Ernst CW, Raney NE, Steibel JP (2014).
Rapid screening for phenotype-genotype associations by linear transformations
of genomic evaluations. BMC Bioinformatics 15: 246.
\end{References}
%
\begin{SeeAlso}\relax
\code{\LinkA{A.mat}{A.mat}}, \code{\LinkA{mmes}{mmes}}, \code{\LinkA{vsm}{vsm}}
\end{SeeAlso}
%
\begin{Examples}
\begin{ExampleCode}
set.seed(7)
M <- matrix(sample(c(-1, 0, 1), 40 * 80, TRUE), 40)
rownames(M) <- paste0("g", 1:40); colnames(M) <- paste0("snp", 1:80)
G <- 0.95 * A.mat(M) + 0.05 * diag(40)
Gi <- solve(G); attr(Gi, "inverse") <- TRUE
d <- data.frame(id = factor(rep(rownames(M), each = 2), levels = rownames(M)))
d$y <- rnorm(80)
fit <- mmes(y ~ 1,
  random = ~ vsm(ism(id), Gu = Gi),
  data = d,
  control = mmes_control(diagnostics = list(verbose = FALSE)))
head(meffects_mmes(fit,
  1,
  M,
  blend = 0.05,
  se = TRUE))
\end{ExampleCode}
\end{Examples}
\inputencoding{utf8}
\HeaderA{metricm}{Distance-based (metric) spatial correlation structures}{metricm}
\keyword{models}{metricm}
%
\begin{Description}\relax
Exponential, power, gaussian, spherical and circular correlation structures
computed from distances between locations, isotropic or anisotropic. The
returned dimensionless covariance shape is intended for use inside
\code{\LinkA{vsm}{vsm}}, which supplies the single overall variance scale.
\end{Description}
%
\begin{Usage}
\begin{verbatim}
metricm(x, model = c("exponential", "power", "gaussian", "spherical", "circular"),
        range = NULL, rho = NULL, anisotropy = c("none", "product", "geometric"),
        angle = pi/2, ratio = 1.5, metric = c("euclidean", "manhattan"),
        distance = NULL, fixed = NULL)
\end{verbatim}
\end{Usage}
%
\begin{Arguments}
\begin{ldescription}
\item[\code{x}] Numeric coordinate vector (1-D) or matrix/data frame whose rows give coordinates. Repeated coordinates share one covariance level.
\item[\code{model}] Correlation function, see Details.
\item[\code{range}] Starting range \eqn{\phi}{} (one per axis with product anisotropy). Defaults to the median positive distance.
\item[\code{rho}] Starting correlation per unit distance for \code{model="power"}, in (0,1). Defaults to \eqn{e^{-1}}{} at the median distance.
\item[\code{anisotropy}] \code{"none"}; \code{"product"}, one parameter per axis, \eqn{K=f_x(|dx|)f_y(|dy|)}{}; or \code{"geometric"}, coordinates rotated by \code{angle} and the second rotated axis divided by \code{ratio}. Anisotropy requires 2-D coordinates.
\item[\code{angle}] Starting anisotropy angle in (0, pi).
\item[\code{ratio}] Starting positive anisotropy ratio.
\item[\code{metric}] Distance metric for isotropic models; \code{"manhattan"} is only valid (positive definite) for the exponential and power models.
\item[\code{distance}] Optional symmetric q by q distance matrix with zero diagonal, aligned to the unique locations (isotropic models only).
\item[\code{fixed}] Logical vector, one value per parameter (in the order range/rho, then angle and ratio), recycled if of length one.
\end{ldescription}
\end{Arguments}
%
\begin{Details}\relax
With \eqn{h=d/\phi}{}:

\Tabular{llll}{
\code{model} & \eqn{K(d)}{} & validity & ASReml \\{}
exponential & \eqn{e^{-h}}{} & any dimension & \code{exp}, \code{iexp} \\{}
power & \eqn{\rho^{d}}{} & any dimension & \code{exp}, \code{iexp} (same parameterization) \\{}
gaussian & \eqn{e^{-h^2}}{} & any dimension & \code{gau}, \code{igau} \\{}
spherical & \eqn{1-\frac32h+\frac12h^3}{} for \eqn{h<1}{}, else 0 & up to 3-D & \code{isp} \\{}
circular & \eqn{\frac{2}{\pi}(\arccos h-h\sqrt{1-h^2})}{} for \eqn{h<1}{}, else 0 & up to 2-D & \code{cir} \\{}
}
Product anisotropy (exponential, power and gaussian) corresponds to ASReml
\code{aexp}/\code{agau}; geometric anisotropy corresponds to ASReml's
anisotropic Matern and is also available in \code{\LinkA{maternm}{maternm}}. The power
model is the exponential model with \eqn{\phi=-1/\log\rho}{}, and the
exponential model equals \code{maternm(nu=0.5)}.

Ranges and ratios are optimized on log scales, correlations and angles through
bounded logits, with analytic derivatives. The structure is dense over the q
locations, so it is intended for up to a few thousand locations. On a complete
row-by-column grid, product anisotropy is a Kronecker product of two 1-D
structures; \code{vsm(metricm(row, "power"), metricm(col, "power"), ism(units))}
fits the same model while keeping the compact grid engine. The gaussian model
is nearly singular for close locations; keep a residual (nugget) term. With
geometric anisotropy the angle is not identifiable when the ratio is close to
one, and the ratio can drift to extreme values when the data are isotropic.
\end{Details}
%
\begin{Value}
A list with the incidence matrix \code{Z} and a CovarianceFactor descriptor \code{covFactor}, for use inside \code{vsm}.
\end{Value}
%
\begin{SeeAlso}\relax
\code{\LinkA{maternm}{maternm}}, \code{\LinkA{ar1m}{ar1m}}, \code{\LinkA{sar}{sar}}, \code{\LinkA{car}{car}}, \code{\LinkA{vsm}{vsm}}.
\end{SeeAlso}
%
\begin{Examples}
\begin{ExampleCode}
set.seed(1)
co <- cbind(x = runif(30, 0, 10), y = runif(30, 0, 10))
d <- data.frame(x = rep(co[, 1], each = 2), y = rep(co[, 2], each = 2), one = factor(1))
d$yield <- rnorm(60)
fit <- mmes(yield ~ 1,
  random = ~ vsm(metricm(cbind(x, y), model = "power"), ism(one)),
  data = d,
  control = mmes_control(diagnostics = list(verbose = FALSE)))
covparams_mmes(fit)
\end{ExampleCode}
\end{Examples}
\inputencoding{utf8}
\HeaderA{mmer}{\strong{m}ixed \strong{m}odel \strong{e}quations for \strong{r} records}{mmer}
%
\begin{Description}\relax

THIS FUNCTION HAS BEEN DEPRECATD PLEASE EXPLORE \code{\LinkA{mmes}{mmes}}. The mmer 
function uses the Direct-Inversion Newton-Raphson or Average Information coded in 
C++ using the Armadillo library to optimize dense matrix operations common in
genomic selection models. These algorithms are \strong{intended to be used for
problems of the type c > r (more coefficients to estimate than records in the
dataset) and/or dense matrices}. For problems with sparse data, or problems of 
the type r > c (more records in the dataset than coefficients to estimate), the 
MME-based algorithm in the \code{\LinkA{mmes}{mmes}} function is faster and we recommend 
to shift to use that function.

\end{Description}
%
\begin{Usage}
\begin{verbatim}

mmer(fixed, random, rcov, data, weights, W, nIters=20, tolParConvLL = 1e-03, 
     tolParInv = 1e-06, init=NULL, constraints=NULL,method="NR", getPEV=TRUE,
     naMethodX="exclude", naMethodY="exclude",returnParam=FALSE, 
     dateWarning=TRUE,date.warning=TRUE,verbose=TRUE, reshapeOutput=TRUE, stepWeight=NULL,
     emWeight=NULL, contrasts=NULL)
\end{verbatim}
\end{Usage}
%
\begin{Arguments}
\begin{ldescription}
\item[\code{fixed}] A formula specifying the \strong{response variable(s)} \strong{and fixed effects}, e.g.:

\emph{response \textasciitilde{} covariate} for univariate models

\emph{cbind(response.i,response.j) \textasciitilde{} covariate} for multivariate models

The \code{fcm} function can be used to constrain fixed effects in multi-response models.



\item[\code{random}] A formula specifying the name of the \strong{random effects}, 
e.g. \emph{random= \textasciitilde{} genotype + year}.

Useful functions can be used to fit heterogeneous variances and other special
models (\emph{see 'Special Functions' in the Details section for more information}):

\code{\LinkA{vsr}{vsr}(...,Gu,Gti,Gtc)} is the main function to specify variance 
models and special structures for random effects. On the ... argument you
provide the unknown variance-covariance structures (e.g., usr,dsr,atr,csr)
and the random effect where such covariance structure will be used (the 
random effect of interest). Gu is used to provide known covariance matrices 
among the levels of the random effect, Gti initial values and Gtc for constraints. 
Auxiliar functions for building the variance models are:

** \code{\LinkA{dsr}{dsr}(x)}, \code{\LinkA{usr}{usr}(x)}, \code{\LinkA{csr}{csr}(x)} and
\code{\LinkA{atr}{atr}(x,levs)} can be used to specify unknown diagonal, 
unstructured and customized unstructured and diagonal covariance structures
to be estimated by REML.

** \code{\LinkA{unsm}{unsm}(x)}, \code{\LinkA{fixm}{fixm}(x)} and \code{\LinkA{diag}{diag}(x)} 
can be used to build easily matrices to specify constraints in the Gtc argument
of the \code{\LinkA{vsr}{vsr}()} function.

** \code{\LinkA{overlay}{overlay}()}, \code{\LinkA{spl2Da}{spl2Da}()}, \code{\LinkA{spl2Db}{spl2Db}()},
and \code{\LinkA{leg}{leg}()} functions can be used to specify overlayed 
of design matrices of random effects, two dimensional spline and random 
regression models within the \code{\LinkA{vsr}{vsr}()} function.

\code{gvsr(...,Gu,Guc,Gti,Gtc)} is an alternative function to specify general
variance structures between different random effects. An special case in the 
indirect genetic effect models. Is similar to the vsr function but in the ... 
argument the different random effects are provided.


\item[\code{rcov}] A formula specifying the name of the \strong{error term}, e.g., \emph{rcov= \textasciitilde{} units}.

Special heterogeneous and special variance models and constraints for the
residual part are the same used on the random term but the name of the random 
effect is always "units" which can be thought as a column with as many levels
as rows in the data, e.g., \emph{rcov=\textasciitilde{}vsr(dsr(covariate),units)}


\item[\code{data}] A data frame containing the variables specified in the formulas for response, fixed, and random effects.
\item[\code{weights}] Name of the covariate for weights. To be used for the product R = Wsi*R*Wsi, where * is the matrix product, Wsi is the square root of the inverse of W and R is the residual matrix.
\item[\code{W}] Alternatively, instead of providing a vector of weights the user can specify an entire W matrix (e.g., when covariances exist). To be used first to produce Wis = solve(chol(W)), and then calculate R = Wsi*R*Wsi.t(), where * is the matrix product, and R is the residual matrix. Only one of the arguments weights or W should be used. If both are indicated W will be given the preference.
\item[\code{nIters}] Maximum number of iterations allowed.
\item[\code{tolParConvLL}] Convergence criteria for the change in log-likelihood.
\item[\code{tolParInv}] Tolerance parameter for matrix inverse used when singularities are encountered in the estimation procedure.
\item[\code{init}] Initial values for the variance components. By default this is NULL and initial values for the variance components are provided by the algorithm, but in case the user want to provide initial values for ALL var-cov components this argument is functional. It has to be provided as a list, where each list element corresponds to one random effect (1x1 matrix) and if multitrait model is pursued each element of the list is a matrix of variance covariance components among traits for such random effect. Initial values can also be provided in the Gti argument of the \LinkA{vsr}{vsr} function. Is highly encouraged to use the Gti and Gtc arguments of the \LinkA{vsr}{vsr} function instead of this argument, but these argument can be used to provide all initial values at once
\item[\code{constraints}] When initial values are provided these have to be accompanied by their constraints. See the \LinkA{vsr}{vsr} function for more details on the constraints. Is highly encouraged to use the Gti and Gtc arguments of the \LinkA{vsr}{vsr} function instead of this argument but these argument can be used to provide all constraints at once.
\item[\code{method}] This refers to the method or algorithm to be used for estimating variance components. Direct-inversion Newton-Raphson \strong{NR} and Average Information \strong{AI} (Tunnicliffe 1989; Gilmour et al. 1995; Lee et al. 2015).
\item[\code{getPEV}] A TRUE/FALSE value indicating if the program should return the predicted error variance and variance for random effects. This option is provided since this can take a long time for certain models where p is > n by a big extent.
\item[\code{naMethodX}] One of the two possible values; "include" or "exclude". If "include" is selected then the function will impute the X matrices for fixed effects with the median value. If "exclude" is selected it will get rid of all rows with missing values for the X (fixed) covariates. The default is "exclude". The "include" option should be used carefully.
\item[\code{naMethodY}] One of the three possible values; "include", "include2" or "exclude" (default) to treat the observations in response variable to be used in the estimation of variance components. The first option "include" will impute the response variables for all rows with the median value, whereas "include2" imputes the responses only for rows where there is observation(s) for at least one of the responses (only available in the multi-response models). If "exclude" is selected (default) it will get rid of rows in response(s) where missing values are present for at least one of the responses.
\item[\code{returnParam}] A TRUE/FALSE value to indicate if the program should return the parameters to be used for fitting the model instead of fitting the model.
\item[\code{dateWarning}] A TRUE/FALSE value to indicate if the program should warn you when is time to update the sommer package.
\item[\code{date.warning}] A TRUE/FALSE value to indicate if the program should warn you when is time to update the sommer package. This argument will be removed soon, just left for backcompatibility.
\item[\code{verbose}] A TRUE/FALSE value to indicate if the program should return the progress of the iterative algorithm.
\item[\code{reshapeOutput}] A TRUE/FALSE value to indicate if the output should be reshaped to be easier to interpret for the user, some information is missing from the multivariate models for an easy interpretation.
\item[\code{stepWeight}] A vector of values (of length equal to the number of iterations) indicating the weight used to multiply the update (delta) for variance components at each iteration. If NULL the 1st iteration will be multiplied by 0.5, the 2nd by 0.7, and the rest by 0.9. This argument can help to avoid that variance components go outside the parameter space in the initial iterations which doesn't happen very often with the NR method but it can be detected by looking at the behavior of the likelihood. In that case you may want to give a smaller weight to the initial 8-10 iterations.
\item[\code{emWeight}] A vector of values (of length equal to the number of iterations) indicating with values between 0 and 1 the weight assigned to the EM information matrix. And the values 1 - emWeight will be applied to the NR/AI information matrix to produce a joint information matrix.
\item[\code{contrasts}] an optional list. See the contrasts.arg of model.matrix.default.

\end{ldescription}
\end{Arguments}
%
\begin{Details}\relax

The use of this function requires a good understanding of mixed models. Please review the 'sommer.quick.start' vignette and pay attention to details like format of your random and fixed variables (e.g. character and factor variables have different properties when returning BLUEs or BLUPs, please see the 'sommer.changes.and.faqs' vignette). 

\strong{For tutorials} on how to perform different analysis with sommer please look at the vignettes by typing in the terminal:

vignette("v1.sommer.quick.start")

vignette("v2.sommer.changes.and.faqs")

vignette("v3.sommer.qg")

vignette("v4.sommer.gxe")


\strong{Citation}

Type \emph{citation("sommer")} to know how to cite the sommer package in your publications. 

\strong{Special variance structures}

\code{\LinkA{vsr}{vsr}(\LinkA{atr}{atr}(x,levels),y)}   

can be used to specify heterogeneous variance for the "y" covariate at specific levels of the covariate "x", e.g., \emph{random=\textasciitilde{}vsr(at(Location,c("A","B")),ID)} fits a variance component for ID at levels A and B of the covariate Location.

\code{\LinkA{vsr}{vsr}(\LinkA{dsr}{dsr}(x),y)}   

can be used to specify a diagonal covariance structure for the "y" covariate for all levels of the covariate "x", e.g., \emph{random=\textasciitilde{}vsr(dsr(Location),ID)}  fits a variance component for ID at all levels of the covariate Location.

\code{\LinkA{vsr}{vsr}(\LinkA{usr}{usr}(x),y)}   

can be used to specify an unstructured covariance structure for the "y" covariate for all levels of the  covariate "x", e.g., \emph{random=\textasciitilde{}vsr(usr(Location),ID)}  fits variance and covariance components for ID at all levels of the covariate Location.

\code{\LinkA{vsr}{vsr}(\LinkA{overlay}{overlay}(...,rlist=NULL,prefix=NULL))}   

can be used to specify overlay of design matrices between consecutive random effects specified, e.g., \emph{random=\textasciitilde{}vsr(overlay(male,female))} overlays (overlaps) the incidence matrices for the male and female random effects to obtain a single variance component for both effects. The `rlist` argument is a list with each element being a numeric value that multiplies the incidence matrix to be overlayed. See \code{\LinkA{overlay}{overlay}} for details.Can be combined with vsr().

\code{\LinkA{vsr}{vsr}(\LinkA{leg}{leg}(x,n),y)}

can be used to fit a random regression model using a numerical variable \code{x} that marks the trayectory for the random effect \code{y}. The leg function can be combined with the special functions \code{dsr}, \code{usr} \code{at} and \code{csr}. For example \emph{random=\textasciitilde{}vsr(leg(x,1),y)} or \emph{random=\textasciitilde{}vsr(usr(leg(x,1)),y)}.

\code{\LinkA{vsr}{vsr}(x,Gtc=fcm(v))}

can be used to constrain fixed effects in the multi-response mixed models. This is a vector that specifies if the fixed effect is to be estimated for such trait. For example \emph{fixed=cbind(response.i, response.j)\textasciitilde{}vsr(Rowf, Gtc=fcm(c(1,0)))} means that the fixed effect Rowf should only be estimated for the first response and the second should only have the intercept.

\code{gvsr(x,y)}  

can be used to fit variance and covariance parameters between two or more random effects. For example, indirect genetic effect models.

\code{\LinkA{spl2Da}{spl2Da}(x.coord, y.coord, at.var, at.levels))}

can be used to fit a 2-dimensional spline (e.g., spatial modeling) using coordinates \code{x.coord} and \code{y.coord} (in numeric class) assuming a single variance component. The 2D spline can be fitted at specific levels using the \code{at.var} and \code{at.levels} arguments. For example \emph{random=\textasciitilde{}spl2Da(x.coord=Row,y.coord=Range,at.var=FIELD)}.

\code{\LinkA{spl2Db}{spl2Db}(x.coord, y.coord, at.var, at.levels))}

can be used to fit a 2-dimensional spline (e.g., spatial modeling) using coordinates \code{x.coord} and \code{y.coord} (in numeric class) assuming multiple variance components. The 2D spline can be fitted at specific levels using the \code{at.var} and \code{at.levels} arguments. For example \emph{random=\textasciitilde{}spl2Db(x.coord=Row,y.coord=Range,at.var=FIELD)}.

\strong{S3 methods}

S3 methods are available for some parameter extraction such as \code{\LinkA{fitted.mmer}{fitted.mmer}}, \code{\LinkA{residuals.mmer}{residuals.mmer}}, \code{\LinkA{summary.mmer}{summary.mmer}}, \code{\LinkA{randef}{randef}}, \code{\LinkA{coef.mmer}{coef.mmer}}, \code{\LinkA{anova.mmer}{anova.mmer}}, \code{\LinkA{plot.mmer}{plot.mmer}}, and \code{\LinkA{predict.mmer}{predict.mmer}} to obtain adjusted means. In addition, the \code{\LinkA{vpredict}{vpredict}} function (replacement of the pin function) can be used to estimate standard errors for linear combinations of variance components (e.g., ratios like h2). 

\strong{Additional Functions}

Additional functions for genetic analysis have been included such as relationship matrix building (\code{\LinkA{A.mat}{A.mat}}, \code{\LinkA{D.mat}{D.mat}}, \code{\LinkA{E.mat}{E.mat}}, \code{\LinkA{H.mat}{H.mat}}), build a genotypic hybrid marker matrix (\code{\LinkA{build.HMM}{build.HMM}}), plot of genetic maps (\code{\LinkA{map.plot}{map.plot}}), and manhattan plots (\code{\LinkA{manhattan}{manhattan}}). If you need to build a pedigree-based relationship matrix use the \code{getA} function from the pedigreemm package. 

\strong{Bug report and contact}

If you have any technical questions or suggestions please post it in https://stackoverflow.com or https://stats.stackexchange.com 

If you have any bug report please go to https://github.com/covaruber/sommer or send me an email to address it asap, just make sure you have read the vignettes carefully before sending your question.

\strong{Example Datasets}

The package has been equiped with several datasets to learn how to use the sommer package: 

* \code{\LinkA{DT\_halfdiallel}{DT.Rul.halfdiallel}}, \code{\LinkA{DT\_fulldiallel}{DT.Rul.fulldiallel}} and \code{\LinkA{DT\_mohring}{DT.Rul.mohring}} datasets have examples to fit half and full diallel designs. 

* \code{\LinkA{DT\_h2}{DT.Rul.h2}} to calculate heritability

* \code{\LinkA{DT\_cornhybrids}{DT.Rul.cornhybrids}} and \code{\LinkA{DT\_technow}{DT.Rul.technow}} datasets to perform genomic prediction in hybrid single crosses 

* \code{\LinkA{DT\_wheat}{DT.Rul.wheat}} dataset to do genomic prediction in single crosses in species displaying only additive effects.

* \code{\LinkA{DT\_cpdata}{DT.Rul.cpdata}} dataset to fit genomic prediction models within a biparental population coming from 2 highly heterozygous parents including additive, dominance and epistatic effects. 

* \code{\LinkA{DT\_polyploid}{DT.Rul.polyploid}} to fit genomic prediction and GWAS analysis in polyploids. 

* \code{\LinkA{DT\_gryphon}{DT.Rul.gryphon}} data contains an example of an animal model including pedigree information.

* \code{\LinkA{DT\_btdata}{DT.Rul.btdata}} dataset contains an animal (birds) model.

* \code{\LinkA{DT\_legendre}{DT.Rul.legendre}} simulated dataset for random regression model.

* \code{\LinkA{DT\_sleepstudy}{DT.Rul.sleepstudy}} dataset to know how to translate lme4 models to sommer models.

* \code{\LinkA{DT\_ige}{DT.Rul.ige}} dataset to show how to fit indirect genetic effect models.

\strong{Models Enabled}

For details about the models enabled and more information about the covariance structures please check the help page of the package (\code{\LinkA{sommer}{sommer}}).

\end{Details}
%
\begin{Value}
If all parameters are correctly indicated the program will return a list with the following information:

\begin{ldescription}
\item[\code{Vi}] the inverse of the phenotypic variance matrix V\textasciicircum{}- = (ZGZ+R)\textasciicircum{}-1
\item[\code{P}] the projection matrix  Vi - [Vi*(X*Vi*X)\textasciicircum{}-*Vi] 
\item[\code{sigma}] a list with the values of the variance-covariance components with one list element for each random effect.
\item[\code{sigma\_scaled}] a list with the values of the scaled variance-covariance components with one list element for each random effect.
\item[\code{sigmaSE}] Hessian matrix containing the variance-covariance for the variance components. SE's can be obtained taking the square root of the diagonal values of the Hessian.
\item[\code{Beta}] a data frame for trait BLUEs (fixed effects).
\item[\code{VarBeta}] a variance-covariance matrix for trait BLUEs
\item[\code{U}] a list (one element for each random effect) with a data frame for trait BLUPs.
\item[\code{VarU}] a list (one element for each random effect) with the variance-covariance matrix for trait BLUPs.
\item[\code{PevU}] a list (one element for each random effect) with the predicted error variance matrix for trait BLUPs.
\item[\code{fitted}]  Fitted values y.hat=XB
\item[\code{residuals}] Residual values e = Y - XB
\item[\code{AIC}] Akaike information criterion
\item[\code{BIC}] Bayesian information criterion
\item[\code{convergence}] a TRUE/FALSE statement indicating if the model converged.
\item[\code{monitor}] The values of log-likelihood and variance-covariance components across iterations during the REML estimation.
\item[\code{percChange}] The percent change of variance components across iterations. There should be one column less than the number of iterations. Calculated as percChange = ((x\_i/x\_i-1) - 1) * 100 where i is the ith iteration.
\item[\code{dL}] The vector of first derivatives of the likelihood with respect to the ith variance-covariance component.
\item[\code{dL2}] The matrix of second derivatives of the likelihood with respect to the i.j th variance-covariance component.
\item[\code{method}] The method for extimation of variance components specified by the user.
\item[\code{call}] Formula for fixed, random and rcov used.
\item[\code{constraints}] contraints used in the mixed models for the random effects.
\item[\code{constraintsF}] contraints used in the mixed models for the fixed effects.
\item[\code{data}] The dataset used in the model after removing missing records for the response variable.
\item[\code{dataOriginal}] The original dataset used in the model.
\item[\code{terms}] The name of terms for responses, fixed, random and residual effects in the model.
\item[\code{termsN}] The number of effects associated to fixed, random and residual effects in the model.
\item[\code{sigmaVector}] a vectorized version of the sigma element (variance-covariance components) to match easily the standard errors of the var-cov components stored in the element sigmaSE.
\item[\code{reshapeOutput}] The value provided to the mmer function for the argument with the same name.

\end{ldescription}
\end{Value}
%
\begin{Author}\relax
Giovanny Covarrubias-Pazaran
\end{Author}
%
\begin{References}\relax

Covarrubias-Pazaran G. Genome assisted prediction of quantitative traits using the R package sommer. PLoS ONE 2016, 11(6): doi:10.1371/journal.pone.0156744   

Covarrubias-Pazaran G. 2018. Software update: Moving the R package sommer to multivariate mixed models for genome-assisted prediction. doi: https://doi.org/10.1101/354639

Bernardo Rex. 2010. Breeding for quantitative traits in plants. Second edition. Stemma Press. 390 pp.

Gilmour et al. 1995. Average Information REML: An efficient algorithm for variance parameter estimation in linear mixed models. Biometrics 51(4):1440-1450.

Kang et al. 2008. Efficient control of population structure in model organism association mapping. Genetics 178:1709-1723.

Lee, D.-J., Durban, M., and Eilers, P.H.C. (2013). Efficient two-dimensional smoothing with P-spline ANOVA mixed models and nested bases. Computational Statistics and Data Analysis, 61, 22 - 37.

Lee et al. 2015. MTG2: An efficient algorithm for multivariate linear mixed model analysis based on genomic information. Cold Spring Harbor. doi: http://dx.doi.org/10.1101/027201.

Maier et al. 2015. Joint analysis of psychiatric disorders increases accuracy of risk prediction for schizophrenia, bipolar disorder, and major depressive disorder. Am J Hum Genet; 96(2):283-294.

Rodriguez-Alvarez, Maria Xose, et al. Correcting for spatial heterogeneity in plant breeding experiments with P-splines. Spatial Statistics 23 (2018): 52-71.

Searle. 1993. Applying the EM algorithm to calculating ML and REML estimates of variance components. Paper invited for the 1993 American Statistical Association Meeting, San Francisco.

Yu et al. 2006. A unified mixed-model method for association mapping that accounts for multiple levels of relatedness. Genetics 38:203-208.

Tunnicliffe W. 1989. On the use of marginal likelihood in time series model estimation. JRSS 51(1):15-27.

Zhang et al. 2010. Mixed linear model approach adapted for genome-wide association studies. Nat. Genet. 42:355-360.

\end{References}
%
\begin{Examples}
\begin{ExampleCode}

####=========================================####
#### For CRAN time limitations most lines in the 
#### examples are silenced with one '#' mark, 
#### remove them and run the examples
####=========================================####

####=========================================####
#### EXAMPLES
#### Different models with sommer
####=========================================####

data(DT_example, package="enhancer")
DT <- DT_example
head(DT)

####=========================================####
#### Univariate homogeneous variance models  ####
####=========================================####

## Compound simmetry (CS) model
ans1 <- mmer(Yield~Env,
             random= ~ Name + Env:Name,
             rcov= ~ units,
             data=DT)
summary(ans1)

####===========================================####
#### Univariate heterogeneous variance models  ####
####===========================================####

## Compound simmetry (CS) + Diagonal (DIAG) model
ans2 <- mmer(Yield~Env,
             random= ~Name + vsr(dsr(Env),Name),
             rcov= ~ vsr(dsr(Env),units),
             data=DT)
summary(ans2)

####===========================================####
####  Univariate unstructured variance models  ####
####===========================================####

ans3 <- mmer(Yield~Env,
             random=~ vsr(usr(Env),Name),
             rcov=~vsr(dsr(Env),units), 
             data=DT)
summary(ans3)



####==========================================####
#### Multivariate homogeneous variance models ####
####==========================================####

## Multivariate Compound simmetry (CS) model
DT$EnvName <- paste(DT$Env,DT$Name)
ans4 <- mmer(cbind(Yield, Weight) ~ Env,
              random= ~ vsr(Name, Gtc = unsm(2)) + vsr(EnvName,Gtc = unsm(2)),
              rcov= ~ vsr(units, Gtc = unsm(2)),
              data=DT)
summary(ans4)

####=============================================####
#### Multivariate heterogeneous variance models  ####
####=============================================####

## Multivariate Compound simmetry (CS) + Diagonal (DIAG) model
ans5 <- mmer(cbind(Yield, Weight) ~ Env,
              random= ~ vsr(Name, Gtc = unsm(2)) + vsr(dsr(Env),Name, Gtc = unsm(2)),
              rcov= ~ vsr(dsr(Env),units, Gtc = unsm(2)),
              data=DT)
summary(ans5)

####===========================================####
#### Multivariate unstructured variance models ####
####===========================================####

ans6 <- mmer(cbind(Yield, Weight) ~ Env,
              random= ~ vsr(usr(Env),Name, Gtc = unsm(2)),
              rcov= ~ vsr(dsr(Env),units, Gtc = unsm(2)),
              data=DT)
summary(ans6)

####=========================================####
####=========================================####
#### EXAMPLE SET 2
#### 2 variance components
#### one random effect with variance covariance structure
####=========================================####
####=========================================####

data("DT_cpdata", package="enhancer")
DT <- DT_cpdata
GT <- GT_cpdata
MP <- MP_cpdata
head(DT)
GT[1:4,1:4]
#### create the variance-covariance matrix
A <- A.mat(GT)
#### look at the data and fit the model
mix1 <- mmer(Yield~1,
             random=~vsr(id, Gu=A) + Rowf,
             rcov=~units,
             data=DT)
summary(mix1)$varcomp

#### multi trait example
mix2 <- mmer(cbind(Yield,color)~1,
              random = ~ vsr(id, Gu=A, Gtc = unsm(2)) + # unstructured at trait level
                            vsr(Rowf, Gtc=diag(2)) + # diagonal structure at trait level
                                vsr(Colf, Gtc=diag(2)), # diagonal structure at trait level
              rcov = ~ vsr(units, Gtc = unsm(2)), # unstructured at trait level
              data=DT)
summary(mix2)




\end{ExampleCode}
\end{Examples}
\inputencoding{utf8}
\HeaderA{mmes}{\strong{m}ixed \strong{m}odel \strong{e}quations \strong{s}olver}{mmes}
%
\begin{Description}\relax

Fits a generalized linear mixed model using model formulas, covariance
structures and observation data. Special covariance structures can be
specified with \code{\LinkA{vsm}{vsm}()}. Numerical fitting choices are supplied
through the strictly validated \code{control} object; see
\code{\LinkA{mmes\_control}{mmes.Rul.control}} and the \code{sommer.mmes.options} vignette.


\end{Description}
%
\begin{Usage}
\begin{verbatim}

mmes(fixed, random, rcov, data, W, weights=NULL,
     naMethodX="exclude", naMethodY="exclude",
     naMethodRandom="exclude", naMethodR="exclude",
     contrasts=NULL, family=stats::gaussian(), vcc=NULL,
     covPar=NULL, control=mmes_control(), .internal=NULL)
\end{verbatim}
\end{Usage}
%
\begin{Arguments}
\begin{ldescription}
\item[\code{fixed}] A formula specifying the \strong{response variable(s)} \strong{and fixed effects}, i.e:

\emph{response \textasciitilde{} covariate}



\item[\code{random}] A formula specifying the \strong{random effects}, e.g.,
\emph{random = \textasciitilde{} genotype + year}. A simple random term that is not wrapped
in \code{\LinkA{vsm}{vsm}} is treated as an identity-covariance random effect.

The \code{\LinkA{vsm}{vsm}} function is the main interface for specifying structured
covariance models for random effects. Its current covariance model is

\deqn{G = \sigma^2 (K_1 \otimes K_2 \otimes \cdots \otimes K_m) \otimes A,}{}

where \eqn{\sigma^2}{} is the single variance scale owned by \code{vsm()},
\eqn{K_1,\ldots,K_m}{} are covariance-shaping factors supplied by covariance
constructors, and \eqn{A}{} is the known covariance relationship for the final
random effect. In the Henderson implementation \code{Gu} supplies the
\strong{inverse} of \eqn{A}{}; it must have row and column names matching the
levels of the final random effect and \code{attr(Gu,"inverse")=TRUE}.

Within \code{vsm(...)}, the \strong{last} constructor supplies the incidence
matrix for the main random effect, while all preceding constructors define
covariance-shaping factors. The number of covariance factors is not limited.
For example,

\code{random = \textasciitilde{} vsm(dsm(Location), ism(Name))}

fits heterogeneous variances across \code{Location} for the random effect
\code{Name}, while

\code{random = \textasciitilde{} vsm(dsm(Environment), ar1m(Row), ar1m(Column),
  ism(Name), Gu=Ainv)}

specifies an arbitrary Kronecker product of covariance structures before the
relationship matrix for \code{Name}.

Covariance constructors currently available include:

\code{\LinkA{ism}{ism}} for identity covariance;

\code{\LinkA{dsm}{dsm}} and \code{\LinkA{atm}{atm}} for diagonal and selected-level
diagonal covariance structures;

\code{\LinkA{usm}{usm}} for an unstructured covariance matrix;

\code{\LinkA{csm}{csm}} and \code{\LinkA{corgm}{corgm}} for correlation structures;

\code{\LinkA{ar1m}{ar1m}}, \code{\LinkA{ar2m}{ar2m}}, and \code{\LinkA{ar3m}{ar3m}} for
autoregressive covariance structures;

\code{\LinkA{mam}{mam}} (including \code{ma1m} and \code{ma2m}) for moving-average
covariance structures;

\code{\LinkA{toeplitzm}{toeplitzm}} for a general positive-definite Toeplitz correlation
structure;

\code{\LinkA{fam}{fam}} for factor-analytic covariance;

\code{\LinkA{rrm}{rrm}} for a reduced-rank random covariance with specific
variances fixed at zero; use \code{fam(..., variance="homogeneous")}
for an estimated common specific variance;

\code{\LinkA{antem}{antem}} for antedependence covariance;

\code{\LinkA{maternm}{maternm}} for Matern spatial covariance;

\code{\LinkA{metricm}{metricm}} for exponential, power, gaussian, spherical and
circular distance-based correlations (isotropic or anisotropic);

\code{\LinkA{sar}{sar}} and \code{\LinkA{car}{car}} for simultaneous and conditional
autoregressive spatial covariance models; and

\code{\LinkA{ownm}{ownm}} for a fixed or user-defined covariance function.

Design-matrix utilities such as \code{\LinkA{overlay}{overlay}},
\code{\LinkA{spl2Dc}{spl2Dc}}, \code{\LinkA{leg}{leg}}, and
\code{\LinkA{redmm}{redmm}} can still be combined with \code{vsm()} when
appropriate.

See \code{\LinkA{vsm}{vsm}} and the individual covariance-constructor help pages
for details on parameterizations and examples.

A single relationship term may request Lee--van der Werf rotation with
\code{vsm(..., Gu=Gu, rotation=TRUE)}. With
\code{control\$formulation="henderson"}, random
coefficients are represented in the eigenbasis and the diagonal relationship
precision is used in the mixed model equations. With
\code{control\$formulation="direct"},
the equivalent transformed observation-covariance parameterization is used.
Public random effects, fitted values, and residuals are returned in the
original basis. The current implementation requires a complete balanced
Gaussian layout, one rotated relationship term, and
\code{control\$output\$inverse=0} during
fitting. Residual structures are allowed when they are rotation invariant:
residuals may be structured across rotation blocks (e.g.,
\code{rcov=\textasciitilde{}vsm(usm(trait), ism(units))} or
\code{rcov=\textasciitilde{}vsm(dsm(Env), ism(units))}), but must be independent and
identically structured across the levels of the rotated relationship term.
Residual structures that correlate observations of different levels (e.g.,
\code{ar1m()}, \code{maternm()}, \code{sar()}, or \code{car()} over plots),
residual variances that differ among levels within a rotation block, and
\code{W} matrices that are not diagonal with a constant weight within each
rotation block are rejected, as are non-Gaussian PQL fits.


\item[\code{rcov}] A formula specifying the \strong{residual covariance structure}.
The default is \emph{rcov = \textasciitilde{} units}, which fits a homogeneous residual
variance.

Structured residual covariance is specified with the same
\code{\LinkA{vsm}{vsm}} and CovarianceFactor constructors used for random effects.
The residual term should be expressed as a single \code{vsm()} term; arbitrary
covariance complexity is represented by placing multiple covariance factors
inside that term rather than by summing multiple residual terms.

The final term is normally \code{ism(units)}, where \code{units} is the
observation-level incidence generated internally by \code{mmes}. For example,

\code{rcov = \textasciitilde{} vsm(dsm(Environment), ism(units))}

fits heterogeneous residual variances across environments, while

\code{rcov = \textasciitilde{} vsm(ar1m(Row), ar1m(Column), ism(units))}

fits a separable row-by-column AR1 residual covariance, with one correlation
per direction shared by every other Kronecker factor. For section-specific
variances and correlations, use

\code{rcov = \textasciitilde{} dsumm(vsm(ar1m(Row), ar1m(Column), ism(units)), by = Trial)}

(see \code{\LinkA{dsumm}{dsumm}}). More generally,
structures such as \code{\LinkA{usm}{usm}}, \code{\LinkA{csm}{csm}},
\code{\LinkA{toeplitzm}{toeplitzm}}, \code{\LinkA{maternm}{maternm}}, \code{\LinkA{sar}{sar}},
\code{\LinkA{car}{car}}, and \code{\LinkA{ownm}{ownm}} can be used when their incidence
layout is appropriate for the residual model.

Residual covariance factors must identify exactly one covariance-product
coordinate for each observation. \code{mmes()} constructs the residual block
and local-coordinate indexing internally, so the data do \strong{not} need to
be manually sorted by the variables defining the residual structure.

\strong{Residual pairing key (multi-trait models).} For a long-format
multi-trait model the residual must know which records belong to the same
experimental unit. With \code{ism(units)} records are paired implicitly by
their order of appearance within each trait; a warning is issued when
repeated residual coordinates must be paired by row order. An explicit key is given by replacing
\code{units} with a column identifying the unit, e.g.

\code{rcov = \textasciitilde{} vsm(usm(trait), ism(record))}

where \code{record} is the column created by \code{\LinkA{stackTraits}{stackTraits}}.
Records sharing a key value form one residual block (missing traits are
simply absent from it); a trait may appear only once per key value.

As with random effects, \code{vsm()} owns one residual variance scale
\eqn{\sigma^2}{}; the preceding covariance constructors provide dimensionless
covariance shapes that are combined through an arbitrary Kronecker product.

See \code{\LinkA{vsm}{vsm}} and the individual covariance-constructor help pages
for details.


\item[\code{data}] Optional data frame containing variables used by the model.
Model expressions are evaluated using normal R scoping rules: variables are first
sought in \code{data} and may also be obtained from the environments associated
with the model formulas or from the calling environment. Consequently, auxiliary
objects such as relationship matrices, coordinate vectors, or grouping variables
do not have to be copied into \code{data}. If \code{data} is omitted, variables
are evaluated from the formula/calling environment.

Observation-level variables used by the response, fixed effects, random effects,
and residual covariance specification must nevertheless have compatible lengths.
Internally, \code{mmes()} constructs a common observation mask and applies it
consistently to all observation-level model components.
\item[\code{W}] Weights matrix (e.g., when covariance among plots exist). Internally
W is squared and inverted as Wsi = solve(chol(W)), then the residual matrix
is calculated as R = Wsi*O*Wsi.t(), where * is the matrix product, and O 
is the original residual matrix. Weights can be combined with any residual
structure, including Kronecker (multi-trait) residuals.
\item[\code{weights}] Optional one-sided formula describing independent row blocks
of the supplied \code{W} matrix, for example \code{weights = \textasciitilde{}trial}. Rows
with different formula groups must have zero cross-group entries in \code{W}.
This is structural metadata only: it does not generate numeric weights.
Declared blocks are processed independently by the Henderson weighting
path; without this formula, \code{W} is handled without a known grouping.
\item[\code{control}] A \code{sommer\_mmes\_control} object created by
\code{\LinkA{mmes\_control}{mmes.Rul.control}}. It owns numerical task, equation formulation,
representation, solver, determinant, score, curvature, covariance update,
proposal, convergence, output, execution and diagnostic settings. Its default
is \code{mmes\_control()}. Removed numerical arguments are not accepted by
\code{mmes()}; configure those choices in this object. See the control help
and vignette for defaults, mathematical ownership, eligibility and migration.

\item[\code{naMethodX}] Missing-data policy for variables entering the fixed-effect
design matrix. The default, \code{"exclude"}, marks observations with missing
fixed-effect information for removal from the common observation set.
\code{"include"} retains the historical sommer behavior for supported cases by
imputing missing fixed-effect covariates. This option should be used carefully,
because imputation changes the fitted design matrix.

\item[\code{naMethodY}] Missing-data policy for the response variable(s). The default,
\code{"exclude"}, removes observations with missing response information from the
common observation set. The historical \code{"include"} and \code{"include2"}
behaviors are retained where supported: \code{"include"} imputes missing responses,
whereas \code{"include2"} is intended for multivariate-response situations in which
an observation may be retained when at least one response is observed.

\item[\code{naMethodRandom}] Missing-data policy for observation-level variables used to
construct random-effect terms. The default is \code{"exclude"}. Under this policy,
an observation for which a required random-effect classification or covariate cannot
be evaluated is excluded from the common observation set. Random terms are evaluated
before the final observation subset is applied, so the same inclusion mask is used
consistently for the response, fixed-effect design, and all random-effect design
matrices.

\item[\code{naMethodR}] Missing-data policy for observation-level variables used to define
the residual covariance structure. The default is \code{"exclude"}. Observations
lacking information required to assign their residual covariance-product coordinate
are excluded from the common observation set. The residual block and local-coordinate
layout are then constructed on the retained observations.

\item[\code{contrasts}] an optional list. See the contrasts.arg of model.matrix.default.
\item[\code{vcc}] Optional data frame of equality/scaling constraints between
covariance parameters (columns \code{parameter}, \code{group} and optionally
\code{scale}), identified through the \code{vcParams} table returned by
\code{control\$output\$setup=TRUE} and stored in the fit. See \code{\LinkA{vcc}{vcc}}.

\item[\code{family}] A \code{stats::family()} object. The default Gaussian identity
family uses the ordinary linear mixed-model implementation. Other families are
fitted by penalized quasi-likelihood (PQL): each outer iteration forms an IRLS
working response and weights, then fits the resulting weighted Gaussian mixed
model with the Henderson solver. The initial implementation supports one numeric
response. Any family with a valid link can be used, e.g. \code{binomial()}
(logit, probit, cloglog), \code{poisson()}, \code{Gamma()},
\code{inverse.gaussian()}, \code{gaussian(link="log")}, \code{quasipoisson()}
and \code{MASS::negative.binomial(theta)}. Binomial counts out of a known
number of trials are supplied as a single proportion response built with
\code{\LinkA{binm}{binm}}, whose trials are used as prior weights. Binomial,
Poisson and negative binomial working residual dispersions are fixed to one.
PQL uses outer-loop controls from \code{control\$task\$pql}; compatibility with
update and solver requests is model-dependent.
For other families the reported \code{dispersion} is the Pearson estimate.
The negative binomial \eqn{\theta}{} is estimated (conditional on the random
effects) with \code{control\$task\$pql=list(estimateTheta=TRUE)} and returned in
\code{theta.nb} and \code{theta.nb.se}. Because the working-model
log-likelihood is not a GLMM likelihood, \code{llik}, \code{AIC} and
\code{BIC} are \code{NA} for PQL fits and \code{anova()} likelihood ratio
tests are refused; use \code{deviance} and Wald tests instead.
Formula offsets are included on the link scale. A supplied symmetric positive
definite \code{W} is combined with the IRLS precision after observation filtering,
and may be sparse and non-diagonal.

Long-format multi-trait data with a different family per trait are fitted by
supplying \code{family = \LinkA{familym}{familym}("trait", Disease = binomial(),
  Yield = gaussian())}. See \code{\LinkA{familym}{familym}} for the required residual
structure; in that case \code{dispersion} is a named vector with one value
per trait.

\item[\code{covPar}] Known covariance parameters for the solve task selected by
\code{control\$task\$mode="solve"}. Accepted forms are a compatible fitted
\code{mmes} object, a named list of term parameters, or a parameter data
frame. When \code{NULL}, covariance values must be fixed in the model.

\item[\code{.internal}] Reserved for package-internal recursive fitting state.
Users should not set this argument.
\end{ldescription}
\end{Arguments}
%
\begin{Details}\relax

\strong{Numerical controls}

\code{mmes()} keeps model specification separate from numerical policy.
The \code{control} argument defaults to \code{mmes\_control()}, whose
modules have strict schemas: unknown field names are errors. The default
configuration retains the former standard bundle: Henderson equations, REML,
exact AI curvature, EM recovery, no extrapolation, 30 inner iterations,
full covariance information, coefficient-matrix retention, and inherited
solver/resource defaults. The \code{mmes\_control} help gives every field,
default, module owner and validation rule.

\strong{Ownership and eligibility}

\begin{description}

\item[Task] \code{control\$task} selects fitting or known-covariance solving,
REML or ML, and outer PQL controls.
\item[Formulation] \code{control\$formulation} selects Henderson or direct
observation-space equations.
\item[Representation and solve] \code{control\$representation} describes
equation/covariance representation; \code{control\$solve} owns linear solver,
preconditioning, reuse and elimination requests.
\item[Evaluation] \code{control\$logdet} owns determinant methods;
\code{control\$score} owns score traces and covariance checks;
\code{control\$curvature} owns exact or hybrid AI and reduced-direction work.
\item[Optimization] \code{control\$update} owns EM/PXEM recovery, blending
and common-scale profiling. \code{control\$proposal} owns step scaling and
Aitken extrapolation. Aitken and hybrid AI are mutually exclusive.
\item[Convergence and outputs] \code{control\$convergence} owns inner
iteration, tolerances and stochastic accuracy; \code{control\$output} owns
coefficient/covariance uncertainty and setup output.
\item[Execution and diagnostics] \code{control\$execution} owns memory,
hardware, threading, precision and seeds. \code{control\$diagnostics} owns
verbosity, package-age warnings and profiling.

\end{description}


These modules do not form a Cartesian product of interchangeable algorithms.
Model-dependent admission remains authoritative; a requested mode may be
ineligible. Consult \code{\LinkA{mmes\_control}{mmes.Rul.control}} for prerequisites, exclusions
and current behavior. No new method combinations are implied by the modular API.
The internal \code{.mmes\_capabilities()} catalog is descriptive and does
not resolve eligibility or predict execution.

\end{Details}
%
\begin{Value}
If all parameters are correctly indicated the program will return a list with the following information:

\begin{ldescription}
\item[\code{data}] the dataset used in model fitting after application of the common
observation mask.
\item[\code{obsInfo}] an observation-level audit table describing the common missing-data
filter. It includes the original row index, validity indicators for
response, fixed, random, and residual model components, the final
inclusion indicator, and the reason for exclusion when applicable.
\item[\code{Dtable}] the table to be used for the predict function to help the program
recognize the factors available.
\item[\code{llik}] the vector of log-likelihoods across iterations
\item[\code{b}] the vector of fixed effect.
\item[\code{u}] the vector of random effect.
\item[\code{bu}] the vector of fixed and random effects together.
\item[\code{Ci}] the inverse of the coefficient matrix.
\item[\code{Ci\_11}] the inverse of the coefficient matrix pertaining to the fixed effects.
\item[\code{theta}] a list of estimated variance covariance matrices. Each element of 
the list corresponds to the different random and residual components
\item[\code{covPar}] the fitted covariance parameters in their descriptor-defined
reported coordinates.
\item[\code{control}] the validated modular numerical controls used for the fit.
See \code{\LinkA{mmes\_control}{mmes.Rul.control}}.
\item[\code{pcgRecycleState}] present only when \code{pcg\$exportRecycle=TRUE}.
A versioned serializable Krylov basis with coefficient-coordinate metadata,
for explicit import into a compatible PCG fit.
\item[\code{covarianceInformation}] \code{"none"} when covariance information was
explicitly omitted. The default retains full covariance information.
\item[\code{retainCoefficientMatrix}] FALSE when an opt-in C-free saddle-point fit
omitted the final compatibility matrix. In that mode \code{C} is NULL, and
C-based post-fit methods require a refit retaining C. The default retains C.
\item[\code{engineDiagnostics}] path-dependent numerical diagnostics.
\code{compression} reports method, original observations, groups, residual
classes, retained statistic bytes and setup time. \code{pcg} reports iterative
work, reuse and stochastic accuracy diagnostics for REML PCG.
\code{profile} contains detailed phase timings when explicitly enabled.
Retained statistic bytes are not process peak memory.
\item[\code{pcgMatrixFree}] logical indicating whether the matrix-free Kronecker PCG
path was used during optimization.
\item[\code{residualCovarianceCompact}] logical indicating that a complete, unweighted,
non-diagonal Kronecker residual covariance was evaluated directly from its
factors without materializing the full covariance matrix. In this case the
residual element of \code{theta} is a \eqn{0 \\times 0}{} matrix; use
\code{covparams\_mmes()} for fitted residual covariance parameters.
\item[\code{covParNative}] a data frame containing all fixed and estimated covariance
parameters in the native model-facing scale defined by each CovarianceFactor.
For example, \code{dsm()} entries are level-specific variances rather than
variance ratios. The table includes delta-method \code{StdError} and
\code{Zratio} columns when covariance information is retained, matching
\code{covparams\_mmes\_se()}. With \code{covarianceInformation="none"} these
columns are NA; \code{covparams\_mmes()} still returns point estimates.
\item[\code{theta\_se}] covariance-parameter uncertainty matrix; NULL when
\code{covarianceInformation="none"}.
\item[\code{InfMat}] information matrix; NULL when
\code{covarianceInformation="none"}.
\item[\code{monitor}] The values of the variance-covariance components across iterations 
during the REML estimation.
\item[\code{AIC}] Akaike information criterion
\item[\code{BIC}] Bayesian information criterion
\item[\code{convergence}] a TRUE/FALSE statement indicating if the model converged.
\item[\code{partitions}] a list where each element contains a matrix indicating where 
each random effect starts and ends.
\item[\code{partitionsX}] a list where each element contains a matrix indicating where 
each fixed effect starts and ends.
\item[\code{percDelta}] the matrix of percentage change in deltas (see \code{control\$convergence\$parameterTol}).
\item[\code{normMonitor}] the matrix of the three norms calculated (see \code{control\$convergence\$parameterTol}).
\item[\code{toBoundary}] the matrix of variance components that were forced to the
boundary across iterations.
\item[\code{Cchol}] the Cholesky decomposition of the coefficient matrix.
\item[\code{y}] the response vector.
\item[\code{W}] the column binded matrix W = [X Z] 
\item[\code{uList}] a list containing the BLUPs in data frame format where rows are 
levels of the random effects and column the different factors at 
which the random effect is fitted. This is specially useful for 
diagonal and unstructured models.
\item[\code{uPevList}] a list containing the BLUPs in data frame format where rows are 
levels of the random effects and column the different factors at 
which the random effect is fitted. This is specially useful for 
diagonal and unstructured models.
\item[\code{args}] the fixed, random and residual formulas from the mmes model.
\item[\code{constraints}] The vector of constraints.

\end{ldescription}
\end{Value}
%
\begin{Author}\relax

Coded by Giovanny Covarrubias-Pazaran with contributions of Christelle Fernandez 
Camacho, Johan Aparicio-Arce, and Claudio Flavio 5.1 jaja (laughing in Spanish).

\end{Author}
%
\begin{References}\relax

Covarrubias-Pazaran G. Genome assisted prediction of quantitative traits using 
the R package sommer. PLoS ONE 2016, 11(6): doi:10.1371/journal.pone.0156744

Jensen, J., Mantysaari, E. A., Madsen, P., and Thompson, R. (1997). Residual 
maximum likelihood estimation of (co) variance components in multivariate mixed 
linear models using average information. Journal of the Indian Society of
Agricultural Statistics, 49, 215-236.

Sanderson, C., \& Curtin, R. (2025). Armadillo: An Efficient Framework for 
Numerical Linear Algebra. arXiv preprint arXiv:2502.03000.

Gilmour et al. 1995. Average Information REML: An efficient algorithm for 
variance parameter estimation in linear mixed models. Biometrics 51(4):1440-1450.

\end{References}
%
\begin{Examples}
\begin{ExampleCode}

####=========================================####
#### For CRAN time limitations most lines in the
#### examples are silenced with one '#' mark,
#### remove them and run the examples
####=========================================####

data(DT_example, package="enhancer")
DT <- DT_example
head(DT)

####=========================================####
#### Univariate homogeneous variance models  ####
####=========================================####

## Compound simmetry (CS) model
ans1 <- mmes(Yield~Env,
             random= ~ Name + Env:Name,
             rcov= ~ units,
             data=DT)
summary(ans1)

## Known variances: solve the MME only (e.g. genetic evaluation)
sol1 <- mmes(Yield~Env,
             random= ~ Name + Env:Name,
             rcov= ~ units,
             data=DT, control=mmes_control(task=list(mode="solve")),
             covPar=ans1)
head(sol1$uList[[1]])



####===========================================####
#### Univariate heterogeneous variance models  ####
####===========================================####
DT=DT[with(DT, order(Env)), ]
## Compound simmetry (CS) + Diagonal (DIAG) model
ans2 <- mmes(Yield~Env,
             random= ~Name + vsm(dsm(Env),ism(Name)),
             rcov= ~ dsumm(vsm(ism(units)), by=Env),
             data=DT)
summary(ans2)

####===========================================####
####  Univariate unstructured variance models  ####
####===========================================####

ans3 <- mmes(Yield~Env,
             random=~ vsm(usm(Env),ism(Name)),
             rcov= ~ dsumm(vsm(ism(units)), by=Env),
             data=DT)
summary(ans3)






## simulated MET
nInds=300
nMarks=1000
nEnvs=3

#random population of nInds lines with nMarks markers
M <- matrix(rep(0,nInds*nMarks),nInds,nMarks)
for (i in 1:nInds) {
  M[i,] <- ifelse(runif(nMarks)<0.5,-1,1)
}
# compute the relationship matrix
MMT <- tcrossprod(M)
m <- sum(diag(MMT))/nrow(MMT)
MMT <- MMT/m
MMT <- MMT + diag(1e-05, ncol(MMT), ncol(MMT))
# get phenotypes
rg <- enhancer::simGECorMat(nEnvs, 1, mu=0.2, v=0.2)
rg <- Matrix::nearPD(rg)$mat
rg # desired genetic correlation
u <- mvrnorm(n = nMarks, rep(0, nEnvs), Sigma=rg) # marker effects
g <- crossprod(t(M),u) # genetic values
h2 <- 0.2  #heritability 
envEffect <- rnorm(nEnvs, mean = abs(rnorm(1))) # environment means
e <- mvrnorm(n = nInds, envEffect, Sigma=(1-h2)/h2*diag(diag(var(g)))) # errors
p = g + e # phenotypes
y <- as.vector(p) # phenotypes as vector
env <- sort(rep(paste0("e",1:nEnvs), nrow(g))) # environment labels
ind <- rep(paste0("i",1:nrow(g)), nEnvs ) # individuals labels
colnames(MMT) <- rownames(MMT) <- paste0("i",1:nrow(g)) # assign dimnames

MMTi <- solve(MMT + diag(1e-4,ncol(MMT),ncol(MMT)))
MMTi[lower.tri(MMTi)] <- t(MMTi)[lower.tri(MMTi)] # fill the lower triangular
MMTi <- as(as(as( MMTi,  "dMatrix"), "generalMatrix"), "CsparseMatrix")
attr(MMTi, 'inverse')=TRUE

DT <- data.frame(ind=ind, env=env, y=y)
# fit the models
mixm2 <- mmes(y ~ env,
              random=~ vsm(dsm(env), ism(ind), Gu=MMTi),
              data=DT
)
summary(mixm2)$varcomp



\end{ExampleCode}
\end{Examples}
\inputencoding{utf8}
\HeaderA{mmes\_control}{Modular Numerical Controls for Mixed Models}{mmes.Rul.control}
%
\begin{Description}\relax
Creates a strictly validated \code{sommer\_mmes\_control} object for the
numerical methods used by \code{\LinkA{mmes}{mmes}}. Model formulas remain separate.
Every module accepts only its documented fields; unknown module or field names
are errors. Constructing controls does not fit a model or establish that
requested methods are eligible for a particular model.
\end{Description}
%
\begin{Usage}
\begin{verbatim}
mmes_control(formulation = "henderson", task = list(),
             representation = list(), solve = list(), logdet = list(),
             score = list(), curvature = list(), update = list(),
             proposal = list(), convergence = list(), output = list(),
             execution = list(), diagnostics = list())
\end{verbatim}
\end{Usage}
%
\begin{Arguments}
\begin{ldescription}
\item[\code{formulation}] Equation domain: \code{"henderson"} (default) or
\code{"direct"}. Henderson forms coefficient-space mixed-model equations;
direct uses observation-space inversion.
\item[\code{task}] Named list with \code{mode}, \code{objective}, and \code{pql}.
Defaults: \code{mode="fit"}, \code{objective="reml"}, and
\code{pql=list(maxit=20L, tol=1e-5, estimateTheta=FALSE, warmStart=TRUE)}.
\code{mode} is \code{"fit"} or \code{"solve"}; \code{objective} is
\code{"reml"} or \code{"ml"}. The nested \code{pql} list accepts only
\code{maxit}, \code{tol}, \code{estimateTheta}, and \code{warmStart}, with
the defaults shown. These fields own task selection and outer PQL iteration,
not inner covariance iterations.
\item[\code{representation}] Named list with \code{compression="off"},
\code{factorScores="none"}, \code{equations="normal"},
\code{sparseFactors="off"}, and \code{matrixFree="auto"} by default,
unless corresponding \code{sommer.mme.*} options override experimental modes.
If \code{sommer.mme.saddlePoint} resolves to \code{"force"}, the default
\code{equations} becomes \code{"saddle-point"}.
\code{compression} is \code{"off"}, \code{"groups"}, or \code{"classes"};
\code{factorScores} is \code{"none"}, \code{"fixed-shape"}, or
\code{"profile"}; \code{equations} is \code{"normal"} or
\code{"saddle-point"}. The remaining fields use experimental mode values
\code{"auto"}, \code{"off"}, or \code{"force"} (logical TRUE/FALSE are
accepted). This module owns equation, compression, covariance-factor, and
factor-score representations.
\item[\code{solve}] Named list. \code{method="auto"} selects among
\code{"auto"}, \code{"ldlt"}, \code{"cholmod"}, and \code{"pcg"}.
PCG fields default to \code{tol=1e-8}, \code{maxIters=0L},
\code{preconditioner="diagonal"}, \code{nystromRank=32L},
\code{warmStart=FALSE}, \code{recycleRank=0L},
\code{recycleMethod="coarse"}, \code{recycleState=NULL},
\code{exportRecycle=FALSE}, \code{rhsMethod="independent"},
\code{blockSize=16L}, and \code{residualRefresh=0L}.
Also accepts the experimental solve requests \code{blockSchur="auto"},
\code{blockBanded="auto"}, \code{factorSchurMultiFactor="auto"}, and
\code{dFirstML="auto"}; global \code{sommer.mme.*} options may override
these mode defaults. The method owns checked linear solves,
preconditioning, reuse, and eligible elimination. PCG settings apply only
when an eligible PCG path uses them.
\item[\code{logdet}] Named list with \code{scaling="none"} and
\code{lanczosSteps=20L} by default. Scaling is \code{"none"} or
\code{"diagonal"}. This module owns determinant approximation and its
correction, not linear-solve accuracy.
\item[\code{score}] Named list with \code{traceProbes=8L},
\code{estimator="hutchinson"}, \code{rank=16L},
\code{covarianceChecks="auto"}, and \code{traceBatch="off"} by default.
The estimator is \code{"hutchinson"} or \code{"control-variate"}.
\code{covarianceChecks} and \code{traceBatch} accept experimental mode
values. This module owns covariance derivatives, trace estimation, and
covariance checks.
\item[\code{curvature}] Named list with \code{method="exact-ai"} and
\code{reduced="off"} by default. The method is \code{"exact-ai"} or
\code{"hybrid-ai"}; \code{reduced} accepts experimental mode values.
This module owns the AI metric and reduced-direction assembly.
\item[\code{update}] Named list with \code{recovery="em"},
\code{blendWeight=NULL}, and \code{commonScale="off"} by default.
\code{recovery} is \code{"em"} or \code{"pxem"}; \code{commonScale}
accepts experimental mode values. This module owns covariance-coordinate
updates, EM/PXEM recovery, blending, and common-scale profiling.
\item[\code{proposal}] Named list with \code{stepWeight=NULL} and
\code{extrapolation="none"} by default. Extrapolation is \code{"none"} or
\code{"aitken"}. This module owns proposal scaling and extrapolation.
Aitken extrapolation cannot be combined with
\code{curvature\$method="hybrid-ai"}.
\item[\code{convergence}] Named list with \code{maxIters=30L},
\code{likelihoodTol=1e-4}, \code{parameterTol=1e-4},
\code{accuracy="fixed"}, \code{maxTraceProbes=NULL},
\code{maxLanczosSteps=NULL}, \code{logDetTol=1e-3}, and
\code{scoreTol=1e-3}. Accuracy is \code{"fixed"} or \code{"adaptive"};
NULL stochastic limits resolve to the initial \code{traceProbes} and
\code{lanczosSteps}. This module owns inner iteration convergence and
stochastic accuracy requests.
\item[\code{output}] Named list with \code{inverse=0L}, \code{covariance="full"},
\code{coefficientMatrix=TRUE},
\code{setup=FALSE}, \code{selectedInverse="auto"}, and
\code{panelSize=64L} by default. \code{inverse} is 0 (no final coefficient
inverse or fit-time PEV computation),
1 (selected inverse and random-effect PEVs), or 2 (full inverse and
random-effect PEVs). Modes 1 and 2 compute PEVs during fitting; mode 1
does not form the complete coefficient-matrix inverse. With mode 0, use
\code{\LinkA{postPEV}{postPEV}} after fitting to compute PEVs or the full inverse,
provided \code{coefficientMatrix=TRUE};
\code{covariance} is \code{"full"} or \code{"none"}; \code{selectedInverse}
accepts experimental mode values. This module owns coefficient uncertainty,
covariance-parameter uncertainty, retained coefficient-matrix output, and
setup output.
\item[\code{execution}] Named list with \code{denseMemoryMB=1000},
\code{memoryPolicy="speed"}, \code{gpu=FALSE}, \code{gpuMinSize=256L},
\code{blasThreads=NULL}, \code{inverseTol=1e-6}, and \code{seed=0L} by
default. A non-NULL \code{sommer.mme.denseMemoryMB},
\code{sommer.mme.memoryPolicy}, \code{sommer.mme.gpu}, or
\code{sommer.mme.gpuMinSize} option supplies its corresponding default.
This module owns resource policy, hardware, threading, numerical inverse
tolerance, and stochastic seed.
\item[\code{diagnostics}] Named list with \code{verbose=TRUE},
\code{dateWarning=TRUE}, and \code{profile="off"} by default.
\code{profile} accepts experimental mode values. This module owns ordinary
fit messages, package-age warnings, and phase profiling.
\end{ldescription}
\end{Arguments}
%
\begin{Details}\relax
Module ownership describes mathematical responsibility, not a promise that
every cross-module combination is supported. For example, formulation and
solver choices are model-dependent; PCG has REML and output restrictions;
compression, saddle-point equations, factor-score augmentation, common-scale
profiling, and reduced AI each have additional prerequisites. Consult
\code{vignette("sommer.mmes.options", package="sommer")} and the internal
\code{.mmes\_capabilities()} catalog for the current conceptual decision graph
and capability inventory. The catalog is descriptive; fitting validation
remains authoritative. The modular API does not imply a kernel rewrite or
arbitrary method composition.

\code{task\$mode="solve"} selects a known-covariance path rather than ordinary
variance-component fitting; it requires an eligible Gaussian identity model
and has narrower PCG/output options. \code{task\$pql\$estimateTheta=TRUE}
requires a negative-binomial family. \code{representation\$compression} needs
eligible matching designs and independent residuals. Saddle-point equations,
factor-score representation, sparse factors, and matrix-free PCG each have
separate model and backend gates. PCG fitting is a Henderson REML path;
selected-inverse output mode 1 is not supported by PCG. The vignette and
capability catalog describe these prerequisites, but the fit-time validator
decides eligibility.

Within PCG, \code{solve\$preconditioner="nystrom"} uses
\code{solve\$nystromRank}; \code{solve\$recycleState} and
\code{solve\$exportRecycle} require \code{solve\$recycleMethod="krylov"} and a
positive \code{solve\$recycleRank}. \code{logdet\$scaling} and Lanczos depth
apply to stochastic determinant evaluation. \code{score\$rank} is used by the
control-variate estimator. Adaptive trace and Lanczos budgets apply to
stochastic PCG accuracy. \code{output\$panelSize} applies to full inverse
panels, not selected-inverse mode.

Experimental mode fields accept \code{"auto"}, \code{"off"}, or
\code{"force"}; logical TRUE and FALSE mean force and off. Package options
may supply defaults for these requests. \code{"force"} requests a path but
does not bypass mathematical, structural, or memory eligibility checks.
Nested lists reject unknown names and duplicate or empty names.

Defaults retain the established numerical bundle. Aitken belongs to
\code{proposal\$extrapolation}; hybrid AI belongs to
\code{curvature\$method}; they are mutually exclusive.

\strong{Migration from removed numerical arguments}

\code{mmes()} now takes numerical settings through \code{control}. The old
formals are not compatibility aliases and must not be passed to \code{mmes()}.
Map former settings to canonical module paths:
\begin{description}

\item[\code{henderson}] \code{formulation="henderson"} or
\code{formulation="direct"}.
\item[\code{solver}, \code{engine}] \code{solve\$method}; pass the complete
control object as \code{control=...}.
\item[\code{nIters}, \code{tolParConvLL}, \code{tolParConvNorm},
\code{tolParInv}] \code{convergence\$maxIters},
\code{convergence\$likelihoodTol}, \code{convergence\$parameterTol}, and
\code{execution\$inverseTol}, respectively.
\item[\code{REML}, \code{solveOnly}, \code{pqlControl}] \code{task\$objective}
(\code{"reml"} or \code{"ml"}), \code{task\$mode}
(\code{"fit"} or \code{"solve"}), and \code{task\$pql}.
\item[\code{computeCi}, \code{getPEV}, \code{returnParam}] \code{output\$inverse}
and \code{output\$setup}. PEV output is controlled by \code{output\$inverse};
the former \code{getPEV} flag has no separate replacement.
\item[\code{blend}, \code{blendWeight}, \code{stepWeight}] \code{update\$recovery},
\code{update\$blendWeight}, and \code{proposal\$stepWeight}.
\item[\code{acceleration="aitken"}] \code{proposal\$extrapolation="aitken"}.
\item[\code{acceleration="hybrid-ai"}] \code{curvature\$method="hybrid-ai"}.
Do not combine with Aitken.
\item[\code{compression}, \code{factorScoreAugmentation},
\code{covarianceInformation}, \code{retainCoefficientMatrix}] 
\code{representation\$compression}, \code{representation\$factorScores},
\code{output\$covariance}, and \code{output\$coefficientMatrix}.
\item[Former \code{pcg} settings] Map solve tolerances, iteration limits,
preconditioning, and reuse to \code{solve}; determinant scaling and Lanczos
depth to \code{logdet}; trace probes, estimator, rank, and covariance checks
to \code{score}; stochastic accuracy budgets to \code{convergence}; seed to
\code{execution}.
\item[Former \code{experimental} settings] Map structured representation
requests to \code{representation}, elimination requests to \code{solve},
trace batching to \code{score}, reduced AI to \code{curvature},
common-scale profiling to \code{update}, selected/full inverse requests to
\code{output}, and profiling to \code{diagnostics}.
\item[\code{denseMemoryMB}, \code{memoryPolicy}, \code{gpu},
\code{gpuMinSize}, \code{blasThreads}] \code{execution} fields with the same
names.
\item[\code{verbose}, \code{dateWarning}] \code{diagnostics} fields with the
same names.
\item[\code{recordDiagnostics}, \code{verboseDiagnostics}] Removed flags with
no direct replacement. Use \code{diagnostics\$verbose} for fit messages and
\code{diagnostics\$profile} for implemented phase profiling.

\end{description}

\end{Details}
%
\begin{Value}
An object of class \code{"sommer\_mmes\_control"}.
\end{Value}
%
\begin{Section}{Compatibility}
Settings are checked against one compatibility table when the control is
built and again in \code{\LinkA{mmes}{mmes}}. A supplied setting that conflicts with
another, or that would have no effect on the selected path, stops the call with
a message that names every offending setting and how to fix it. A setting counts
as supplied only when it differs from its default, so restating defaults is
always accepted. The main rules are:
\begin{itemize}

\item{} \code{formulation="direct"} accepts only \code{task}, \code{update\$recovery},
\code{update\$blendWeight},
\code{proposal\$stepWeight}, convergence iterations and tolerances,
\code{output\$setup}, \code{output\$inverse} 0 or 2, memory
and threading \code{execution} settings, and \code{diagnostics}.
\item{} \code{task\$mode="solve"} accepts only \code{solve\$tol}, \code{solve\$maxIters},
\code{solve\$warmStart}, \code{solve\$residualRefresh}, a non-Nystrom
\code{solve\$preconditioner}, \code{output\$setup}, CPU \code{execution} settings and
\code{diagnostics}.
\item{} PCG-only settings (solve tolerances, preconditioner, recycling,
\code{logdet}, \code{score} estimators, adaptive accuracy,
\code{representation\$matrixFree}, \code{execution\$seed}) require
\code{solve\$method="pcg"}; \code{"auto"} never selects PCG.
\item{} \code{solve\$method="pcg"} rejects EM/PXEM recovery, blending,
common-scale profiling, reduced AI, Aitken, hybrid AI, saddle-point
equations, factor scores and \code{task\$objective="ml"}.
\item{} Aitken extrapolation, hybrid AI, saddle-point equations, factor scores,
blending and \code{output\$coefficientMatrix=FALSE} each reject the modules
they cannot combine with.

\end{itemize}

Rules that depend on the data, family, \code{vcc} or rotation are checked by
\code{\LinkA{mmes}{mmes}} at fit time; for example a Gaussian identity fit stops when
\code{task\$pql} is changed.
\end{Section}
%
\begin{SeeAlso}\relax
\code{\LinkA{mmes}{mmes}}; \code{vignette("sommer.mmes.options", package="sommer")}
\end{SeeAlso}
%
\begin{Examples}
\begin{ExampleCode}
ctrl <- mmes_control(
  formulation = "henderson",
  solve = list(method = "pcg", tol = 1e-8),
  score = list(traceProbes = 16L),
  convergence = list(accuracy = "adaptive")
)
print(ctrl)
\end{ExampleCode}
\end{Examples}
\inputencoding{utf8}
\HeaderA{mme\_workload\_grid}{Generate a Factorial Grid of MME Workloads}{mme.Rul.workload.Rul.grid}
%
\begin{Description}\relax
Generate a Factorial Grid of MME Workloads
\end{Description}
%
\begin{Usage}
\begin{verbatim}
mme_workload_grid(
  structures = c("ar1", "us", "diag"),
  n_ind = c(40L, 100L),
  n_levels = c(4L, 10L),
  relationships = c("none", "dense"),
  reml = TRUE,
  computeCi = 0L
)
\end{verbatim}
\end{Usage}
%
\begin{Arguments}
\begin{ldescription}
\item[\code{structures}] Vector of covariance structures to test.

\item[\code{n\_ind}] Vector of individual counts.

\item[\code{n\_levels}] Vector of factor level counts.

\item[\code{relationships}] Vector of relationship types.

\item[\code{reml}] Logical; whether to fit by REML.

\item[\code{computeCi}] Vector of computeCi modes.
\end{ldescription}
\end{Arguments}
%
\begin{Value}
A named list of \code{mme\_workload} objects.
\end{Value}
\inputencoding{utf8}
\HeaderA{MNR}{Multivariate Newton-Raphson algorithm }{MNR}
%
\begin{Description}\relax
Multivariate Newton-Raphson algorithm used behind mmes when the henderson argument 
is set to FALSE. Algorithm made available for users that want to avoid the user-friendly interface 
and provide their matrices directly.
\end{Description}
%
\begin{Usage}
\begin{verbatim}
MNR(Y, # response
    X, Gx, # fixed effects
    Z, K, # random effects
    R,  # residual effects
    Ge, GeI, # initial values and constraints
    W, isInvW, # weights matrix
    iters, tolpar, tolparinv, # other params
    ai, pev,
    verbose, retscaled,
    stepweight, emweight
    )
\end{verbatim}
\end{Usage}
%
\begin{Arguments}
\begin{ldescription}
\item[\code{Y}] 
A matrix with rows for records and columns for traits. Expected class is a 'matrix'.

\item[\code{X}] 
Design matrices for fixed effects. Expected class is a 'list' with as many design 
matrices as fixed effects. Matrices within the list should be of class 'matrix'

\item[\code{Gx}] 
Trait multiplier matrix for fixed effects. Expected class is a 'list' with as many matrices as 
fixed effects. Each matrix is a square matrix with as many rows and columns as
number of traits. Each matrix is of class 'matrix'

\item[\code{Z}] 
Design matrices for random effects. Expected class is a 'list' with as many design 
matrices as random effects. Each element in the list should be a matrix of class
'dgCMatrix'

\item[\code{K}] 
Covariance matrices for design matrices of random effects. Expected class is a 'list' with as many  
covariance matrices as random effects specified in Z. Each element in the list should be a matrix of class
'dgCMatrix'

\item[\code{R}] 
Residual matrices for residual effects. Expected class is a 'list' with as many  
residual matrices as residual effects. Each element in the list should be a square 
matrix of dimensions n x n, where n is the number of records. Each matrix should be of class
'matrix'

\item[\code{Ge}] 
Initial values for variance components. Expected class is a 'list' with as many  
variance component matrices as random effects specified in Z. Each element in the 
list should be a matrix of dimensions t x t, where t is the number of traits and of 
class 'matrix'. Initial values are any real values.

\item[\code{GeI}] 
Initial constraints for variance components. Expected class is a 'list' with as many  
variance component constraints matrices as random effects specified in Z. Each element in the 
list should be a matrix of dimensions t x t, where t is the number of traits and of 
class 'matrix'. Values expected are:

0: not to be estimated

1: estimated and constrained to be positive (i.e. variance component)

2: estimated and unconstrained (can be negative or positive, i.e. covariance component)

3: not to be estimated but fixed (value has to be provided in the Gti argument)

Please notice that lower triangular values of these matrices have to be equal to zero since
only the values in the upper triangular are estimated.


\item[\code{W}] 
Design matrix for weights. A matrix of class 'matrix' for weighting the records.

\item[\code{isInvW}] 
A value of class 'logical' to indicate if the W matrix provided is already an inverse or
not. This aims to speed up the computations.

\item[\code{iters}] 
Maximum number of iterations allowed in REML. Value is of class 'integer'.

\item[\code{tolpar}] 
Convergence criteria for the change in log-likelihood. Value is of class 'numeric'.

\item[\code{tolparinv}] 
Tolerance parameter for matrix inverse used when singularities are encountered in
the estimation procedure. Value is of class 'numeric'.

\item[\code{ai}] 
A value of class 'logical' to indicate if the Average Information algorithm
should be used instead. That is faster but much less stable.

\item[\code{pev}] 
A value of class 'logical' to indicate if the predicted error variance should be 
computed or not. If FALSE computations are speeded up for models with many effects
to be estimated.

\item[\code{verbose}] 
A value of class 'logical' to value to indicate if the program should return
the progress of the iterative algorithm.

\item[\code{retscaled}] 
A value of class 'logical' to indicate if we should avoid scaling the traits.

\item[\code{stepweight}] 
Vector of class 'numeric' with length equal to niters to specify the relative weight
given to the second derivative Newton update.

\item[\code{emweight}] 
Vector of class 'numeric' with length equal to niters to specify the relative weight
given to the EM update compared to the NR update.



\end{ldescription}
\end{Arguments}
%
\begin{Details}\relax

This is the Rcpp-coded Direct-Inversion LMM REML algorithm used behind the mmer function.

\end{Details}
%
\begin{Value}
If all parameters are correctly indicated the program will return a list with the following information:

\begin{ldescription}
\item[\code{Vi}] the inverse of the phenotypic variance matrix V\textasciicircum{}- = (ZGZ+R)\textasciicircum{}-1
\item[\code{P}] the projection matrix  Vi - [Vi*(X*Vi*X)\textasciicircum{}-*Vi] 
\item[\code{sigma}] a list with the values of the variance-covariance components with one list element for each random effect.
\item[\code{sigma\_scaled}] a list with the values of the scaled variance-covariance components with one list element for each random effect.
\item[\code{sigmaSE}] Hessian matrix containing the variance-covariance for the variance components. SE's can be obtained taking the square root of the diagonal values of the Hessian.
\item[\code{Beta}] a data frame for trait BLUEs (fixed effects).
\item[\code{VarBeta}] a variance-covariance matrix for trait BLUEs
\item[\code{U}] a list (one element for each random effect) with a data frame for trait BLUPs.
\item[\code{VarU}] a list (one element for each random effect) with the variance-covariance matrix for trait BLUPs.
\item[\code{PevU}] a list (one element for each random effect) with the predicted error variance matrix for trait BLUPs.
\item[\code{fitted}]  Fitted values y.hat=XB
\item[\code{residuals}] Residual values e = Y - XB
\item[\code{AIC}] Akaike information criterion
\item[\code{BIC}] Bayesian information criterion
\item[\code{convergence}] a TRUE/FALSE statement indicating if the model converged.
\item[\code{monitor}] The values of log-likelihood and variance-covariance components across iterations during the REML estimation.
\item[\code{percChange}] The percent change of variance components across iterations. There should be one column less than the number of iterations. Calculated as percChange = ((x\_i/x\_i-1) - 1) * 100 where i is the ith iteration.
\item[\code{dL}] The vector of first derivatives of the likelihood with respect to the ith variance-covariance component.
\item[\code{dL2}] The matrix of second derivatives of the likelihood with respect to the i.j th variance-covariance component.

\end{ldescription}
\end{Value}
%
\begin{References}\relax

Mrode, R. A. (2014). Linear models for the prediction of animal breeding values. Cabi.

Covarrubias-Pazaran G (2016) Genome assisted prediction of quantitative traits using the R package sommer. PLoS ONE 11(6): doi:10.1371/journal.pone.0156744

\end{References}
%
\begin{SeeAlso}\relax
\code{\LinkA{mmes}{mmes}} -- the core function of the package 
\end{SeeAlso}
%
\begin{Examples}
\begin{ExampleCode}

data(DT_cpdata, package="enhancer")
DT <- DT_cpdata


# response matrix
y <- cbind(imputev(DT$Yield), imputev(DT$Firmness))
# fixed effect incidence matrix
X <- model.matrix(~Rowf, data=DT)
# random effect incidence matrix
Z <- Matrix::sparse.model.matrix(~id - 1, data=DT)
colnames(Z) <- gsub("id","",colnames(Z))
# covariance for random effect
GT <- GT_cpdata
A <- A.mat(GT) # additive relationship matrix
A <- A[colnames(Z),colnames(Z)] # make sure of the order
A <- A + diag(1e-4,nrow(A), nrow(A))
# residual effect incidence matrix (dimensions equal nrow(y))
R1 <- Matrix::Diagonal(n=nrow(y)) 
# weights matrix (dimensions equal nrow(y))
W <- diag(nrow(y)) 
# some other parameters
maxIter=3
stepWeight <- rep(0.9, maxIter)
stepWeight[1:2] <- c(0.5, 0.7)
emWeights <- rep(0,maxIter)
# model fit
res <- MNR( Y=y, # multi-trait response
            X=list(X), Gx=list(diag(2)), # fixed effects
            Z=list(Z), K=list(A), # random effects
            R=list(R1), # residual effects
            Ge=list( (diag(2)*.3)+.15 , diag(2)*0.75 ), # inital vc 
            GeI=list(unsm2(2),diag(2)), # vc constraints
            W=W, isInvW=TRUE, # weights for records
            iters=maxIter,
            tolpar=1e-4, tolparinv=1e-6, # tolerance
            ai=FALSE, pev=FALSE, # algorithm specifics
            verbose=TRUE, 
            retscaled=FALSE, 
            stepweight=stepWeight, # second derivatives weights
            emweight=emWeights # em update weights
)
res$sigma
res$Beta
res$U[[1]]


\end{ExampleCode}
\end{Examples}
\inputencoding{utf8}
\HeaderA{ownm}{User-defined covariance structure}{ownm}
%
\begin{Description}\relax
User-defined covariance structure. The returned dimensionless covariance shape is intended for use inside \code{vsm}, which supplies the single overall variance scale.
\end{Description}
%
\begin{Usage}
\begin{verbatim}
ownm(x, K = NULL, fun = NULL, par = numeric(), fixed = NULL,
  dfun = NULL, par_names = NULL, native_report = NULL)
\end{verbatim}
\end{Usage}
%
\begin{Arguments}
\begin{ldescription}
\item[\code{x}] Variable or design defining q covariance levels.
\item[\code{K}] Optional known positive-definite q by q covariance matrix.  If supplied, the structure is fixed.
\item[\code{fun}] Optional function fun(par) returning a finite q by q covariance matrix.
\item[\code{par}] Starting parameter vector for fun.
\item[\code{fixed}] Logical vector of length par indicating fixed parameters.
\item[\code{dfun}] Optional derivative function dfun(par, k) returning the derivative of the raw covariance matrix with respect to parameter k.
\item[\code{par\_names}] Optional names for the parameters.
\item[\code{native\_report}] Optional function with arguments \code{scale},
\code{par}, \code{factor}, and \code{absorb\_scale}. It must return a named
numeric vector giving the fitted parameters in the model's preferred native
scale. This callback is used by \code{covparams\_mmes()}.
\end{ldescription}
\end{Arguments}
%
\begin{Details}\relax
There are two modes.  With \code{K}, \code{ownm} creates a fixed covariance factor.  With \code{fun}, the matrix is evaluated at every optimizer point.  The returned matrix is symmetrized and normalized by its first diagonal element so that \code{vsm} retains the unique overall variance scale.  If \code{dfun} is supplied, its derivative is normalized analytically using the quotient rule; otherwise a central numerical derivative of the normalized factor is used.  This is the general extension mechanism for covariance structures that do not require a native C++ evaluator. When \code{native\_report} is omitted, the overall variance and naturally transformed \code{par} values are reported.
\end{Details}
%
\begin{Value}
A list containing the incidence/design matrix in \code{Z} and a compiled CovarianceFactor v2 descriptor in \code{covFactor}, for use inside \code{vsm}.
\end{Value}
%
\begin{SeeAlso}\relax
\code{\LinkA{vsm}{vsm}}, \code{\LinkA{covparams\_mmes}{covparams.Rul.mmes}}, \code{\LinkA{rrm}{rrm}}, \code{\LinkA{maternm}{maternm}}, \code{\LinkA{sar}{sar}}, \code{\LinkA{car}{car}}.
\end{SeeAlso}
%
\begin{Examples}
\begin{ExampleCode}
## Not run: 
K <- matrix(c(1,.3,.3,1),2,2)
vsm(ownm(group, K=K), ism(id))

cfun <- function(p) {
  r <- tanh(p[1]); matrix(c(1,r,r,1),2,2)
}
native <- function(scale, par, factor, absorb_scale=TRUE) {
  c(variance=scale, correlation=tanh(par[1]))
}
vsm(ownm(group, fun=cfun, par=0, native_report=native), ism(id))

## End(Not run)
\end{ExampleCode}
\end{Examples}
\inputencoding{utf8}
\HeaderA{plot.mmes}{plot form a LMM plot with mmes}{plot.mmes}
\aliasA{plot.mmer}{plot.mmes}{plot.mmer}
\keyword{models}{plot.mmes}
%
\begin{Description}\relax
\code{plot} method for class \code{"mmes"}.
\end{Description}
%
\begin{Usage}
\begin{verbatim}
## S3 method for class 'mmes'
plot(x,stnd=TRUE, ...)
\end{verbatim}
\end{Usage}
%
\begin{Arguments}
\begin{ldescription}
\item[\code{x}] an object of class \code{"mmes"}
\item[\code{stnd}] argument for ploting the residuals to know if they should be standarized.
\item[\code{...}] Further arguments to be passed
\end{ldescription}
\end{Arguments}
%
\begin{Value}
vector of plot

\end{Value}
%
\begin{Author}\relax
Giovanny Covarrubias \email{covarrubiasp@wisc.edu}
\end{Author}
%
\begin{SeeAlso}\relax
\code{\LinkA{plot}{plot}}, \code{\LinkA{mmes}{mmes}}
\end{SeeAlso}
%
\begin{Examples}
\begin{ExampleCode}
data(DT_yatesoats, package="enhancer")
DT <- DT_yatesoats
head(DT)
m3 <- mmes(fixed=Y ~ V + N + V:N,
           random = ~ B + B:MP,
           rcov=~units,
           data = DT)

\end{ExampleCode}
\end{Examples}
\inputencoding{utf8}
\HeaderA{pmonitor}{plot the change of VC across iterations}{pmonitor}
\keyword{models}{pmonitor}
%
\begin{Description}\relax
\code{plot} for monitoring.
\end{Description}
%
\begin{Usage}
\begin{verbatim}
pmonitor(object, ...)
\end{verbatim}
\end{Usage}
%
\begin{Arguments}
\begin{ldescription}
\item[\code{object}] model object of class \code{"mmes"}
\item[\code{...}] Further arguments to be passed to the plot function.
\end{ldescription}
\end{Arguments}
%
\begin{Value}
vector of plot

\end{Value}
%
\begin{Author}\relax
Giovanny Covarrubias
\end{Author}
%
\begin{SeeAlso}\relax
\code{\LinkA{plot}{plot}}, \code{\LinkA{mmes}{mmes}}
\end{SeeAlso}
%
\begin{Examples}
\begin{ExampleCode}
data(DT_yatesoats, package="enhancer")
DT <- DT_yatesoats
head(DT)
m3 <- mmes(fixed=Y ~ V + N + V:N,
           random = ~ B + B:MP,
           rcov=~units,
           data = DT)
pmonitor(m3)
\end{ExampleCode}
\end{Examples}
\inputencoding{utf8}
\HeaderA{postPEV}{Post-fit prediction error variances and inverse coefficient matrix}{postPEV}
%
\begin{Description}\relax
Computes prediction error variances (PEVs) or the complete inverse of the mixed model coefficient matrix after a model has been fitted with \code{mmes}. For PEV-only calculations, the function uses the Takahashi sparse inverse subset approach, avoiding computation of the complete inverse of the coefficient matrix.
\end{Description}
%
\begin{Usage}
\begin{verbatim}
postPEV(object, mode = 1L)
\end{verbatim}
\end{Usage}
%
\begin{Arguments}
\begin{ldescription}
\item[\code{object}] 
A fitted model object of class \code{"mmes"}. The object must contain the final sparse mixed model coefficient matrix and its associated scaling factor.

\item[\code{mode}] 
Integer specifying the post-fit inverse calculation to perform. Use \code{0} to perform no inverse calculation and clear previously computed inverse-derived outputs, \code{1} to compute the prediction error variances using the Takahashi sparse inverse subset approach without forming the complete inverse, or \code{2} to compute the complete inverse of the coefficient matrix and the prediction error variances. The default is \code{1}.

\end{ldescription}
\end{Arguments}
%
\begin{Details}\relax
The function is intended for post-processing models fitted with \code{mmes}, allowing prediction error variances or the complete inverse of the mixed model coefficient matrix to be calculated only when they are needed.

When \code{mode = 1}, the sparse mixed model coefficient matrix stored in the fitted object is factorized and the Takahashi equations are used to obtain the selected elements of its inverse required for the diagonal prediction error variances. The complete inverse is not formed.

When \code{mode = 2}, the complete inverse of the mixed model coefficient matrix is computed. The prediction error variances are then obtained from the corresponding diagonal elements of this inverse.

When \code{mode = 0}, no matrix factorization or inverse calculation is performed and previously stored inverse-derived results are cleared.

\code{\LinkA{predict.mmes}{predict.mmes}} does not require any particular \code{mode}: it computes exact standard errors for arbitrary linear combinations of fixed and random effects directly from the stored coefficient matrix \code{C}, without using \code{Ci} at all (see \code{\LinkA{predict.mmes}{predict.mmes}} for details). \code{postPEV()} is only needed when the diagonal prediction error variances (\code{uPevList}) or the complete inverse (\code{Ci}) are wanted for their own sake.

For the sampling variance of the BLUP itself, use \code{\LinkA{postVarU}{postVarU}}:
\eqn{\operatorname{Var}(\hat u)=G-\operatorname{Var}(u-\hat u)}{}.
PEV and BLUP sampling variance are different quantities. Neither includes
uncertainty from estimating covariance parameters. \code{postVarU()} leaves
existing PEV outputs unchanged.

For this post-fit calculation to be available, fit the \code{mmes} model with
\code{control = mmes\_control(output = list(coefficientMatrix = TRUE))} so the
final coefficient matrix is retained. Models can use
\code{control\$output\$inverse = 0} during fitting, then calculate PEVs or the
complete inverse afterwards only if required.
\end{Details}
%
\begin{Value}

An object of class \code{"mmes"} containing the original fitted model together with the requested post-fit inverse information.

With \code{mode = 1}, \code{uPevList} contains random-effect prediction error variances. \code{Ci} may contain the selected inverse entries used to obtain those PEVs, but it does not contain the complete inverse.

With \code{mode = 2}, \code{Ci} contains the complete inverse of the mixed model coefficient matrix and \code{uPevList} contains the corresponding prediction error variances for the random effects.

With \code{mode = 0}, inverse-derived outputs are cleared.

The returned object also contains \code{CiComputed} and \code{CiMode}. \code{CiComputed} is TRUE only when the complete inverse was computed; \code{CiMode} records the selected mode.

\end{Value}
%
\begin{References}\relax

Takahashi, K., Fagan, J., and Chin, M. S. (1973). Formation of a sparse bus impedance matrix and its application to short circuit study. Proceedings of the 8th PICA Conference, Minneapolis, Minnesota.

Covarrubias-Pazaran G (2016) Genome assisted prediction of quantitative traits using the R package sommer. PLoS ONE 11(6): doi:10.1371/journal.pone.0156744

\end{References}
%
\begin{Examples}
\begin{ExampleCode}
####=========================================####
#### Fit a mixed model without computing PEVs
#### or the complete inverse during estimation
####=========================================####

# mod <- mmes(
#   fixed = Yield ~ 1,
#   random = ~ vsm(ism(Genotype)),
#   rcov = ~ units,
#   data = DT,
#   control = mmes_control(output = list(inverse = 0))
# )

####=========================================####
#### Compute PEVs afterwards using the
#### Takahashi sparse inverse subset
####=========================================####

# mod <- postPEV(mod, mode = 1)
# mod$uPevList

####=========================================####
#### Compute the complete inverse afterwards
####=========================================####

# mod <- postPEV(mod, mode = 2)
# mod$Ci
# mod$uPevList
\end{ExampleCode}
\end{Examples}
\inputencoding{utf8}
\HeaderA{postVarU}{Post-fit Sampling Covariance of BLUPs}{postVarU}
%
\begin{Description}\relax
Computes the plug-in sampling variance or covariance of fitted random-effect
predictors, \eqn{\hat u}{}, using Henderson's mixed-model coefficient system.
This is distinct from the prior variance of the random effects and their
prediction error variance (PEV).
\end{Description}
%
\begin{Usage}
\begin{verbatim}
postVarU(object, mode=1L)
\end{verbatim}
\end{Usage}
%
\begin{Arguments}
\begin{ldescription}
\item[\code{object}] A Gaussian \code{mmes} Henderson fit with stored \code{C} and
\code{Cscale}. Direct-engine, PQL and matrix-free \code{solveOnly} fits are
not supported.
\item[\code{mode}] \code{1} (default) returns diagonal BLUP sampling variances in
\code{uVarList}; \code{2} additionally returns the joint sampling covariance
matrix \code{VarU}; \code{0} clears only VarU-derived outputs.
\end{ldescription}
\end{Arguments}
%
\begin{Details}\relax
For \eqn{y=X\beta+Zu+e}{}, with \eqn{\operatorname{Var}(u)=G}{} and
\eqn{\operatorname{Var}(e)=R}{}, the unscaled Henderson matrix is
\deqn{C_0=\left[\begin{array}{cc}
X'R^{-1}X & X'R^{-1}Z \\
Z'R^{-1}X & Z'R^{-1}Z+G^{-1}
\end{array}\right].}{}
Its random-effect inverse block is
\eqn{Q=(C_0^{-1})_{uu}=\operatorname{Var}(u-\hat u)}{}. Consequently,
\deqn{\operatorname{Var}(\hat u)=G-Q.}{}
For known covariance parameters this identity is exact. At estimated
parameters the result is a plug-in covariance, omitting covariance-parameter
estimation uncertainty. The same identity is used for ML and REML fits.

The implementation does not construct or invert the observation covariance
\eqn{V=ZGZ'+R}{}. Error covariance is computed by solving the stored Henderson
system and multiplying by \code{Cscale} to recover the original response
scale. For each independent random term,
\eqn{G_k=\Sigma_k\otimes A_k}{}, where \eqn{\Sigma_k}{} is the fitted
covariance-coordinate matrix and \eqn{A_k}{} the relationship covariance.
Products with \eqn{A_k}{} are evaluated by solving its retained sparse
precision, not by explicitly forming its inverse or the full joint \eqn{G}{}.

\code{mode=1} processes identity contrasts in batches of at most 128 rows,
so neither a full coefficient inverse nor a full random-effect covariance
matrix is stored. It requires solves for all requested effect diagonals and
can be expensive for large models. \code{mode=2} requires quadratic storage
in the number of random effects. Although independent terms have zero prior
cross-covariance, their BLUPs can have nonzero sampling cross-covariance; the
joint \code{VarU} includes these cross-term entries.

Rotation and factor-score augmentation are mapped to the public
\code{uList} coordinates before computing error covariance. New fits retain
\code{randomPrecision} in original level order, so external \code{Gu}
objects are not needed after fitting. Older ordinary fits attempt to rebuild
relationship precisions from saved model inputs, requiring the original
formula objects; dimension checks cannot detect all changes to those objects.
Older rotated or augmented fits must be refitted. No clipping is performed:
small negative variances from numerical cancellation remain visible.

Existing \code{Ci} and \code{uPevList} are left unchanged. The variance here
is not the empirical variance across BLUP values and is not a posterior
variance of \eqn{u}{}. It is the repeated-sampling covariance of \eqn{\hat u}{}.
\end{Details}
%
\begin{Value}
The original fitted object with \code{uVarList}, a named list matching the
dimensions and labels of \code{uList}, and \code{uVarMode}. With
\code{mode=2}, \code{VarU} is the joint covariance in term order, with each
\code{uList} matrix stacked column by column. Its labels contain term,
covariance-coordinate and main-effect labels. Modes 0 and 1 remove any
previous full \code{VarU}; mode 0 also removes \code{uVarList}.
\end{Value}
%
\begin{SeeAlso}\relax
\code{\LinkA{postPEV}{postPEV}}, \code{\LinkA{predict.mmes}{predict.mmes}},
\code{\LinkA{covmatrix\_mmes}{covmatrix.Rul.mmes}}
\end{SeeAlso}
%
\begin{Examples}
\begin{ExampleCode}
data(DT_example, package="enhancer")
model <- mmes(Yield~Env,
  random=~Name,
  data=DT_example,
  control = mmes_control(output = list(inverse = 0),
  diagnostics = list(verbose = FALSE)))
model <- postVarU(model)
head(model$uVarList[[1]])
full <- postVarU(model, mode=2)
head(diag(full$VarU))

prediction <- predict(model, D="Name", PEV=TRUE, VarU=TRUE)
head(prediction$pvals)
prediction$PEV[1:3, 1:3]
prediction$VarU[1:3, 1:3]
prediction$sampling.vcov[1:3, 1:3]
\end{ExampleCode}
\end{Examples}
\inputencoding{utf8}
\HeaderA{predict.mmes}{Predict form of a LMM fitted with mmes}{predict.mmes}
\aliasA{predict.mmer}{predict.mmes}{predict.mmer}
\keyword{models}{predict.mmes}
%
\begin{Description}\relax
\code{predict} method for class \code{"mmes"}.
\end{Description}
%
\begin{Usage}
\begin{verbatim}
## S3 method for class 'mmes'
predict(object, Dtable=NULL, D, levels=NULL, sed=FALSE,
        pairwise=FALSE, adjust="none", df=Inf, PEV=TRUE, VarU=TRUE, ...)
\end{verbatim}
\end{Usage}
%
\begin{Arguments}
\begin{ldescription}
\item[\code{object}] a mixed model of class \code{"mmes"}
\item[\code{Dtable}] a table specifying the terms to be included or averaged.

An "include" term uses the cells matching each prediction row. For
categorical random effects these are mapped within the fitted term, including
fitted cells without records. Numeric fixed terms retain their evaluated
covariate values rather than being replaced by indicators.

An "include and average" term means that the included cells of that effect are averaged: each prediction row is divided by the number of cells of the term present with it (for fixed terms this count includes the reference cells absorbed in the intercept).

An "average" term alone averages over the selected factor levels. For
categorical terms the weights are 1/\#levels; numeric fixed terms are
evaluated at their covariate settings before averaging.

If a term is not considered "include" or "average" is then totally ignored in the BLUP and SE calculation. 

The default rule to invoke when the user doesn't provide the Dtable is to include and average all terms that match the argument D.

The \code{levels} column is a list-column. \code{NULL} means all factor
levels and default numeric settings. Entries restrict levels or specify
covariate settings; see the \code{levels} argument. The returned Dtable also
records resolved settings for numeric terms and can be reused.


\item[\code{D}] a character string specifying the variable used to extract levels for the rows of the D matrix and its construction. Alternatively, the D matrix (of class dgCMatrix) specifying the matrix to be used for the predictions directly.
\item[\code{levels}] an optional named list whose names are values of \code{Dtable\$term}.
Each element is stored in the \code{levels} column of the returned Dtable,
so the Dtable remains the complete description of how D was built. Levels
are written as data levels for fixed factors (\code{"a:b"} for
interactions, e.g. \code{"Victory:0.2"}), as cell labels for categorical random
terms (e.g. \code{"CA.2011:genotype1"} for an Env by Name term), and as numeric
settings for covariate terms. A term containing one numeric variable accepts
a numeric scalar, e.g. \code{levels=list("Env:cov"=5)}. Numeric or mixed
terms also accept a named per-variable list, e.g.
\code{levels=list("Env:cov"=list(Env=c("CA.2011", "CA.2012"), cov=5))}.
Each numeric setting must be one finite value. Use the exact term spelling
in \code{object\$Dtable\$term}. For an included categorical term matching the classify
variable(s) the levels select the prediction rows, in the given order, and
may include levels without records (e.g. levels only present in \code{Gu}).
For other included or averaged terms, only the listed levels enter the
inclusion or the average. Restrictions are term-specific: selecting levels
of \code{V} does not automatically restrict \code{V:N}.
\item[\code{sed}] if \code{TRUE}, return the matrix of standard errors of differences
\code{sed} and its summary \code{avsed} (mean, minimum and maximum).
\item[\code{pairwise}] \code{TRUE} for all pairwise differences between predictions,
or a single prediction level to compare every other prediction against it.
\item[\code{adjust}] multiplicity adjustment passed to \code{\LinkA{p.adjust}{p.adjust}}.
\item[\code{df}] degrees of freedom for pairwise tests: \code{Inf} (normal-based,
default), a positive number, \code{"residual"} for n - rank(X), or
\code{"satterthwaite"}/\code{"kr"} (see \code{\LinkA{wald\_mmes}{wald.Rul.mmes}}). The latter
two apply to differences involving only fixed effects; differences involving
random effects keep normal-based tests (\code{df=Inf}). With \code{"kr"} the
SED of those differences uses the Kenward-Roger adjusted covariance.
\item[\code{PEV}] Logical, default \code{TRUE}. Return the full target prediction-error
covariance as \code{PEV} and its diagonal as \code{pvals\$pev}. When
\code{FALSE}, only these additional outputs are omitted; \code{vcov} and
\code{std.error} remain available for compatibility.
\item[\code{VarU}] Logical, default \code{TRUE}. Return the sampling covariance of
the random contribution as \code{VarU}, its diagonal as \code{pvals\$var.u},
and the sampling covariance of the entire predictor as \code{sampling.vcov}.
Set \code{FALSE} to skip the additional relationship solves and outputs.
These calculations require a Gaussian Henderson fit. For PQL predictions
use \code{VarU=FALSE}; their existing working-model uncertainty remains an
approximation.
\item[\code{...}] Further arguments to be passed.
\end{ldescription}
\end{Arguments}
%
\begin{Details}\relax
This function allows to produce predictions specifying those variables that define the margins of the hypertable to be predicted (argument D). Predictions are obtained for each combination of values of the specified variables that is present in the data set used to fit the model. See vignettes for more details. 

For categorical structured random terms, the fitted covariance-coordinate
and main-effect levels define the genetic prediction weights, even when a
cell has no observations. The \code{levels} argument can additionally request
fitted levels present only in \code{Gu}. Unsupported random designs that
cannot be mapped for all prediction rows produce a warning; an explicit
\code{D} matrix can be used for these designs.

Numeric fixed terms, including factor-by-covariate interactions, are
evaluated using the fitted fixed formula and contrasts. Unless overridden,
each numeric variable is set to its mean in the retained fitting data.
Factor levels are equally weighted when averaged. Thus an averaged
\code{Env:cov} term with one slope per environment uses
\code{mean(cov)/nlevels(Env)} for each slope, not \code{1/nlevels(Env)}.

Write \eqn{D=[L_b,L_u]}{}. The target is \eqn{t=L_b\beta+L_u u}{} and its
predictor is \eqn{\hat t=L_b\hat\beta+L_u\hat u}{}. The existing
prediction \code{PEV} and \code{VarU} matrices are for the linear
combinations represented by the final returned \code{D}, after applying
\code{Dtable} and \code{levels}, not copies of model-level effect variances.
Each covariance matrix has one row and column per requested prediction.
The existing
\code{vcov}, optional \code{PEV}, and squared \code{std.error} refer to
\eqn{\operatorname{Var}(t-\hat t)}{}, not generally to
\eqn{\operatorname{Var}(\hat t)}{}. They are computed from
\deqn{\operatorname{Var}(t-\hat t)=D C_0^{-1}D',}{}
where \eqn{C_0}{} is the original-scale Henderson coefficient matrix defined
in \code{\LinkA{postVarU}{postVarU}}. The stored system \code{C} is solved against
contrast transposes and the result multiplied by \code{Cscale}. Arbitrary
contrasts work regardless of \code{computeCi}; no stored \code{Ci} is used.

With \eqn{Q=(C_0^{-1})_{uu}}{} and \eqn{B=(C_0^{-1})_{bb}}{}, the new outputs are
\deqn{\mathrm{VarU}=L_u(G-Q)L_u',}{}
\deqn{\mathrm{sampling.vcov}=L_b B L_b'+L_u(G-Q)L_u'.}{}
The BLUE and BLUP have zero sampling cross-covariance, whereas the
fixed/random blocks of the inverse Henderson matrix describe error
covariance and generally are not zero. In particular, subtracting the full
target PEV from \eqn{L_uGL_u'}{} is incorrect if fixed effects are included.
A fixed-only prediction has zero \code{VarU}, but generally nonzero
\code{sampling.vcov} and PEV. Random-only predictions satisfy
\eqn{L_uGL_u'=\mathrm{PEV}+\mathrm{VarU}}{}.

The implementation batches full, random-only and fixed-only contrasts into
one Henderson factorization. Prior products use the term structure
\eqn{G_k=\Sigma_k\otimes A_k}{} and sparse relationship-precision solves.
Neither the observation covariance \eqn{V}{} nor its inverse is constructed.
Additional matrices require quadratic storage in the number of predictions.
Rotation and factor-score mappings operate in public effect coordinates.
New fits retain relationship precisions; older ordinary fits attempt model
input reconstruction as described in \code{\LinkA{postVarU}{postVarU}}.

All covariance outputs use the fitted variance components as fixed,
omitting their estimation uncertainty. They are exact for known components
and plug-in approximations for estimated components, for both ML and REML.
SEDs, pairwise tests and \code{std.error} retain their PEV-based meaning;
enabling VarU does not change those calculations. Direct and matrix-free
\code{solveOnly} fits lack the stored coefficient system required here.


\end{Details}
%
\begin{Value}

\begin{ldescription}
\item[\code{pvals}] the table of predictions according to the specified arguments.

\item[\code{vcov}] The full target prediction-error covariance, always returned.
\item[\code{PEV}] When requested, identical to \code{vcov}; its diagonal is
\code{pvals\$pev}.
\item[\code{VarU}] When requested, sampling covariance of the predicted random
contribution only; its diagonal is \code{pvals\$var.u}.
\item[\code{sampling.vcov}] When \code{VarU=TRUE}, sampling covariance of the entire
fixed-plus-random predictor. This is not the covariance used for standard errors.

\item[\code{D}] the model matrix for predictions as defined in Welham et al.(2004).

\item[\code{Dtable}] the table specifying the terms to include and terms to be averaged, including the \code{levels} used.

\item[\code{sed, avsed}] when \code{sed=TRUE}: the standard errors of differences, sqrt(v\_ii + v\_jj - 2 v\_ij) from \code{vcov}, and their mean, minimum and maximum.

\item[\code{pairwise}] when requested: a data frame with \code{level1}, \code{level2}, \code{difference} (level1 - level2), \code{SED}, \code{statistic}, \code{df}, \code{p.value} and \code{p.adjusted}.

\end{ldescription}
\end{Value}
%
\begin{Author}\relax
Giovanny Covarrubias-Pazaran
\end{Author}
%
\begin{References}\relax
Welham, S., Cullis, B., Gogel, B., Gilmour, A., and Thompson, R. (2004). Prediction in linear mixed models. Australian and New Zealand Journal of Statistics, 46, 325 - 347.

\end{References}
%
\begin{SeeAlso}\relax
\code{\LinkA{predict}{predict}}, \code{\LinkA{mmes}{mmes}}
\end{SeeAlso}
%
\begin{Examples}
\begin{ExampleCode}

data(DT_yatesoats, package="enhancer")
DT <- DT_yatesoats
m3 <- mmes(fixed=Y ~ V + N + V:N ,
           random = ~ B + B:MP,
           rcov=~units,
           data = DT)

m3 <- postPEV(m3, mode = 2)
 
#############################
## predict means for nitrogen
#############################
Dt <- m3$Dtable; Dt
# first fixed effect just average
Dt[1,"average"] = TRUE
# second fixed effect include
Dt[2,"include"] = TRUE
# third fixed effect include and average
Dt[3,"include"] = TRUE
Dt[3,"average"] = TRUE
Dt

pp=predict(object=m3, Dtable=Dt, D="N")
pp$pvals

#############################
## predict means for variety
#############################

Dt <- m3$Dtable; Dt
# first fixed effect include
Dt[1,"include"] = TRUE
# second fixed effect just average
Dt[2,"average"] = TRUE
# third fixed effect include and average
Dt[3,"include"] = TRUE
Dt[3,"average"] = TRUE
Dt

pp=predict(object=m3, Dtable=Dt, D="V")
pp$pvals

#############################
## predict means for nitrogen:variety
#############################
# prediction matrix D based on (equivalent to classify in asreml)
Dt <- m3$Dtable; Dt
# first fixed effect include and average
Dt[1,"include"] = TRUE
Dt[1,"average"] = TRUE
# second fixed effect include and average
Dt[2,"include"] = TRUE
Dt[2,"average"] = TRUE
# third fixed effect include and average
Dt[3,"include"] = TRUE
Dt[3,"average"] = TRUE
Dt

pp=predict(object=m3, Dtable=Dt, D="N:V")
pp$pvals

#############################
## select levels through the argument or the Dtable
#############################
Dt <- m3$Dtable
Dt$average[Dt$term %in% c("1", "V", "V:N")] <- TRUE
Dt$include[Dt$term %in% c("N", "V:N")] <- TRUE

p <- predict(m3, D="N", Dtable=Dt,
             levels=list(N=c("0.6", "0.2")))
p$pvals
p$Dtable$levels[[match("N", p$Dtable$term)]]

Dt$levels[[match("N", Dt$term)]] <- c("0.6", "0.2")
p2 <- predict(m3, D="N", Dtable=Dt)
all.equal(p$D, p2$D)

## Covariances of the requested linear combinations
p <- predict(m3, D="N", Dtable=Dt, PEV=TRUE, VarU=TRUE)
p$PEV
p$VarU
p$sampling.vcov
randomColumns <- seq.int(nrow(m3$b)+1L, nrow(m3$bu))
randomD <- p$D[,randomColumns,drop=FALSE]
full <- postVarU(m3, mode=2)
all.equal(unname(p$VarU),
          unname(as.matrix(randomD %*% full$VarU %*% t(randomD))))

\end{ExampleCode}
\end{Examples}
\inputencoding{utf8}
\HeaderA{print.sommer\_mmes\_control}{Print Modular Numerical Controls}{print.sommer.Rul.mmes.Rul.control}
%
\begin{Description}\relax
Prints the formulation and each named numerical-control module. Nested lists
are summarized rather than expanded; long vectors are represented by their
length. The method returns the control object invisibly.
\end{Description}
%
\begin{Usage}
\begin{verbatim}
## S3 method for class 'sommer_mmes_control'
print(x, ...)
\end{verbatim}
\end{Usage}
%
\begin{Arguments}
\begin{ldescription}
\item[\code{x}] An object of class \code{sommer\_mmes\_control}.

\item[\code{...}] Ignored.
\end{ldescription}
\end{Arguments}
\inputencoding{utf8}
\HeaderA{print.sommer\_mme\_benchmark}{Print method for sommer\_mme\_benchmark}{print.sommer.Rul.mme.Rul.benchmark}
%
\begin{Description}\relax
Print method for sommer\_mme\_benchmark
\end{Description}
%
\begin{Usage}
\begin{verbatim}
## S3 method for class 'sommer_mme_benchmark'
print(x, ...)
\end{verbatim}
\end{Usage}
%
\begin{Arguments}
\begin{ldescription}
\item[\code{x}] An object of class \code{sommer\_mme\_benchmark}.

\item[\code{...}] Ignored.
\end{ldescription}
\end{Arguments}
\inputencoding{utf8}
\HeaderA{print.sommer\_mme\_crossover}{Print method for sommer\_mme\_crossover}{print.sommer.Rul.mme.Rul.crossover}
%
\begin{Description}\relax
Print method for sommer\_mme\_crossover
\end{Description}
%
\begin{Usage}
\begin{verbatim}
## S3 method for class 'sommer_mme_crossover'
print(x, ...)
\end{verbatim}
\end{Usage}
%
\begin{Arguments}
\begin{ldescription}
\item[\code{x}] An object of class \code{sommer\_mme\_crossover}.

\item[\code{...}] Ignored.
\end{ldescription}
\end{Arguments}
\inputencoding{utf8}
\HeaderA{print.sommer\_mme\_profile}{Print method for sommer\_mme\_profile}{print.sommer.Rul.mme.Rul.profile}
%
\begin{Description}\relax
Print method for sommer\_mme\_profile
\end{Description}
%
\begin{Usage}
\begin{verbatim}
## S3 method for class 'sommer_mme_profile'
print(x, ...)
\end{verbatim}
\end{Usage}
%
\begin{Arguments}
\begin{ldescription}
\item[\code{x}] An object of class \code{sommer\_mme\_profile}.

\item[\code{...}] Ignored.
\end{ldescription}
\end{Arguments}
\inputencoding{utf8}
\HeaderA{profile\_mme\_workload}{Profile an MME Workload by Phase}{profile.Rul.mme.Rul.workload}
%
\begin{Description}\relax
Fits one workload with phase profiling enabled and tabulates where the time goes.
\end{Description}
%
\begin{Usage}
\begin{verbatim}
profile_mme_workload(
  workload,
  engine = mmes_control(),
  nIters = 5L,
  dominantShare = 0.3
)
\end{verbatim}
\end{Usage}
%
\begin{Arguments}
\begin{ldescription}
\item[\code{workload}] An \code{mme\_workload} from \code{\LinkA{generate\_mme\_workload}{generate.Rul.mme.Rul.workload}}.

\item[\code{engine}] A \code{\LinkA{mmes\_control}{mmes.Rul.control}} object; profiling is forced on.

\item[\code{nIters}] Maximum REML/ML iterations. Default is 5.

\item[\code{dominantShare}] Share of fit time above which a phase is flagged dominant. Default is 0.3.
\end{ldescription}
\end{Arguments}
%
\begin{Value}
A list of class \code{"sommer\_mme\_profile"} with a per-phase \code{table},
workload \code{counters}, per-iteration seconds, and the total elapsed time.
\end{Value}
\inputencoding{utf8}
\HeaderA{r2}{Reliability}{r2}
%
\begin{Description}\relax
Calculates the reliability of BLUPs in a sommer model.
\end{Description}
%
\begin{Usage}
\begin{verbatim}
r2(object, object2=NULL)
\end{verbatim}
\end{Usage}
%
\begin{Arguments}
\begin{ldescription}
\item[\code{object}] 
Model fitted with the mmes function.

\item[\code{object2}] 
An optional model identical to object in the first argument but fitted with the argument returnParam set to TRUE to access the relationship matrices from the fitted model.

\end{ldescription}
\end{Arguments}
%
\begin{Details}\relax

The reliability method calculated is the classical animal model: R2=(G-PEV)/G

\end{Details}
%
\begin{Value}
\begin{description}

\item[result] a list with as many elements as random effects fitted containing reliabilities for individual BLUPs.

\end{description}

\end{Value}
%
\begin{References}\relax

Mrode, R. A. (2014). Linear models for the prediction of animal breeding values. Cabi.

Covarrubias-Pazaran G (2016) Genome assisted prediction of quantitative traits using the R package sommer. PLoS ONE 11(6): doi:10.1371/journal.pone.0156744

\end{References}
%
\begin{SeeAlso}\relax
\code{\LinkA{mmes}{mmes}} -- the core function of the package 
\end{SeeAlso}
%
\begin{Examples}
\begin{ExampleCode}
####=========================================####
#### Example population
####=========================================####
data(DT_example, package="enhancer")
DT <- DT_example
head(DT)
ans1 <- mmes(Yield~Env, 
             random= ~ Name + Env:Name,
             rcov= ~ units,
             data=DT)
ans1 <- postPEV(ans1, mode = 1)
rel=r2(ans1)
\end{ExampleCode}
\end{Examples}
\inputencoding{utf8}
\HeaderA{randef}{extracting random effects}{randef}
%
\begin{Description}\relax
This function is extracts the random effects from a mixed model fitted by mmer.
\end{Description}
%
\begin{Usage}
\begin{verbatim}
randef(object)
\end{verbatim}
\end{Usage}
%
\begin{Arguments}
\begin{ldescription}
\item[\code{object}] an mmer object
\end{ldescription}
\end{Arguments}
%
\begin{Value}
\begin{description}

\item[\$randef]  a list structure with the random effects or BLUPs.

\end{description}

\end{Value}
%
\begin{Examples}
\begin{ExampleCode}
# randef(model)
\end{ExampleCode}
\end{Examples}
\inputencoding{utf8}
\HeaderA{residuals.mmes}{Residuals form a GLMM fitted with mmes}{residuals.mmes}
\aliasA{residuals.mmer}{residuals.mmes}{residuals.mmer}
\keyword{models}{residuals.mmes}
%
\begin{Description}\relax
\code{residuals} method for class \code{"mmes"}.
\end{Description}
%
\begin{Usage}
\begin{verbatim}
## S3 method for class 'mmes'
residuals(object,
                          type=c("response", "deviance", "pearson", "working"), ...)
\end{verbatim}
\end{Usage}
%
\begin{Arguments}
\begin{ldescription}
\item[\code{object}] an object of class \code{"mmes"}
\item[\code{type}] For a non-Gaussian PQL fit, return response residuals (default),
signed deviance residuals, Pearson residuals, or final IRLS working residuals.
Deviance and Pearson residuals use the prior weights (binomial trials, see
\code{\LinkA{binm}{binm}}).
\item[\code{...}] Further arguments to be passed
\end{ldescription}
\end{Arguments}
%
\begin{Value}
For Gaussian fits, residuals of the form e = y - Xb - Zu. For non-Gaussian
PQL fits, residuals on the requested scale.

\end{Value}
%
\begin{Author}\relax
Giovanny Covarrubias
\end{Author}
%
\begin{SeeAlso}\relax
\code{\LinkA{residuals}{residuals}}, \code{\LinkA{mmes}{mmes}}
\end{SeeAlso}
\inputencoding{utf8}
\HeaderA{rotate\_back\_mmes}{Back-transform rotated mmes random effects}{rotate.Rul.back.Rul.mmes}
%
\begin{Description}\relax
Back-transforms the selected relationship random effect from the solver
eigenbasis to the original relationship-level basis.
\end{Description}
%
\begin{Usage}
\begin{verbatim}
rotate_back_mmes(object)
\end{verbatim}
\end{Usage}
%
\begin{Arguments}
\begin{ldescription}
\item[\code{object}] An \code{mmes} model fitted with
\code{vsm(..., rotation=TRUE)}.
\end{ldescription}
\end{Arguments}
%
\begin{Details}\relax
If \eqn{Gu=U\Lambda U'}{} and \eqn{a}{} denotes the solver-basis effect, the
reported effect is \eqn{g=Ua}{}. Fitted \code{mmes} objects already expose this
original-basis effect in \code{uList}; this function provides the explicit
transformation from the retained engine output.
\end{Details}
%
\begin{Value}
A matrix of original-basis random effects, with relationship levels in rows
and covariance-product coordinates in columns.
\end{Value}
%
\begin{SeeAlso}\relax
\code{\LinkA{mmes}{mmes}}, \code{\LinkA{vsm}{vsm}}
\end{SeeAlso}
\inputencoding{utf8}
\HeaderA{rrm}{Reduced-rank covariance structure}{rrm}
%
\begin{Description}\relax
Reduced-rank covariance structure. The returned dimensionless covariance shape is intended for use inside \code{vsm}, which supplies the single overall variance scale.
\end{Description}
%
\begin{Usage}
\begin{verbatim}
rrm(x, k = 1L, loadings = NULL, fixed = NULL)
\end{verbatim}
\end{Usage}
%
\begin{Arguments}
\begin{ldescription}
\item[\code{x}] Variable defining q covariance levels.
\item[\code{k}] Reduced rank, satisfying 1 <= k < q.
\item[\code{loadings}] Optional finite q by k starting loading matrix. The leading
k by k block is made lower triangular and must have positive diagonal entries.
Loadings are divided by their first diagonal entry to remove the internal scale.
\item[\code{fixed}] Logical vector of length \eqn{qk-k(k-1)/2}{}, controlling the
triangular loading entries in column-major order. The first entry is always
fixed: the reference loading is one, and \code{vsm()} estimates the overall
variance scale. Other fixed entries constrain relative loadings.
\end{ldescription}
\end{Arguments}
%
\begin{Details}\relax
The reduced-rank structure uses
\deqn{K=\Lambda\Lambda^{\mathsf T},\qquad \Lambda_{11}=1,}{}
and \code{vsm()} supplies the overall variance scale:
\deqn{G=\sigma^2 K.}{}
All specific variances are fixed at exactly zero; no diagonal remainder or
numerical nugget is added. This is the same covariance family as ASReml's
default \code{rr(x, k)} model. Individual loadings can differ by sign or
rotation, so compare reconstructed covariance matrices when comparing fits.

The leading \eqn{k\times k}{} loading block is lower triangular for rotational
identification and its diagonal entries are positive. There are
\eqn{qk-k(k-1)/2-1}{} free shape parameters plus the overall scale.
Because \eqn{k<q}{}, the random-effect covariance is positive semidefinite
and singular by design.

The Henderson engine fits this structure through latent factors rather than
inverting the singular environment covariance. The direct-inversion formulation
(\code{control\$formulation="direct"}) also supports it, provided the total observation
covariance is positive definite. \code{rrm()} is intended for random effects,
not residual covariance structures.

Native latent-factor Henderson fits support \code{control\$solve\$method="ldlt"} and
\code{control\$solve\$method="cholmod"} (including automatic selection). They use
loading-aware likelihood scores and average-information updates in the
original loading/variance coordinates. The loading derivatives account for
the parameter-dependent latent design; the singular environment covariance
is never inverted. \code{control\$proposal\$stepWeight} damps covariance-parameter proposals,
while \code{control\$update\$blendWeight} blends EM/GEM or scalar PX-GEM with
average-information update directions, selected through \code{control\$update\$recovery},
as in ordinary \code{mmes()} fitting. Parameter standard errors use the unblended
average-information matrix at the final parameter values.
Prediction errors are mapped back from the latent coefficient equations
to the original random effects.

The likelihood trace terms are computed exactly by sharing random-design
and residual-block contractions across covariance parameters. Bounded panels
of Henderson solves avoid constructing a dense record-by-record covariance
or projection matrix. Random trace solves use the loading span (rank \code{k}),
rather than every environment direction. Reducing the coefficient dimension does not guarantee
a faster fit for every dataset.

Section-specific spatial residuals supplied by \code{dsumm()} are evaluated
on their observed residual blocks, including grids with missing cells and
unlisted sections with independent residuals. Residual covariance derivatives
and trace contractions never expand the full section-by-grid covariance.
\code{engineDiagnostics\$residualEvaluation} reports \code{"observed-blocks"}
and \code{engineDiagnostics\$largestResidualBlock} records the largest observed
block dimension used by the native RR AI evaluator.

Likelihood and AI calculations share the selected residual precision operator
and the live coefficient factorization, including grid residuals with
missing-cell absorption and eligible block-Schur solves. No separate
coefficient factorization or dense observed-block precision is built for AI.
Random trace panels use the latent Henderson identity to cancel relationship
covariance actions. Grid residual traces use block-local projection panels
rather than repeated full-observation precision applications.
\code{engineDiagnostics\$residualEngine} identifies the
shared residual engine, and \code{sharedConditionalOperators} confirms its
use. \code{aiElapsedSeconds} gives elapsed timings for the final AI
evaluation's setup, random traces, residual traces, and information matrix.

The positive-specific-variance Woodbury Schur formulation used for
\code{fam()} is unnecessary for exact RR: the specific-effect blocks are
absent, not inverted at zero. The smaller latent coefficient system still
uses the ordinary LDLT/CHOLMOD and eligible block-Schur dispatch.
With \code{control\$diagnostics\$verbose=TRUE}, the native fit reports its RR path, selected
coefficient solver, and then the progress table. Each accepted iteration
reports likelihood, wall-clock time, cumulative CPU seconds, the scheduled
\code{EMweight}, accepted-step \code{trustAlpha} and \code{lineAlpha},
and score norm. The step scales describe the update that reached the current
row (both are one at initialization); \code{lineAlpha} includes
\code{control\$proposal\$stepWeight} and any backtracking. Boundary-restraint and \code{notPD}
columns are omitted because the latent optimizer does not use those
ordinary covariance-update diagnostics. Numeric weights, step scales and
CPU times are retained in \code{engineDiagnostics\$iterationProgress}.

For native Henderson fits with weights \code{W}, let
\eqn{U=\mathrm{chol}(W)}{} in observation order. The effective residual
precision is \eqn{U R_0^{-1} U^{\mathsf T}}{}, where \eqn{R_0}{} is the
unweighted residual covariance. The score and AI use derivatives of the
corresponding covariance \eqn{U^{-\mathsf T}R_0U^{-1}}{}.

The native path does not currently support PCG, \code{vcc} constraints,
covariance-update acceleration, \code{control\$task\$mode="solve"}, or
\code{control\$representation\$factorScores} augmentation. These combinations
produce explicit errors.
The stored coefficient matrix \code{C} describes the latent equations;
\code{coefficientMap} maps latent coefficients to the original effect space.
Use \code{predict()}, \code{postPEV()}, and the fitted \code{uPevList} for
prediction uncertainty rather than directly indexing \code{C} by original
effect indices.

Earlier versions used \eqn{\Lambda\Lambda^{\mathsf T}+I_q}{}, estimating a
nonzero common specific variance on the fitted scale. To fit that former
covariance family, use \code{fam(x, k, variance="homogeneous")}.
\end{Details}
%
\begin{Value}
A list containing the incidence/design matrix in \code{Z} and a compiled CovarianceFactor v2 descriptor in \code{covFactor}, for use inside \code{vsm}.
\end{Value}
%
\begin{SeeAlso}\relax
\code{\LinkA{fam}{fam}}, \code{\LinkA{ownm}{ownm}}, \code{\LinkA{vsm}{vsm}}.
\end{SeeAlso}
%
\begin{Examples}
\begin{ExampleCode}
## Not run: 
vsm(rrm(environment, k=2), ism(genotype))
# Former rrm() model: an estimated common specific variance
vsm(fam(environment, k=2, variance="homogeneous"), ism(genotype))

## End(Not run)
\end{ExampleCode}
\end{Examples}
\inputencoding{utf8}
\HeaderA{run\_mme\_benchmark}{Run Experimental MME Benchmark Suite}{run.Rul.mme.Rul.benchmark}
%
\begin{Description}\relax
Executes a battery of workloads across candidate solver engines, measuring wall-clock time,
iterations, peak memory (if available), active numerical paths, and asserting numerical parity.
\end{Description}
%
\begin{Usage}
\begin{verbatim}
run_mme_benchmark(
  workloads,
  engines = NULL,
  replications = 3L,
  verifyParity = TRUE,
  tolerance = 1e-05,
  nIters = 5L,
  verbose = TRUE
)
\end{verbatim}
\end{Usage}
%
\begin{Arguments}
\begin{ldescription}
\item[\code{workloads}] A list of \code{mme\_workload} objects (or a single workload).

\item[\code{engines}] A named list of \code{\LinkA{mmes\_control}{mmes.Rul.control}} objects. Defaults to comparing
\code{baseline} (standard auto solver) against candidate experimental toggles.

\item[\code{replications}] Number of timed repetitions (warmup run excluded). Default is 3.

\item[\code{verifyParity}] Logical; whether to assert that experimental candidates match the reference
engine in likelihood, variance components, and BLUPs. Default is TRUE.

\item[\code{tolerance}] Numerical tolerance for parity verification. Default is 1e-5.

\item[\code{nIters}] Maximum REML/ML iterations for benchmark fits. Default is 5.

\item[\code{verbose}] Logical; whether to print progress. Default is TRUE.
\end{ldescription}
\end{Arguments}
%
\begin{Value}
An object of class \code{"sommer\_mme\_benchmark"}.
\end{Value}
\inputencoding{utf8}
\HeaderA{sar}{Simultaneous autoregressive spatial covariance structure}{sar}
%
\begin{Description}\relax
Simultaneous autoregressive spatial covariance structure. The returned dimensionless covariance shape is intended for use inside \code{vsm}, which supplies the single overall variance scale.
\end{Description}
%
\begin{Usage}
\begin{verbatim}
sar(x, W, rho = 0.10, fixed = FALSE)
\end{verbatim}
\end{Usage}
%
\begin{Arguments}
\begin{ldescription}
\item[\code{x}] Factor or design defining q spatial levels.
\item[\code{W}] Square spatial weights matrix aligned to the levels of x.  Row/column names are used when available.
\item[\code{rho}] Starting SAR dependence parameter.
\item[\code{fixed}] Logical indicating whether rho is fixed.
\end{ldescription}
\end{Arguments}
%
\begin{Details}\relax
The simultaneous autoregressive covariance is constructed from \deqn{B=I-\rho W,\qquad M=B^{-1}B^{-\mathsf T},\qquad K=M/M_{11}.}{} For a general real weights matrix, \code{rho} is restricted to the conservative interval \eqn{(-1/r(W),1/r(W))}{}, where \eqn{r(W)}{} is the spectral radius.  A bounded-logit working coordinate enforces the interval.  The first covariance derivative is supplied analytically to the generic CovarianceFactor engine.
\end{Details}
%
\begin{Value}
A list containing the incidence/design matrix in \code{Z} and a compiled CovarianceFactor v2 descriptor in \code{covFactor}, for use inside \code{vsm}.
\end{Value}
%
\begin{SeeAlso}\relax
\code{\LinkA{car}{car}}, \code{\LinkA{maternm}{maternm}}, \code{\LinkA{vsm}{vsm}}.
\end{SeeAlso}
%
\begin{Examples}
\begin{ExampleCode}
## Not run: 
vsm(sar(location, W), ism(genotype))

## End(Not run)
\end{ExampleCode}
\end{Examples}
\inputencoding{utf8}
\HeaderA{scores\_mmes}{Predict factor-analytic scores from a fitted mmes model}{scores.Rul.mmes}
%
\begin{Description}\relax
Predicts latent factor scores for each level of the main random-effect term
crossed with a \code{fam()} or \code{rrm()} covariance-shaping factor,
combining the fitted loadings (see \code{\LinkA{loadings\_mmes}{loadings.Rul.mmes}}), the fitted
covariance matrix, and the BLUPs already stored in \code{object\$uList}.
\end{Description}
%
\begin{Usage}
\begin{verbatim}
scores_mmes(object, term = NULL, method = c("regression", "bartlett"), 
            varianceScale=TRUE, rotation=TRUE)
\end{verbatim}
\end{Usage}
%
\begin{Arguments}
\begin{ldescription}
\item[\code{object}] a fitted model of class \code{"mmes"}.
\item[\code{term}] character name of the random term built with a single
\code{fam()} or \code{rrm()} covariance-shaping factor. If \code{NULL}, the
unique such term in the model is used automatically.
\item[\code{method}] \code{"regression"} (default) uses the full fitted covariance
\eqn{\Sigma}{} (Thomson's method); \code{"bartlett"} uses only the specific
(residual) variances \eqn{\Psi}{} (Bartlett's classic unbiased estimator).
\item[\code{varianceScale}] a logical argument to indicate if loadings should be returned 
in variance scale (multiplied by sqrt(sigma2)).
\item[\code{rotation}] a logical value to indicate if loadings should be rotated by its
singular vectors.
\end{ldescription}
\end{Arguments}
%
\begin{Details}\relax
Let \eqn{U}{} be the levels-by-levels BLUP matrix in \code{object\$uList[[term]]}
and \eqn{\Lambda}{} the loadings from \code{\LinkA{loadings\_mmes}{loadings.Rul.mmes}}. The
regression (Thomson) scores are
\deqn{F = U\Sigma^{-1}\Lambda,}{}
and the Bartlett scores are
\deqn{F = U\Psi^{-1}\Lambda(\Lambda^{\mathsf T}\Psi^{-1}\Lambda)^{-1}.}{}
Because \eqn{U}{} already contains shrunken (BLUP) effects, the resulting
scores are similarly regularized.

For \code{rrm()}, the covariance is singular and the specific variances are
zero. Regression scores instead recover the latent BLUP coordinates using
\deqn{F=U\Lambda(\Lambda^{\mathsf T}\Lambda)^{-1},}{}
so that \eqn{F\Lambda^{\mathsf T}=U}{}. This uses the returned loading scale
and rotation. Bartlett scores are undefined for this zero-specific-variance
model and produce an informative error.
\end{Details}
%
\begin{Value}
A matrix with one row per level of the main random-effect term and one
column per latent factor.
\end{Value}
%
\begin{SeeAlso}\relax
\code{\LinkA{loadings\_mmes}{loadings.Rul.mmes}}, \code{\LinkA{fam}{fam}}, \code{\LinkA{rrm}{rrm}}, \code{\LinkA{vsm}{vsm}}, \code{\LinkA{mmes}{mmes}}.
\end{SeeAlso}
%
\begin{Examples}
\begin{ExampleCode}
## Not run: 
mix <- mmes(BLUEs ~ trial,
            random = ~ vsm(fam(trial, 2), ism(genotype)),
            rcov = ~ units, data = dt)
scores_mmes(mix)

## End(Not run)
\end{ExampleCode}
\end{Examples}
\inputencoding{utf8}
\HeaderA{spl2Dc}{Two-dimensional penalised tensor-product of marginal B-Spline basis.}{spl2Dc}
\aliasA{spl2Da}{spl2Dc}{spl2Da}
\aliasA{spl2Db}{spl2Dc}{spl2Db}
%
\begin{Description}\relax

Auxiliary function used for modelling the spatial or environmental effect as a 
two-dimensional penalised tensor-product (isotropic approach) based on Lee et al. 
(2013) and Rodriguez-Alvarez et al. (2018). This is a modified wrapper of some 
portions of the SpATS package to build a single incidence matrix containing all
the columns from tensor products of the x and y coordinates and it fits such matrix 
as a single random effect. Then the heterogeneous covariances structure capabilities
of sommer can be used to enhance the model fit. You may be interested in reading and 
citing not only sommec but also Wageningen publications if using this 2D spline methodology.

\end{Description}
%
\begin{Usage}
\begin{verbatim}
spl2Dc(x.coord,y.coord,at.var=NULL,at.levels=NULL, type="PSANOVA", 
      nsegments = c(10,10), penaltyord = c(2,2), degree = c(3,3), 
      nestorder = c(1,1), thetaC=NULL, theta=NULL, sp=FALSE)
\end{verbatim}
\end{Usage}
%
\begin{Arguments}
\begin{ldescription}
\item[\code{x.coord}] vector of coordinates on the x-axis direction (i.e. row) to use in the 2 dimensional spline.
\item[\code{y.coord}] vector of coordinates on the y-axis direction (i.e. range or column) to use in the 2 dimensional spline.
\item[\code{at.var}] vector of indication variable where heterogeneous variance is required (e.g., a different spl2D for each field).
\item[\code{at.levels}] character vector with the names of the leves for the at term that should be used, if missing all levels are used.
\item[\code{type}] one of the two methods "PSANOVA" or "SAP". See details below.
\item[\code{nsegments}] numerical vector of length 2 containing the number of segments for each marginal (strictly \code{nsegments} - 1 is the number of internal knots in the domain of the covariate). Atomic values are also valid, being recycled. Default set to 10.
\item[\code{penaltyord}] numerical vector of length 2 containing the penalty order for each marginal. Atomic values are also valid, being recycled. Default set to 2 (second order). Currently, only second order penalties are allowed.
\item[\code{degree}] numerical vector of length 2 containing the order of the polynomial of the B-spline basis for each marginal. Atomic values are also valid, being recycled. Default set to 3 (cubic B-splines).
\item[\code{nestorder}] numerical vector of length 2 containing the divisor of the number of segments (\code{nsegments}) to be used for the construction of the nested B-spline basis for the smooth-by-smooth interaction component. In this case, the nested B-spline basis will be constructed assuming a total of \code{nsegments}/\code{nestorder} segments. Default set to 1, which implies that nested basis are not used. See \code{SAP} for more details.
\item[\code{thetaC}] an optional matrix for constraints in the variance components.
\item[\code{theta}] an optional matrix for initial values of the variance components.
\item[\code{sp}] a TRUE/FALSE statement to indicate if the VC from this structure should be multiplied by the scale parameter added in the mmes function through the addScaleParam argument in the \code{mmes} function .

\end{ldescription}
\end{Arguments}
%
\begin{Details}\relax

\strong{The following documentation is taken from the SpATS package. Please refer to this package and associated publications if you are interested in going deeper on this technique:}

Within the P-spline framework, anisotropic low-rank tensor-product smoothers have become the general approach for modelling multidimensional surfaces (Eilers and Marx 2003; Wood 2006). In the original SpATS package, was proposed to model the spatial or environmental effect by means of the tensor-product of B-splines basis functions. In other words, was proposed to model the spatial trend as a smooth bivariate surface jointly defined over the the spatial coordinates. Accordingly, the current function has been designed to allow the user to specify the spatial coordinates that the spatial trend is a function of. There is no restriction about how the spatial coordinates shall be specified: these can be the longitude and latitude of the position of the plot on the field or the column and row numbers. The only restriction is that the variables defining the spatial coordinates should be numeric (in contrast to factors).

As far as estimation is concerned, we have used in this package the equivalence between P-splines and linear mixed models (Currie and Durban, 2002). Under this approach, the smoothing parameters are expressed as the ratio between variance components. Moreover, the smooth components are decomposed in two parts: one which is not penalised (and treated as fixed) and one with is penalised (and treated as random). For the two-dimensional case, the mixed model representation leads also to a very interesting decomposition of the penalised part of the bivariate surface in three different components (Lee and Durban, 2011): (a) a component that contains the smooth main effect (smooth trend)	along one of the covariates that the surface is a function of (as, e.g, the x-spatial coordinate or column position of the plot in the field), (b) a component that contains the smooth main effect (smooth trend) along the other covariate (i.e., the y-spatial coordinate or row position); and (c) a smooth interaction component (sum of the linear-by-smooth interaction components and the smooth-by-smooth interaction component). 

The original implementation of SpATS assumes two different smoothing parameters, i.e., one for each covariate in the smooth component. Accordingly, the same smoothing parameters are used for both, the main effects and the smooth interaction. However, this approach can be extended to deal with the ANOVA-type decomposition presented in Lee and Durban (2011). In their approach, four different smoothing parameters are considered for the smooth surface, that are in concordance with the aforementioned decomposition: (a) two smoothing parameter, one for each of the main effects; and (b) two smoothing parameter for the smooth interaction component.

It should be noted that, the computational burden associated with the estimation of the two-dimensional tensor-product smoother might be prohibitive if the dimension of the marginal bases is large. In these cases, Lee et al. (2013) propose to reduce the computational cost by using nested bases. The idea is to reduce the dimension of the marginal bases (and therefore the associated number of parameters to be estimated), but only for the smooth-by-smooth interaction component. As pointed out by the authors, this simplification can be justified by the fact that the main effects would in fact explain most of the structure (or spatial trend) presented in the data, and so a less rich representation of the smooth-by-smooth interaction component could be needed. In order to ensure that the reduced bivariate surface is in fact nested to the model including only the main effects, Lee et al. (2013) show that the number of segments used for the nested basis should be a divisor of the number of segments used in the original basis (\code{nsegments} argument). In the present function, the divisor of the number of segments is specified through the argument \code{nestorder}. For a more detailed review on this topic, see Lee (2010) and Lee et al. (2013). The "PSANOVA" approach represents an alternative method. In this case, the smooth bivariate surface (or spatial trend) is decomposed in five different components each of them depending on a single smoothing parameter (see Lee et al., 2013).

-----------------

As mentioned at the beginning, the piece of documentation stated above was taken completely from the SpATS package in order to provide a deeper explanation. In practice, sommec uses some pieces of code from SpATS to build the design matrix containing all the columns from tensor products of the x and y coordinates and it fits such matrix as a single random effect. As a result the same variance component is assumed for the linear, linear by linear, linear by spline, and spline by spline interactions. This results in a less flexible approach than the one proposed by Rodriguez-Alvarez et al. (2018) but still makes a pretty good job to model the spatial variation. Use under your own risk.

\end{Details}
%
\begin{References}\relax

Rodriguez-Alvarez, M.X, Boer, M.P., van Eeuwijk, F.A., and Eilers, P.H.C. (2018). SpATS: Spatial Analysis of Field Trials with Splines. R package version 1.0-9. https://CRAN.R-project.org/package=SpATS.

Rodriguez-Alvarez, M.X., et al. (2015) Fast smoothng parameter separaton n multdmensonal generalzed P-splnes: the SAP algorthm. Statistics and Computing 25.5: 941-957.

Lee, D.-J., Durban, M., and Eilers, P.H.C. (2013). Efficient two-dimensional smoothing with P-spline ANOVA mixed models and nested bases. Computational Statistics and Data Analysis, 61, 22 - 37.

Gilmour, A.R., Cullis, B.R., and Verbyla, A.P. (1997). Accounting for Natural and Extraneous Variation in the Analysis of Field Experiments. Journal of Agricultural, Biological, and Environmental Statistics, 2, 269 - 293.

\end{References}
%
\begin{SeeAlso}\relax
\code{\LinkA{mmes}{mmes}}
\end{SeeAlso}
%
\begin{Examples}
\begin{ExampleCode}
## ============================ ##
## example to use spl2Dc() 
## ============================ ## 
data(DT_cpdata, package="enhancer")
# DT <- DT_cpdata
# GT <- GT_cpdata
# MP <- MP_cpdata
# A <- A.mat(GT)
## ============================ ##
## mimic 2 fields
## ============================ ## 
# aa <- DT; bb <- DT
# aa$FIELD <- "A";bb$FIELD <- "B"
# set.seed(1234)
# aa$Yield <- aa$Yield + rnorm(length(aa$Yield),0,4)
# DT2 <- rbind(aa,bb)
# head(DT2)
# mix <- mmes(Yield~1, control = mmes_control(formulation = "direct"),
#             random=~
#               vsm(dsm(FIELD),ism(Rowf)) +
#               vsm(dsm(FIELD),ism(Colf)) +
#                 spl2Dc(Row,Col,at.var=FIELD),
#             rcov=~vsm(dsm(FIELD),ism(units)),
#             data=DT2)
# 
# # extract spatial effects
# blup <- mix$uList$`spl2Dc(Row, Col, at.var = FIELD`
# head(blup) # 2 fields
# # recreate the incidence matrices
# xx=with(DT2, spl2Dc(Row,Col,at.var=FIELD))
# # get fitted values Zu for spatial effects and add them to the dataset
# field1 <- xx$Z$`A:all` %*% blup[,1]
# field2 <- xx$Z$`B:all` %*% blup[,2]
# DT2$spat <- field1+field2
# # plots the spatial effects
# lattice::levelplot(spat~Row*Col|FIELD, data=DT2)

\end{ExampleCode}
\end{Examples}
\inputencoding{utf8}
\HeaderA{spl2Dmats}{Get Tensor Product Spline Mixed Model Incidence Matrices}{spl2Dmats}
%
\begin{Description}\relax
\code{spl2Dmats} gets Tensor-Product P-Spline Mixed Model Incidence Matrices
for use with \code{sommer} and its main function \code{mmes}. We thank Sue Welham for making the TPSbits package available to the community. If you're using this function for your research please cite her TPSbits package :) this is mostly a wrapper of her tpsmmb function to enable the use in sommer.
\end{Description}
%
\begin{Usage}
\begin{verbatim}
spl2Dmats(
  x.coord.name,
  y.coord.name,
  data,
  at.name,
  at.levels, 
  nsegments=NULL,
  minbound=NULL,
  maxbound=NULL,
  degree = c(3, 3),
  penaltyord = c(2,2), 
  nestorder = c(1,1),
  method = "Lee"
)
\end{verbatim}
\end{Usage}
%
\begin{Arguments}
\begin{ldescription}
\item[\code{x.coord.name}] A string. Gives the name of \code{data} element
holding column locations.

\item[\code{y.coord.name}] A string. Gives the name of \code{data} element
holding row locations.

\item[\code{data}] A dataframe. Holds the dataset to be used for fitting.

\item[\code{at.name}] name of a variable defining if the 2D spline matrices should be created at different units (e.g., at different environments).

\item[\code{at.levels}] a vector of names indicating which levels of the at.name variable should be used for fitting the 2D spline function.

\item[\code{nsegments}] A list of length 2. Number of segments to split column and
row ranges into, respectively (= number of internal knots + 1). If only
one number is specified, that value is used in both dimensions. If not
specified, (number of unique values - 1) is used in each dimension;
for a grid layout (equal spacing) this gives a knot at each data value.

\item[\code{minbound}] A list of length 2. The lower bound to be used for column
and row dimensions respectively; default calculated as the minimum value
for each dimension.

\item[\code{maxbound}] A list of length 2. The upper bound to be used for column
and row dimensions respectively; default calculated as the maximum value
for each dimension.

\item[\code{degree}] A list of length 2. The degree of polynomial spline to be used
for column and row dimensions respectively; default=3.

\item[\code{penaltyord}] A list of length 2. The order of differencing for column
and row dimensions, respectively; default=2.

\item[\code{nestorder}] A list of length 2. The order of nesting for column and row
dimensions, respectively; default=1 (no nesting). A value of 2 generates
a spline with half the number of segments in that dimension, etc. The
number of segments in each direction must be a multiple of the order
of nesting.

\item[\code{method}] A string. Method for forming the  penalty; default=\code{"Lee"}
ie the penalty from Lee, Durban \& Eilers (2013, CSDA 61, 22-37). The
alternative method is \code{"Wood"} ie. the method from Wood et al (2012,
Stat Comp 23, 341-360).
This option is a research tool and requires further investigation.

\end{ldescription}
\end{Arguments}
%
\begin{Value}
List of length 7 elements:
\begin{enumerate}

\item{} \code{data} = the input data frame augmented with structures required
to fit tensor product splines in \code{asreml-R}. This data frame can be used
to fit the TPS model.

Added columns:
\begin{itemize}

\item{} \code{TP.col}, \code{TP.row} = column and row coordinates
\item{} \code{TP.CxR} = combined index for use with smooth x smooth term
\item{} \code{TP.C.n} for n=1:(diff.c) = X parts of column spline for use
in random model (where diff.c is the order of column differencing)
\item{} \code{TP.R.n} for n=1:(diff.r) = X parts of row spline for use in
random model (where diff.r is the order of row differencing)
\item{} \code{TP.CR.n} for n=1:((diff.c*diff.r)) = interaction between the
two X parts for use in fixed model. The first variate is
a constant term which should be omitted from the model when the constant
(1) is present. If all elements are
included in the model then the constant term should be omitted,
eg. \code{y \textasciitilde{} -1 + TP.CR.1 + TP.CR.2 + TP.CR.3 + TP.CR.4 + other terms...}
\item{} when \code{asreml="grp"} or \code{"sepgrp"}, the spline basis
functions are also added into the data frame. Column numbers for each
term are given in the \code{grp} list structure.

\end{itemize}



\item{} \code{fR} = Xr1:Zc
\item{} \code{fC} = Xr2:Zc
\item{} \code{fR.C} = Zr:Xc1
\item{} \code{R.fC} = Zr:Xc2 
\item{} \code{fR.fC} = Zc:Zr
\item{} \code{all} = Xr1:Zc | Xr2:Zc | Zr:Xc1 | Zr:Xc2 | Zc:Zr


\end{enumerate}


\end{Value}
%
\begin{Examples}
\begin{ExampleCode}

data("DT_cpdata", package="enhancer")
DT <- DT_cpdata
GT <- GT_cpdata
MP <- MP_cpdata
#### create the variance-covariance matrix
A <- A.mat(GT) # additive relationship matrix

M <- spl2Dmats(x.coord.name = "Col", y.coord.name = "Row", data=DT, nseg =c(14,21))
head(M$data)
# m1g <- mmes(Yield~1+TP.CR.2+TP.CR.3+TP.CR.4,
#             random=~Rowf+Colf+vsm(ism(M$fC))+vsm(ism(M$fR))+
#               vsm(ism(M$fC.R))+vsm(ism(M$C.fR))+vsm(ism(M$fC.fR))+
#               vsm(ism(id),Gu=A),
#             data=M$data, tolpar = 1e-6,
#             iters=30)
# 
# summary(m1g)$varcomp

\end{ExampleCode}
\end{Examples}
\inputencoding{utf8}
\HeaderA{stackTraits}{Stack several trait columns into the long format used by multi-trait models}{stackTraits}
\keyword{manip}{stackTraits}
%
\begin{Description}\relax
Converts a wide data frame (one column per trait) into the long format used by
multi-trait \code{\LinkA{mmes}{mmes}} models: one row per record and trait, a factor
identifying the trait, a single response column and a record key that pairs
the traits measured on the same experimental unit.
\end{Description}
%
\begin{Usage}
\begin{verbatim}
stackTraits(data, traits, keep = NULL, trait = "trait", value = "value",
            record = "record")
\end{verbatim}
\end{Usage}
%
\begin{Arguments}
\begin{ldescription}
\item[\code{data}] A data frame in wide format.
\item[\code{traits}] Character vector with the names of the trait columns to stack.
Their order defines the levels of the trait factor.
\item[\code{keep}] Columns of \code{data} copied to every stacked row. By default,
all columns that are not in \code{traits}.
\item[\code{trait}] Name of the new trait factor column.
\item[\code{value}] Name of the new response column.
\item[\code{record}] Name of the new record-key factor column.
\end{ldescription}
\end{Arguments}
%
\begin{Details}\relax
Rows are stacked trait by trait, so row \eqn{i}{} of trait \eqn{t}{} is row
\eqn{(t-1)n+i}{} of the output. The record key is taken from the row names of
\code{data} when these are unique and not the default \code{1..n}; otherwise
it is \code{"r1", "r2", ...}. Missing trait values are kept as \code{NA} (they
are handled by the \code{naMethodY} rules of \code{mmes}, so incomplete records
are allowed).

Using the record key in the residual,
\code{rcov = \textasciitilde{} vsm(usm(trait), ism(record))}, pairs the residuals of the traits
of each unit explicitly. With \code{ism(units)} the pairing is implied by the
order of the rows within each trait instead.
\end{Details}
%
\begin{Value}
A data frame with the \code{keep} columns and the \code{record},
\code{trait} and \code{value} columns, with \code{length(traits) * nrow(data)}
rows.
\end{Value}
%
\begin{SeeAlso}\relax
\code{\LinkA{mmes}{mmes}}, \code{\LinkA{covmatrix\_mmes}{covmatrix.Rul.mmes}}, \code{\LinkA{familym}{familym}}
\end{SeeAlso}
%
\begin{Examples}
\begin{ExampleCode}
set.seed(17)
n <- 60
DT <- data.frame(id = factor(rep(seq_len(n), each = 3)))
u <- matrix(rnorm(n * 2), n) 
DT$trait_a <- u[DT$id, 1] + rnorm(nrow(DT), sd = 0.7)
DT$trait_b <- u[DT$id, 2] + rnorm(nrow(DT), sd = 1)
DTL <- stackTraits(DT, traits = c("trait_a", "trait_b"), keep = "id")
head(DTL)
table(DTL$trait)

# bivariate model with the residual paired through the record key
fit <- mmes(value ~ trait,
  random = ~ vsm(usm(trait), ism(id)),
  rcov = ~ vsm(usm(trait), ism(record)),
  data = DTL,
  control = mmes_control(diagnostics = list(verbose = FALSE)))
covmatrix_mmes(fit, 1)$correlation   # genetic correlation
\end{ExampleCode}
\end{Examples}
\inputencoding{utf8}
\HeaderA{strm}{Covariance structure spanning several random terms}{strm}
\keyword{models}{strm}
%
\begin{Description}\relax
Combines two or more random terms built with \code{\LinkA{vsm}{vsm}} into one
random structure with an estimated covariance between the terms. Typical uses are
direct-maternal animal models, correlated
permanent-environment effects and indirect genetic effects.
\end{Description}
%
\begin{Usage}
\begin{verbatim}
strm(..., cov = usm, Gu = NULL, labels = NULL, sigma2 = NULL,
     fixedSigma2 = FALSE)
\end{verbatim}
\end{Usage}
%
\begin{Arguments}
\begin{ldescription}
\item[\code{...}] Two or more random \code{vsm()} terms, optionally named (names
become the term labels), e.g. \code{dir = vsm(ism(id)), mat = vsm(ism(dam))}.
\item[\code{cov}] A covariance constructor applied to the term index, e.g.
\code{usm} (default), \code{dsm}, \code{csm}, \code{corgm}, or
\code{function(x) fam(x, k = 1)}.
\item[\code{Gu}] Precision matrix shared by all terms (\code{attr(Gu, "inverse") = TRUE}).
If \code{NULL}, a \code{Gu} given inside the terms is used (all terms must
then use the same one) or an identity over the union of the terms' levels.
\item[\code{labels}] Term labels; default the argument names or \code{t1, t2, ...}.
\item[\code{sigma2}] Starting overall variance scale; \code{NULL} lets \code{mmes} choose.
\item[\code{fixedSigma2}] Fix the overall variance scale.
\end{ldescription}
\end{Arguments}
%
\begin{Details}\relax
For terms \eqn{u_1,\dots,u_T}{} with incidence matrices \eqn{Z_1,\dots,Z_T}{}
the structure is
\deqn{\mathrm{Var}(u_1,\dots,u_T)=\sigma^2\,K_{terms}\otimes K_{inner}\otimes A,}{}
where \eqn{K_{terms}}{} is given by \code{cov}, \eqn{K_{inner}}{} is the Kronecker
product of the covariance factors used inside the terms (e.g.
\code{dsm(env)}), and \eqn{A}{} is the relationship matrix. The requirements are
those of a direct product: all terms share the
coefficient levels (their levels are aligned on the union, or on the levels of
\code{Gu}), the same \code{Gu}, and the same inner covariance constructors with
the same levels; the parameters of the inner factors are shared. The
\code{sigma2} of individual terms is ignored. BLUPs are returned in one
\code{uList} element with a column per term (and inner level).
\code{covm(ran1, ran2)} is the two-term special case with \code{cov = usm}.

The terms do not need to be adjacent in the model
formula and their order is checked by level names rather than assumed.
\end{Details}
%
\begin{Value}
A random structure for the \code{random} formula of \code{\LinkA{mmes}{mmes}}.
\end{Value}
%
\begin{SeeAlso}\relax
\code{\LinkA{covm}{covm}}, \code{\LinkA{vsm}{vsm}}, \code{\LinkA{usm}{usm}}, \code{\LinkA{mmes}{mmes}}
\end{SeeAlso}
%
\begin{Examples}
\begin{ExampleCode}
set.seed(3)
ids <- paste0("i", 1:30)
A <- diag(30) + diag(0.2, 30); dimnames(A) <- list(ids, ids)
Ai <- solve(A); attr(Ai, "inverse") <- TRUE
d <- data.frame(id = factor(sample(ids, 200, TRUE), levels = ids),
                dam = factor(sample(ids[1:15], 200, TRUE), levels = ids))
d$y <- rnorm(200)
fit <- mmes(y ~ 1,
  random = ~ strm(dir = vsm(ism(id)), mat = vsm(ism(dam)), Gu = Ai),
  data = d,
  control = mmes_control(diagnostics = list(verbose = FALSE)))
covparams_mmes(fit, 1)
\end{ExampleCode}
\end{Examples}
\inputencoding{utf8}
\HeaderA{summary.mmer}{summary form a GLMM fitted with mmer}{summary.mmer}
\aliasA{print.mmer}{summary.mmer}{print.mmer}
\aliasA{print.summary.mmer}{summary.mmer}{print.summary.mmer}
\keyword{models}{summary.mmer}
%
\begin{Description}\relax
\code{summary} method for class \code{"mmer"}.
\end{Description}
%
\begin{Usage}
\begin{verbatim}
## S3 method for class 'mmer'
summary(object, ...)
\end{verbatim}
\end{Usage}
%
\begin{Arguments}
\begin{ldescription}
\item[\code{object}] an object of class \code{"mmer"}
\item[\code{...}] Further arguments to be passed
\end{ldescription}
\end{Arguments}
%
\begin{Value}
vector of summary

\end{Value}
%
\begin{Author}\relax
Giovanny Covarrubias-Pazaran
\end{Author}
%
\begin{SeeAlso}\relax
\code{\LinkA{summary}{summary}}, \code{\LinkA{mmer}{mmer}}
\end{SeeAlso}
\inputencoding{utf8}
\HeaderA{summary.mmes}{summary form a GLMM fitted with mmes}{summary.mmes}
\aliasA{print.mmes}{summary.mmes}{print.mmes}
\aliasA{print.summary.mmes}{summary.mmes}{print.summary.mmes}
\keyword{models}{summary.mmes}
%
\begin{Description}\relax
\code{summary} method for class \code{"mmes"}.
\end{Description}
%
\begin{Usage}
\begin{verbatim}
## S3 method for class 'mmes'
summary(object, ...)
\end{verbatim}
\end{Usage}
%
\begin{Arguments}
\begin{ldescription}
\item[\code{object}] an object of class \code{"mmes"}
\item[\code{...}] Further arguments to be passed
\end{ldescription}
\end{Arguments}
%
\begin{Value}
vector of summary

\end{Value}
%
\begin{Author}\relax
Giovanny Covarrubias-Pazaran
\end{Author}
%
\begin{SeeAlso}\relax
\code{\LinkA{summary}{summary}}, \code{\LinkA{mmes}{mmes}}
\end{SeeAlso}
\inputencoding{utf8}
\HeaderA{toeplitzm}{General positive-definite Toeplitz correlation structure}{toeplitzm}
%
\begin{Description}\relax
General positive-definite Toeplitz correlation structure. The returned dimensionless covariance shape is intended for use inside \code{vsm}, which supplies the single overall variance scale.
\end{Description}
%
\begin{Usage}
\begin{verbatim}
toeplitzm(x, pacf = NULL, fixed = NULL)
\end{verbatim}
\end{Usage}
%
\begin{Arguments}
\begin{ldescription}
\item[\code{x}] Ordered factor or design defining q covariance levels.
\item[\code{pacf}] Optional q-1 starting reflection coefficients / partial autocorrelations.  Defaults to 0.10.
\item[\code{fixed}] Logical vector of length q-1.
\end{ldescription}
\end{Arguments}
%
\begin{Details}\relax
A full q by q Toeplitz correlation matrix is parameterized by q-1 reflection coefficients.  Each PACF is mapped from an unconstrained working coordinate with \code{tanh}; the reflection coefficients are converted to a stable autoregressive representation and then to the implied Toeplitz correlation sequence.  This parameterization preserves positive definiteness while allowing every lag correlation to vary indirectly.
\end{Details}
%
\begin{Value}
A list containing the incidence/design matrix in \code{Z} and a compiled CovarianceFactor v2 descriptor in \code{covFactor}, for use inside \code{vsm}.
\end{Value}
%
\begin{SeeAlso}\relax
\code{\LinkA{ar1m}{ar1m}}, \code{\LinkA{ar2m}{ar2m}}, \code{\LinkA{ar3m}{ar3m}}, \code{\LinkA{mam}{mam}}, \code{\LinkA{vsm}{vsm}}.
\end{SeeAlso}
%
\begin{Examples}
\begin{ExampleCode}
## Not run: 
vsm(toeplitzm(time), ism(id))

## End(Not run)
\end{ExampleCode}
\end{Examples}
\inputencoding{utf8}
\HeaderA{tpsmmbwrapper}{Get Tensor Product Spline Mixed Model Incidence Matrices}{tpsmmbwrapper}
%
\begin{Description}\relax
\code{tpsmmbwrapper} is a wrapper of tpsmmb function from the TPSbits package to avoid version dependencies but if you're using this function for your research please cite the TPSbits package. This function is internally used by the \code{spl2Dmatrices} function to get Tensor-Product P-Spline Mixed Model Bits
(design matrices) for use with \code{sommer}.
\end{Description}
%
\begin{Usage}
\begin{verbatim}
tpsmmbwrapper(
  columncoordinates,
  rowcoordinates,
  data,
  nsegments=NULL,
  minbound=NULL,
  maxbound=NULL,
  degree = c(3, 3),
  penaltyord = c(2, 2),
  nestorder = c(1, 1),
  asreml = "mbf",
  eigenvalues = "include",
  method = "Lee",
  stub = NULL
)
\end{verbatim}
\end{Usage}
%
\begin{Arguments}
\begin{ldescription}
\item[\code{columncoordinates}] A string. Gives the name of \code{data} element
holding column locations.

\item[\code{rowcoordinates}] A string. Gives the name of \code{data} element
holding row locations.

\item[\code{data}] A dataframe. Holds the dataset to be used for fitting.

\item[\code{nsegments}] A list of length 2. Number of segments to split column and
row ranges into, respectively (= number of internal knots + 1). If only
one number is specified, that value is used in both dimensions. If not
specified, (number of unique values - 1) is used in each dimension;
for a grid layout (equal spacing) this gives a knot at each data value.

\item[\code{minbound}] A list of length 2. The lower bound to be used for column
and row dimensions respectively; default calculated as the minimum value
for each dimension.

\item[\code{maxbound}] A list of length 2. The upper bound to be used for column
and row dimensions respectively; default calculated as the maximum value
for each dimension.

\item[\code{degree}] A list of length 2. The degree of polynomial spline to be used
for column and row dimensions respectively; default=3.

\item[\code{penaltyord}] A list of length 2. The order of differencing for column
and row dimensions, respectively; default=2.

\item[\code{nestorder}] A list of length 2. The order of nesting for column and row
dimensions, respectively; default=1 (no nesting). A value of 2 generates
a spline with half the number of segments in that dimension, etc. The
number of segments in each direction must be a multiple of the order
of nesting.

\item[\code{asreml}] A string. Indicates the types of structures to be generated
for use in asreml models; default \code{"mbf"}. The
appropriate eigenvalue scaling is included within the Z matrices unless
setting \code{scaling="none"} is used, and then the scaling factors are
supplied separately in the returned object.
\begin{itemize}

\item{} \code{asreml="mbf"} indicates the function should put the
spline design matrices into structures for use with \code{"mbf"};
\item{} \code{asreml="grp"} indicates the function should add the
composite spline design matrices (eg. for second-order differencing,
matrices Xr1:Zc, Xr2:Zc, Zr:Xc1, Zr:Xc2 and Zc:Zr) into the data frame
and provide a group list structure for each term;
\item{} \code{asreml="sepgrp"} indicates the function should generate the
individual X and Z spline design matrices separately (ie. Xc, Xr, Zc and
Zr), plus the smooth x smooth interaction term as a whole (ie. Zc:Zr),
and provide a group list structure for each term.
\item{} \code{asreml="own"} indicates the function should generate the
composite matrix ( Xr:Zc | Zr:Xc | Zc:Zr ) as a single set of columns.

\end{itemize}


\item[\code{eigenvalues}] A string. Indicates whether eigenvalues should be
included within the Z design matrices \code{eigenvalues="include"}, or
whether this scaling should be omitted (\code{eigenvalues="omit"});
default \code{eigenvalues="include"}. If the eigenvalue scaling is
omitted from the Z design matrices, then it should instead be included in
the model as a variance structure to obtain the correct TPspline model.

\item[\code{method}] A string. Method for forming the  penalty; default=\code{"Lee"}
ie the penalty from Lee, Durban \& Eilers (2013, CSDA 61, 22-37). The
alternative method is \code{"Wood"} ie. the method from Wood et al (2012,
Stat Comp 23, 341-360).
This option is a research tool and requires further investigation.

\item[\code{stub}] A string. Stub to be attached to names in the \code{mbf} list to
avoid over-writing structures and general confusion.
\end{ldescription}
\end{Arguments}
%
\begin{Value}
List of length 7, 8 or 9 (according to the \code{asreml} and
\code{eigenvalues} parameter settings).
\begin{enumerate}

\item{} \code{data} = the input data frame augmented with structures required
to fit tensor product splines in \code{asreml-R}. This data frame can be used
to fit the TPS model.

Added columns:
\begin{itemize}

\item{} \code{TP.col}, \code{TP.row} = column and row coordinates
\item{} \code{TP.CxR} = combined index for use with smooth x smooth term
\item{} \code{TP.C.n} for n=1:(diff.c) = X parts of column spline for use
in random model (where diff.c is the order of column differencing)
\item{} \code{TP.R.n} for n=1:(diff.r) = X parts of row spline for use in
random model (where diff.r is the order of row differencing)
\item{} \code{TP.CR.n} for n=1:((diff.c*diff.r)) = interaction between the
two X parts for use in fixed model. The first variate is
a constant term which should be omitted from the model when the constant
(1) is present. If all elements are
included in the model then the constant term should be omitted,
eg. \code{y \textasciitilde{} -1 + TP.CR.1 + TP.CR.2 + TP.CR.3 + TP.CR.4 + other terms...}
\item{} when \code{asreml="grp"} or \code{"sepgrp"}, the spline basis
functions are also added into the data frame. Column numbers for each
term are given in the \code{grp} list structure.

\end{itemize}

\item{} \code{mbflist} = list that can be used in call to asreml (so long as Z
matrix data frames extracted with right names, eg BcZ<stub>.df)
\item{} \code{BcZ.df} = mbf data frame mapping onto smooth part of column
spline, last column (labelled \code{TP.col}) gives column index
\item{} \code{BrZ.df} = mbf data frame mapping onto smooth part of row spline,
last column (labelled \code{TP.row}) gives row index
\item{} \code{BcrZ.df} = mbf data frame mapping onto smooth x smooth term, last
column (labelled \code{TP.CxR}) maps onto col x row combined index
\item{} \code{dim} = list structure, holding dimension values relating to the
model:
\begin{enumerate}

\item{} \code{"diff.c"} = order of differencing used in column dimension
\item{} \code{"nbc"} = number of random basis functions in column dimension
\item{} \code{"nbcn"} = number of nested random basis functions in column dimension
used in smooth x smooth term
\item{} \code{"diff.r"} = order of differencing used in column dimension
\item{} \code{"nbr"} = number of random basis functions in column dimension
\item{} \code{"nbrn"} = number of nested random basis functions in column dimension
used in smooth x smooth term

\end{enumerate}

\item{} \code{trace} = list of trace values for ZGZ' for the random TPspline
terms, where Z is the design matrix and G
is the known diagonal variance matrix derived from eigenvalues. This can
be used to rescale the spline design matrix (or equivalently variance
components).
\item{} \code{grp} = list structure, only added for settings
\code{asreml="grp"},  \code{asreml="sepgrp"} or \code{asreml="own"}.
For \code{asreml="grp"}, provides column indexes for each of the 5
random components of the 2D splines.
For \code{asreml="sepgrp"}, provides column indexes for each of the X and
Z component matrices for the 1D splines, plus the composite smooth x
smooth interaction term. For \code{asreml="own"}, provides column indexes
for the composite random model.
Dimensions of the components can be derived from the values in the
\code{dim} item.  The Z terms are scaled by the associated
eigenvalues when \code{eigenvalues="include"}, but not when
\code{eigenvalues="omit"}.
\item{} \code{eigen} = list structure, only added for option setting
\code{eigenvalues="omit"}. Holds the diagonal elements of the inverse
variance matrix for the terms Xc:Zr (called \code{diagr}), Zc:Xr
(called \code{diagc}) and Zc:Zr (called \code{diagcr}).

\end{enumerate}

\end{Value}
\inputencoding{utf8}
\HeaderA{unsm}{unstructured indication matrix}{unsm}
\aliasA{unsm2}{unsm}{unsm2}
%
\begin{Description}\relax
\code{unsm} creates a square matrix with ones in the diagonals and 2's in the off-diagonals to quickly specify an unstructured constraint in the Gtc argument of the \code{\LinkA{vsm}{vsm}} function.
\end{Description}
%
\begin{Usage}
\begin{verbatim}
  unsm(x, reps=NULL)
\end{verbatim}
\end{Usage}
%
\begin{Arguments}
\begin{ldescription}
\item[\code{x}] integer specifying the number of traits to be fitted for a given random effect.
\item[\code{reps}] integer specifying the number of times the matrix should be repeated in a list format to provide easily the constraints in complex models that use the ds(), us() or cs() structures.
\end{ldescription}
\end{Arguments}
%
\begin{Value}
\begin{description}


\item[\$res]  a matrix or a list of matrices with the constraints to be provided in the Gtc argument of the \code{\LinkA{vsm}{vsm}} function.

\end{description}

\end{Value}
%
\begin{Author}\relax
Giovanny Covarrubias-Pazaran
\end{Author}
%
\begin{References}\relax
Covarrubias-Pazaran G (2016) Genome assisted prediction of quantitative traits using the R package sommer. PLoS ONE 11(6): doi:10.1371/journal.pone.0156744 
\end{References}
%
\begin{Examples}
\begin{ExampleCode}
unsm(3)
unsm(3,2)
\end{ExampleCode}
\end{Examples}
\inputencoding{utf8}
\HeaderA{usm}{Unstructured positive-definite covariance structure}{usm}
%
\begin{Description}\relax
Unstructured positive-definite covariance structure. The returned dimensionless covariance shape is intended for use inside \code{vsm}, which supplies the single overall variance scale.
\end{Description}
%
\begin{Usage}
\begin{verbatim}
usm(x, theta = NULL, fixed = NULL)
\end{verbatim}
\end{Usage}
%
\begin{Arguments}
\begin{ldescription}
\item[\code{x}] Variable or design defining q covariance levels.
\item[\code{theta}] Optional positive-definite q by q starting covariance matrix.
\item[\code{fixed}] Logical vector of length q(q+1)/2 - 1 controlling the estimable normalized-Cholesky parameters.
\end{ldescription}
\end{Arguments}
%
\begin{Details}\relax
The covariance shape is parameterized as \deqn{K=LL^{\mathsf T},}{} with lower-triangular \eqn{L}{} and \eqn{L_{11}=1}{}.  Diagonal elements after the first are positive and represented on log scales; lower off-diagonal elements are unrestricted.  This guarantees positive definiteness without repeatedly repairing an unconstrained covariance matrix.
\end{Details}
%
\begin{Value}
A list containing the incidence/design matrix in \code{Z} and a compiled CovarianceFactor v2 descriptor in \code{covFactor}, for use inside \code{vsm}.
\end{Value}
%
\begin{SeeAlso}\relax
\code{\LinkA{vsm}{vsm}}, \code{\LinkA{mmes}{mmes}}.
\end{SeeAlso}
%
\begin{Examples}
\begin{ExampleCode}
## Not run: 
vsm(usm(environment), ism(genotype))

## End(Not run)
\end{ExampleCode}
\end{Examples}
\inputencoding{utf8}
\HeaderA{usr}{Unstructured covariance structure for \code{mmer} and \code{vsr}}{usr}
%
\begin{Description}\relax
Creates an unstructured variance-covariance component pattern for use by the \code{mmer} covariance-model interface, typically inside \code{\LinkA{vsr}{vsr}}.
\end{Description}
%
\begin{Usage}
\begin{verbatim}
usr(x)
\end{verbatim}
\end{Usage}
%
\begin{Arguments}
\begin{ldescription}
\item[\code{x}] 
A factor, character vector, numeric vector, or design/incidence matrix defining the covariance dimension. Factor and character inputs are expanded to incidence columns. A matrix is used directly.

\end{ldescription}
\end{Arguments}
%
\begin{Details}\relax
\code{usr()} is part of the covariance-structure interface used by \code{mmer}/\code{vsr}; it is distinct from \code{\LinkA{usm}{usm}}, which belongs to the newer \code{mmes}/\code{vsm} CovarianceFactor interface.

After constructing the design matrix \eqn{Z}{}, the function calls \code{unsm(q)}, where \eqn{q}{} is the number of columns of \eqn{Z}{}, to create the \code{thetaC} pattern for an unstructured covariance matrix. Thus the associated covariance model permits a separate variance for every level and a separate covariance for every pair of levels:
\deqn{
\Sigma =
\begin{bmatrix}
\sigma_1^2 & \sigma_{12} & \cdots & \sigma_{1q}\\
\sigma_{12} & \sigma_2^2 & \cdots & \sigma_{2q}\\
\vdots & \vdots & \ddots & \vdots\\
\sigma_{1q} & \sigma_{2q} & \cdots & \sigma_q^2
\end{bmatrix}.
}{}
The exact component coding is stored in the matrix returned by \code{unsm(q)} and passed to the \code{vsr} machinery through \code{thetaC}.

For a factor or character vector, missing observations are retained through the construction of the incidence matrix. For a numeric non-factor input, \code{x} is treated as a single design column.
\end{Details}
%
\begin{Value}
A list with components:
\begin{itemize}

\item{} \code{Z}: the design/incidence matrix associated with the covariance dimension.
\item{} \code{thetaC}: the unstructured covariance-component pattern produced by \code{unsm()}.

\end{itemize}

\end{Value}
%
\begin{SeeAlso}\relax
\code{\LinkA{vsr}{vsr}}, \code{\LinkA{dsr}{dsr}}, \code{\LinkA{csr}{csr}}, \code{\LinkA{atr}{atr}},
\code{\LinkA{usm}{usm}}, \code{\LinkA{mmer}{mmer}}
\end{SeeAlso}
%
\begin{Examples}
\begin{ExampleCode}
# Unstructured covariance pattern across three levels
U <- usr(factor(c("A","B","A","C")))
U$Z
U$thetaC
\end{ExampleCode}
\end{Examples}
\inputencoding{utf8}
\HeaderA{vcc}{Equality and scaling constraints between variance parameters}{vcc}
\aliasA{vcParams}{vcc}{vcParams}
\keyword{models}{vcc}
%
\begin{Description}\relax
The \code{vcc} argument of \code{\LinkA{mmes}{mmes}} constrains groups of covariance
parameters to be equal or to keep fixed ratios.
Parameters are identified through the \code{vcParams} table.
\end{Description}
%
\begin{Details}\relax
Parameters are estimated on transformed scales that guarantee valid
covariance matrices (log for variances, atanh or bounded logits for
correlations). A constraint is enforced exactly when it is affine on that
scale, which covers equalities of any parameter and ratios of positive
parameters (a log-scale offset). Relations such as
\eqn{\sigma^2_3=\sigma^2_1+\sigma^2_2}{} are not affine on the log scale and are not supported.

Internally the free working parameters are written as \eqn{w=o+T\phi}{}. Each
average-information update is projected onto this set in the information
metric, which equals the constrained Newton step, and then repaired exactly;
both the Henderson and the direct engines support constraints. The covariance
of the estimates is \eqn{T(T^\top\mathcal{I}T)^{-1}T^\top}{}, so constrained
parameters share their standard errors.
\end{Details}
%
\begin{Section}{The vcParams table}
\code{mmes(..., control = mmes\_control(output = list(setup = TRUE)))\$vcParams}
(before fitting) and
\code{fit\$vcParams} (after fitting) list every covariance parameter with
columns \code{index}, \code{term}, \code{parameter}, \code{kind},
\code{start} (natural scale), \code{free}, \code{structure} and \code{position}
(and \code{group} when constraints were used). \code{kind} is \code{"scale"}
for the variance scale (\code{sigma2}) owned by each \code{vsm()} term, or the
transform used to estimate the parameter: \code{"exp"} (positive quantities
such as ranges and variance ratios), \code{"identity"} (e.g. Cholesky or
factor-analytic loadings), \code{"tanh"} or \code{"bounded\_logit"}
(correlation-type parameters).
\end{Section}
%
\begin{Section}{Specifying constraints}
\code{vcc} is a data frame with columns
\begin{description}

\item[parameter] the \code{index} of the parameter, or its name
\code{"term:parameter"} (e.g. \code{"vsm(ism(B)):sigma2"}); the
\code{parameter} name alone is accepted when unique.
\item[group] parameters with the same group value are constrained together.
\item[scale] optional, default 1: each member equals \code{scale} times the
first member of its group (natural scale). Scaling is allowed for
\code{"scale"}, \code{"exp"} and \code{"identity"} parameters; correlation-type
parameters can only be equated.

\end{description}

All members of a group must have the same \code{kind}: variance scales can be
constrained with other variance scales, and other parameters with parameters
estimated through the same transform. If any member of a group is fixed, the
whole group is fixed at the value implied by that member. Starting values of
the members are made consistent with the first (or fixed) member.
\end{Section}
%
\begin{SeeAlso}\relax
\code{\LinkA{mmes}{mmes}}, \code{\LinkA{vsm}{vsm}}, \code{\LinkA{covparams\_mmes}{covparams.Rul.mmes}}
\end{SeeAlso}
%
\begin{Examples}
\begin{ExampleCode}
data(DT_yatesoats, package = "enhancer")
DT <- DT_yatesoats
p <- mmes(Y ~ V * N,
  random = ~ B + B:MP,
  rcov = ~ units,
  data = DT,
  control = mmes_control(output = list(setup = TRUE)))
p$vcParams
# block and whole-plot variances equal
m1 <- mmes(Y ~ V * N,
  random = ~ B + B:MP,
  rcov = ~ units,
  data = DT,
  vcc = data.frame(parameter = c(1, 2), group = 1),
  control = mmes_control(diagnostics = list(verbose = FALSE)))
# block variance twice the whole-plot variance
m2 <- mmes(Y ~ V * N,
  random = ~ B + B:MP,
  rcov = ~ units,
  data = DT,
  vcc = data.frame(parameter = c("vsm(ism(B:MP)):sigma2", "vsm(ism(B)):sigma2"),
                            group = 1, scale = c(1, 2)),
  control = mmes_control(diagnostics = list(verbose = FALSE)))
unlist(m2$covPar)
\end{ExampleCode}
\end{Examples}
\inputencoding{utf8}
\HeaderA{vpredict}{vpredict form of a LMM fitted with mmes}{vpredict}
\methaliasA{vpredict.mmes}{vpredict}{vpredict.mmes}
\keyword{models}{vpredict}
%
\begin{Description}\relax
\code{vpredict} method for class \code{"mmes"}.

Post-analysis procedure to calculate linear combinations of variance components. Its intended use is when the variance components are either simple variances or are variances and covariances in an unstructured matrix. The functions covered are linear combinations of the variance components (for example, phenotypic variance), a ratio of two components (for example, heritabilities) and the correlation based on three components (for example, genetic correlation). 

The calculations are based on the estimated variance parameters and their variance matrix as represented by the inverse of the Fisher or Average information matrix. Note that this matrix has zero values for fixed variance parameters including those near the parameter space boundary.

The transform is specified with a formula. On the left side of the formula is a name for the transformation. On the right side of the formula is a transformation specified with shortcut names like `V1`, `V2`, etc.  The easiest way to identify these shortcut names is to use `summary(object)\$varcomp`. The rows of this object can referred to with shortcuts `V1`, `V2`, etc.  See the example below.

\end{Description}
%
\begin{Usage}
\begin{verbatim}

vpredict(object, transform)
## S3 method for class 'mmes'
vpredict(object, transform)

\end{verbatim}
\end{Usage}
%
\begin{Arguments}
\begin{ldescription}

\item[\code{object}] a model fitted with the mmes function.

\item[\code{transform}] a formula to calculate the function.

\end{ldescription}
\end{Arguments}
%
\begin{Details}\relax

The delta method (e.g., Lynch and Walsh 1998, Appendix 1; Ver Hoef 2012) uses
a Taylor series expansion to approximate the moments of a function of
parameters. Here, a second-order Taylor series expansion is implemented to
approximate the standard error for a function of (co)variance parameters.
Partial first derivatives of the function are calculated by algorithmic
differentiation with \code{\LinkA{deriv}{deriv}}.

Though \code{vpredict} can calculate standard errors for non-linear functions
of (co)variance parameters from a fitted \code{mmes} model, it is limited
to non-linear functions constructed by mathematical operations such as the
arithmetic operators \code{+}, \code{-}, \code{*}, \code{/} and \code{\textasciicircum{}},
and single-variable functions such as  \code{exp} and \code{log}. See 
\code{\LinkA{deriv}{deriv}} for more information.


\end{Details}
%
\begin{Value}

\begin{ldescription}
\item[\code{dd}] the parameter and its standard error.

\end{ldescription}
\end{Value}
%
\begin{Author}\relax
Giovanny Covarrubias
\end{Author}
%
\begin{References}\relax
Covarrubias-Pazaran G (2016) Genome assisted prediction of quantitative traits using the R package sommer. PLoS ONE 11(6): doi:10.1371/journal.pone.0156744 

Lynch, M. and B. Walsh 1998. Genetics and Analysis of Quantitative Traits. Sinauer Associates, Inc., Sunderland, MA, USA.

Ver Hoef, J.M. 2012. Who invented the delta method? The American Statistician 66:124-127. DOI: 10.1080/00031305.2012.687494


\end{References}
%
\begin{SeeAlso}\relax
\code{\LinkA{vpredict}{vpredict}}, \code{\LinkA{mmes}{mmes}}
\end{SeeAlso}
%
\begin{Examples}
\begin{ExampleCode}

####=========================================####
####=========================================####
#### EXAMPLE 1
#### simple example with univariate models
####=========================================####
####=========================================####
# data(DT_cpdata, package="enhancer")
# DT <- DT_cpdata
# GT <- GT_cpdata
# MP <- MP_cpdata
# #### create the variance-covariance matrix 
# A <- A.mat(GT)
# #### look at the data and fit the model
# head(DT)
# mix1 <- mmes(Yield~1,
#               random=~vsm(ism(id),Gu=A), 
#               data=DT)
# summary(mix1)$varcomp
# #### run the vpredict function
# vpredict(mix1, h2 ~ V1 / ( V1 + V2 ) )


\end{ExampleCode}
\end{Examples}
\inputencoding{utf8}
\HeaderA{vs}{variance structure specification}{vs}
%
\begin{Description}\relax
\code{vs} DEPRECATED NOW. was the main function to build the variance-covariance structure for the random effects to be fitted in the \code{\LinkA{mmer}{mmer}} solver. 
\end{Description}
%
\begin{Usage}
\begin{verbatim}
  vs(..., Gu=NULL, Gti=NULL, Gtc=NULL, reorderGu=TRUE, buildGu=TRUE)
\end{verbatim}
\end{Usage}
%
\begin{Arguments}
\begin{ldescription}
\item[\code{...}] variance structure to be specified following the logic desired in the internal kronecker product. For example, if user wants to define a diagonal variance structure for the random effect 'genotypes'(g) with respect to a random effect 'environments'(e), this is:

\code{var(g) = G.e @ I.g}

being \code{G.e}  a matrix containing the variance covariance components for g (genotypes) in each level of e (environments), \code{I.g} is the covariance among levels of g (genotypes; i.e. relationship matrix), and @ is the kronecker product. This would be specified in the mmer solver as:

\code{random=\textasciitilde{}vs(dsr(e),g)}

One strength of sommer is the ability to specify very complex structures with as many kronecker products as desired. For example:

\code{var(g) = G.e @ G.f @ G.h @ I.g}

is equivalent to

\code{random=\textasciitilde{}vs(e,f,h,g)}

where different covariance structures can be applied to the levels of \code{e,f,h} or a combination of these). For more examples please see the vignettes 'sommer.start' available in the package.

\item[\code{Gu}] matrix with the known variance-covariance values for the levels of the u.th random effect (i.e. relationship matrix among individuals or any other known covariance matrix). If \code{NULL}, then an identity matrix is assumed. The Gu matrix can have more levels than the ones present in the random effect linked to it but not the other way around. Otherwise, an error message of missing level in Gu will be returned.

\item[\code{Gti}] matrix with dimensions t x t (t equal to number of traits) with initial values of the variance-covariance components for the random effect specified in the .... argument. If \code{NULL} the program will provide the initial values. The values need to be scaled, see Details section.

\item[\code{Gtc}] matrix with dimensions t x t (t equal to number of traits) of constraints for the variance-covariance components for the random effect specified in the ... argument according to the following rules:

\code{0: not to be estimated}

\code{1: estimated and constrained to be positive (i.e. variance component)}

\code{2: estimated and unconstrained (can be negative or positive, i.e. covariance component)}

\code{3: not to be estimated but fixed (value has to be provided in the Gti argument)}

In the multi-response scenario if the user doesn't specify this argument the default is to build an unstructured matrix (using the \code{\LinkA{unsm}{unsm}}() function). This argument needs to be used wisely since some covariance among responses may not make sense. Useful functions to specify constraints are; \code{\LinkA{diag}{diag}}(), \code{\LinkA{unsm}{unsm}}(), \code{\LinkA{fixm}{fixm}}().


\item[\code{reorderGu}] a \code{TRUE/FALSE} statement if the Gu matrix should be reordered based on the names of the design matrix of the random effect or passed with the custom order of the user. This may be important when fitting covariance components in a customized fashion. Only for advanced users.

\item[\code{buildGu}] a \code{TRUE/FALSE} statement to indicate if the Gu matrix should be built in R when the value for the argument \code{Gu=NULL}. Repeat, only when when the value for the argument \code{Gu} is equal to \code{NULL}. In some cases when the incidence matrix is wide (e.g. rrBLUP models) the covariance structure is a huge p x p matrix that can be avoided when performing matrix operations. By setting this argument to \code{FALSE} it allows to skip forming this covariance matrix.

\end{ldescription}
\end{Arguments}
%
\begin{Details}\relax
When providing initial values in the \code{Gti} argument the user has to provide scaled variance component values. The user can provide values from a previous model by accessing the \code{sigma\_scaled} output from an \code{mmer} model or if an specific value is desired the user can obtain the scaled value as:

\code{m = x/var(y)}

where \code{x} is the desired initial value and \code{y} is the response variable. You can find an example in the \code{\LinkA{DT\_cpdata}{DT.Rul.cpdata}} dataset.

\end{Details}
%
\begin{Value}
\begin{description}

\item[\$res] a list with all neccesary elements (incidence matrices, known var-cov structures, unknown covariance structures to be estimated and constraints) to be used in the mmer solver.

\end{description}

\end{Value}
%
\begin{Author}\relax
Giovanny Covarrubias-Pazaran
\end{Author}
%
\begin{References}\relax
Covarrubias-Pazaran G (2016) Genome assisted prediction of quantitative traits using the R package sommer. PLoS ONE 11(6): doi:10.1371/journal.pone.0156744 

Covarrubias-Pazaran G (2018) Software update: Moving the R package sommer to
multivariate mixed models for genome-assisted prediction. doi:
https://doi.org/10.1101/354639
\end{References}
%
\begin{SeeAlso}\relax
The core function of the package: \code{\LinkA{mmer}{mmer}}
\end{SeeAlso}
%
\begin{Examples}
\begin{ExampleCode}

data(DT_example, package="enhancer")
DT <- DT_example
A <- A_example

## ============================ ##
## example to without structure
## ============================ ##

mix <- mmer(Yield~Env, 
            random= ~ vs(Name),
            rcov=~ vs(units),
            data=DT)

\end{ExampleCode}
\end{Examples}
\inputencoding{utf8}
\HeaderA{vsm}{Variance-structure model with arbitrary Kronecker covariance factors}{vsm}
%
\begin{Description}\relax
\code{vsm} constructs random-effect or residual covariance structures for
\code{mmes}.  Covariance-shaping terms are represented by CovarianceFactor v2
descriptors and may be combined in an arbitrary-depth Kronecker product.  A
single product-level variance, \eqn{\sigma^2}{}, is owned by \code{vsm}, which
avoids scale confounding among the individual covariance factors.
\end{Description}
%
\begin{Usage}
\begin{verbatim}
vsm(..., Gu = NULL, sigma2 = NULL, fixedSigma2 = FALSE,
  rotation = FALSE, isFixed = FALSE, verbose = TRUE)
\end{verbatim}
\end{Usage}
%
\begin{Arguments}
\begin{ldescription}
\item[\code{...}] One or more covariance-constructor terms.  The last term supplies
the main-effect incidence matrix.  All preceding terms are covariance-shaping
factors.  Examples include \code{dsm(environment)}, \code{ar1m(row)}, and
\code{ism(genotype)}.
\item[\code{Gu}] Optional known precision matrix for the levels of the final
main-effect term. For \code{mmes}, \code{Gu} must have row and column
names matching the main-effect levels and \code{attr(Gu, "inverse")} must be
\code{TRUE}. If omitted, an identity precision matrix is used.

\code{Gu} can also be supplied in 3-column (row, column, value) format as a
numeric matrix or data frame, e.g. the output of ASReml-R \code{ainverse()}
or the \code{listAinv} element of \code{nadiv::makeAinv()}. The first two
columns are 1-based level indices and the third the precision values. The
level names must be given in \code{attr(Gu, "rowNames")}. The lower
triangle, the upper triangle, or both may be given (entries present in both
must agree), and every level needs a diagonal entry. This format is detected
automatically, is always taken to be a precision matrix, and does not require
the \code{inverse} attribute. Symmetric precision matrices are stored with
one triangle only.
\item[\code{sigma2}] Positive starting value for the single overall covariance
scale of this \code{vsm} term. If \code{NULL} (the default), \code{mmes}
replaces it with a data-driven starting value computed from the response
and fixed effects; supplying a value here always takes precedence and is
never overridden.
\item[\code{fixedSigma2}] Logical indicating whether the product-level
\eqn{\sigma^2}{} is fixed at its starting value.
\item[\code{rotation}] Logical indicating whether the supplied \code{Gu} precision
matrix should be eigen-decomposed for Lee--van der Werf rotation in
\code{mmes}. Rotation is opt-in and currently requires a balanced Gaussian
random-effect term and a rotation-invariant residual structure (see
\code{\LinkA{mmes}{mmes}}).
\item[\code{isFixed}] Logical retained for compatibility.  If \code{TRUE}, return
the combined design matrix instead of the structured \code{vsm} object.
\item[\code{verbose}] Logical controlling messages, including notification when
levels present in \code{Gu} are appended to the model matrix.
\end{ldescription}
\end{Arguments}
%
\begin{Details}\relax
Let \eqn{K_1,\ldots,K_m}{} be the dimensionless covariance shapes supplied by
the covariance constructors.  The covariance represented by one \code{vsm}
term is
\deqn{\Sigma = \sigma^2(K_1 \otimes K_2 \otimes \cdots \otimes K_m).}{}
For a random effect with known relationship precision \code{Gu}, the factor
product determines the covariance among the crossed covariance coordinates and
\code{Gu} determines the relationship structure among the levels of the final
main effect.

The final argument in \code{...} is the main-effect incidence term.  Every term
before it is a covariance-shaping factor.  Thus
\code{vsm(dsm(location), ar1m(row), ism(genotype), Gu=Ainv)} represents a
location-by-row covariance shape crossed with genotype.  Kronecker ordering is
left to right: earlier factors are the slow index and later factors are the fast
index.

All optimizer coordinates are stored on unconstrained working scales.  Each
CovarianceFactor descriptor also carries its evaluator, derivative strategy,
reporting transformation, and trust-region caps.  Built-in high-use structures
may use native C++ covariance primitives, while newer or user-defined
structures can use generic R callbacks without requiring changes to
\code{ai\_mme\_sp2}.

For residual covariance models, the covariance-shaping factors must define
exactly one product coordinate for each observation.  \code{mmes} combines
that local coordinate with an independently constructed residual block index,
so a large observation-level identity design is not materialized.

When \code{rotation=TRUE}, \code{Gu} is decomposed as
\deqn{Gu=U\Lambda U',}{}
where \eqn{\Lambda}{} is diagonal. The aligned original precision, diagonal
precision, and eigenvectors are retained in the returned descriptor. The
\code{mmes} front end validates the observation layout and applies the
engine-specific rotation after observation filtering. Residual covariance
must satisfy \eqn{R_{(i,a),(l,b)} = c_{ab}\delta_{il}}{}, where \eqn{i,l}{} index
relationship levels and \eqn{a,b}{} index rotation blocks, so that
\eqn{U'RU = R}{}; otherwise \code{mmes} stops with an error.
\end{Details}
%
\begin{Value}
A list containing the random-effect design blocks in \code{Z}, the precision
matrix in \code{Gu}, the flattened Kronecker descriptor in \code{covStruct},
the product design, and, for residual models, the local covariance-coordinate
index. When rotation is requested, \code{GuRot} contains the diagonal
precision and \code{rotation} contains its eigen metadata. The
\code{covStruct} object contains descriptor version 2
CovarianceFactor objects.
\end{Value}
%
\begin{SeeAlso}\relax
\code{\LinkA{mmes}{mmes}}, \code{\LinkA{ism}{ism}}, \code{\LinkA{dsm}{dsm}}, \code{\LinkA{atm}{atm}}, \code{\LinkA{usm}{usm}},
\code{\LinkA{csm}{csm}}, \code{\LinkA{ar1m}{ar1m}}, \code{\LinkA{ar2m}{ar2m}}, \code{\LinkA{ar3m}{ar3m}},
\code{\LinkA{mam}{mam}},
\code{\LinkA{corgm}{corgm}}, \code{\LinkA{fam}{fam}}, \code{\LinkA{antem}{antem}}, \code{\LinkA{rrm}{rrm}},
\code{\LinkA{maternm}{maternm}}, \code{\LinkA{toeplitzm}{toeplitzm}}, \code{\LinkA{sar}{sar}}, \code{\LinkA{car}{car}},
and \code{\LinkA{ownm}{ownm}}. See \code{\LinkA{dsumm}{dsumm}} for section-specific (direct-sum)
residual structures.
\end{SeeAlso}
%
\begin{Examples}
\begin{ExampleCode}

####=========================================####
#### For CRAN time limitations most lines in the
#### examples are silenced with one '#' mark,
#### remove them and run the examples
####=========================================####

data(DT_example, package="enhancer")
DT <- DT_example
head(DT)


str(with(DT, vsm(dsm(Env),ism(Name))))
str(with(DT, vsm(csm(Env),ism(Name))))
str(with(DT, vsm(toeplitzm(Env),ism(Name))))

## Gu in 3-column (row, column, value) format with level names in rowNames
Ainv <- data.frame(Row=c(1,2,2,3), Column=c(1,1,2,3),
                   Ainverse=c(2,-1,2,1))
attr(Ainv, "rowNames") <- c("a","b","c")
id <- factor(c("a","b","c","a"))
str(vsm(ism(id), Gu=Ainv)$Gu)

\end{ExampleCode}
\end{Examples}
\inputencoding{utf8}
\HeaderA{vsr}{variance structure specification}{vsr}
%
\begin{Description}\relax
\code{vsr} DEPRECATED NOW. was the main function to build the variance-covariance structure for the random effects to be fitted in the \code{\LinkA{mmer}{mmer}} solver. 
\end{Description}
%
\begin{Usage}
\begin{verbatim}
  vsr(..., Gu=NULL, Gti=NULL, Gtc=NULL, reorderGu=TRUE, buildGu=TRUE)
\end{verbatim}
\end{Usage}
%
\begin{Arguments}
\begin{ldescription}
\item[\code{...}] variance structure to be specified following the logic desired in the internal kronecker product. For example, if user wants to define a diagonal variance structure for the random effect 'genotypes'(g) with respect to a random effect 'environments'(e), this is:

\code{var(g) = G.e @ I.g}

being \code{G.e}  a matrix containing the variance covariance components for g (genotypes) in each level of e (environments), \code{I.g} is the covariance among levels of g (genotypes; i.e. relationship matrix), and @ is the kronecker product. This would be specified in the mmer solver as:

\code{random=\textasciitilde{}vsr(dsr(e),g)}

One strength of sommer is the ability to specify very complex structures with as many kronecker products as desired. For example:

\code{var(g) = G.e @ G.f @ G.h @ I.g}

is equivalent to

\code{random=\textasciitilde{}vsr(e,f,h,g)}

where different covariance structures can be applied to the levels of \code{e,f,h}. For more examples please see the vignettes 'sommer.start' available in the package.

\item[\code{Gu}] matrix with the known variance-covariance values for the levels of the u.th random effect (i.e. relationship matrix among individuals or any other known covariance matrix). If \code{NULL}, then an identity matrix is assumed. The Gu matrix can have more levels than the ones present in the random effect linked to it but not the other way around. Otherwise, an error message of missing level in Gu will be returned.

\item[\code{Gti}] matrix with dimensions t x t (t equal to number of traits) with initial values of the variance-covariance components for the random effect specified in the .... argument. If \code{NULL} the program will provide the initial values. The values need to be scaled, see Details section.

\item[\code{Gtc}] matrix with dimensions t x t (t equal to number of traits) of constraints for the variance-covariance components for the random effect specified in the ... argument according to the following rules:

\code{0: not to be estimated}

\code{1: estimated and constrained to be positive (i.e. variance component)}

\code{2: estimated and unconstrained (can be negative or positive, i.e. covariance component)}

\code{3: not to be estimated but fixed (value has to be provided in the Gti argument)}

In the multi-response scenario if the user doesn't specify this argument the default is to build an unstructured matrix (using the \code{\LinkA{unsm}{unsm}}() function). This argument needs to be used wisely since some covariance among responses may not make sense. Useful functions to specify constraints are; \code{\LinkA{diag}{diag}}(), \code{\LinkA{unsm}{unsm}}(), \code{\LinkA{fixm}{fixm}}().


\item[\code{reorderGu}] a \code{TRUE/FALSE} statement if the Gu matrix should be reordered based on the names of the design matrix of the random effect or passed with the custom order of the user. This may be important when fitting covariance components in a customized fashion. Only for advanced users.

\item[\code{buildGu}] a \code{TRUE/FALSE} statement to indicate if the Gu matrix should be built in R when the value for the argument \code{Gu=NULL}. Repeat, only when when the value for the argument \code{Gu} is equal to \code{NULL}. In some cases when the incidence matrix is wide (e.g. rrBLUP models) the covariance structure is a huge p x p matrix that can be avoided when performing matrix operations. By setting this argument to \code{FALSE} it allows to skip forming this covariance matrix.

\end{ldescription}
\end{Arguments}
%
\begin{Details}\relax
When providing initial values in the \code{Gti} argument the user has to provide scaled variance component values. The user can provide values from a previous model by accessing the \code{sigma\_scaled} output from an \code{mmer} model or if an specific value is desired the user can obtain the scaled value as:

\code{m = x/var(y)}

where \code{x} is the desired initial value and \code{y} is the response variable. You can find an example in the \code{\LinkA{DT\_cpdata}{DT.Rul.cpdata}} dataset.

\end{Details}
%
\begin{Value}
\begin{description}

\item[\$res] a list with all neccesary elements (incidence matrices, known var-cov structures, unknown covariance structures to be estimated and constraints) to be used in the mmer solver.

\end{description}

\end{Value}
%
\begin{Author}\relax
Giovanny Covarrubias-Pazaran
\end{Author}
%
\begin{References}\relax
Covarrubias-Pazaran G (2016) Genome assisted prediction of quantitative traits using the R package sommer. PLoS ONE 11(6): doi:10.1371/journal.pone.0156744 

Covarrubias-Pazaran G (2018) Software update: Moving the R package sommer to
multivariate mixed models for genome-assisted prediction. doi:
https://doi.org/10.1101/354639
\end{References}
%
\begin{Examples}
\begin{ExampleCode}

data(DT_example, package="enhancer")
DT <- DT_example
A <- A_example

## ============================ ##
## example to without structure
## ============================ ##

mix <- mmer(Yield~Env, 
            random= ~ vsr(Name),
            rcov=~ vsr(units),
            data=DT)

\end{ExampleCode}
\end{Examples}
\inputencoding{utf8}
\HeaderA{wald\_mmes}{Wald tests for fixed effects with optional small-sample denominator df}{wald.Rul.mmes}
\aliasA{print.wald.mmes}{wald\_mmes}{print.wald.mmes}
\keyword{models}{wald\_mmes}
%
\begin{Description}\relax
Pseudo analysis of variance for the fixed terms of an \code{\LinkA{mmes}{mmes}} fit,
based on incremental or conditional Wald statistics, with chi-square,
residual, Satterthwaite or Kenward-Roger reference distributions.
\end{Description}
%
\begin{Usage}
\begin{verbatim}
wald_mmes(object, ssType = c("incremental", "conditional"),
          denDF = c("none", "residual", "satterthwaite", "kr"), terms = NULL)
\end{verbatim}
\end{Usage}
%
\begin{Arguments}
\begin{ldescription}
\item[\code{object}] a model of class \code{"mmes"}.
\item[\code{ssType}] \code{"incremental"} tests each term adjusted for the terms
preceding it in the fixed formula. \code{"conditional"} tests each term
adjusted for every other term that does not contain it, so marginality is
respected: a term is never adjusted for a term that contains it
structurally (A:B contains A) or implicitly (its columns, together with the
intercept, span the tested term, e.g. locations coded separately within
regions).
\item[\code{denDF}] \code{"none"} gives chi-square tests of the Wald statistic.
\code{"residual"} uses F tests with n - rank(X) denominator df.
\code{"satterthwaite"} and \code{"kr"} (Kenward-Roger) compute approximate
denominator df; \code{"kr"} also uses the Kenward-Roger adjusted covariance
of the fixed effects and scaled F statistic.
\item[\code{terms}] optional subset of fixed terms (as named in \code{object\$Dtable}).
\end{ldescription}
\end{Arguments}
%
\begin{Details}\relax
Let \eqn{\Phi}{} be the covariance of the fixed effects and \eqn{M=\Phi^{-1}}{}.
For a given order of the terms, the upper Cholesky factor \eqn{U}{} of \eqn{M}{}
defines \eqn{z=U b}{}; the incremental Wald statistic of a term is the sum of the
squared elements of \eqn{z}{} belonging to it, i.e. a test of \eqn{L b=0}{} with
\eqn{L}{} the corresponding rows of \eqn{U}{}. Conditional tests use the order
[terms not containing the tested term, tested term, terms containing it].

For Satterthwaite and Kenward-Roger df the model inputs are rebuilt from the
stored call and \eqn{X^\top V^{-1}X}{} is re-evaluated at perturbed covariance
parameters (on the reported scale of \code{covPar}); its numerical first and
second derivatives give the quantities \eqn{P_i}{} and the Kenward-Roger
adjustment. The covariance of the covariance parameters is \code{theta\_se},
the inverse average information. The second-derivative term of \eqn{V}{}
enters through the Hessian of \eqn{X^\top V^{-1}X}{}; it vanishes for
variance-component models, where the method is the exact Kenward-Roger
approximation. Because the average (rather than expected) information is
used, df can differ slightly from software using the expected information.
The re-evaluation requires the data and objects used in the original call
to be available, REML estimation, and a Gaussian fit.

\code{anova(object)} with a single model returns \code{wald\_mmes(object, ...)}.
\end{Details}
%
\begin{Value}
A data frame of class \code{"wald.mmes"} with columns \code{Df},
\code{denDF}, \code{Wald}, \code{F.value} and \code{p.value}, one row per term.
\end{Value}
%
\begin{References}\relax
Kenward MG and Roger JH (1997). Small sample inference for fixed effects from
restricted maximum likelihood. Biometrics 53: 983-997.
\end{References}
%
\begin{SeeAlso}\relax
\code{\LinkA{mmes}{mmes}}, \code{\LinkA{predict.mmes}{predict.mmes}}, \code{\LinkA{anova.mmes}{anova.mmes}}
\end{SeeAlso}
%
\begin{Examples}
\begin{ExampleCode}
data(DT_yatesoats, package = "enhancer")
m <- mmes(Y ~ V * N,
  random = ~ B + B:MP,
  rcov = ~ units,
  data = DT_yatesoats,
  control = mmes_control(diagnostics = list(verbose = FALSE)))
wald_mmes(m)
wald_mmes(m, ssType = "conditional", denDF = "kr")
\end{ExampleCode}
\end{Examples}
\printindex{}
\end{document}
